% concatenated from main.tex
\documentclass[a4paper,11p]{amsart}

%\usepackage{amsmath, amsthm, amssymb, amsfonts, mathtools, bbold, stmaryrd}
\usepackage[all,arc]{xy}
\usepackage[dvipsnames]{xcolor}
\usepackage{graphicx} % Required for inserting images
\usepackage[centertags]{amsmath}
\usepackage{amscd}
\usepackage{amsthm}
\usepackage{amssymb}
\usepackage{dsfont} %%Package for evaluation functions
\usepackage{MnSymbol} %%Package for the star operator
\usepackage{algorithm}
\usepackage{algpseudocode}
\usepackage{tikz-cd}
\usetikzlibrary{cd}
\usepackage{tikz}
\usepackage{relsize}
\usepackage{ytableau}
\usepackage{mdframed}
\usepackage{comment}
\usepackage{enumitem}
\usepackage[utf8]{inputenc}
\usepackage[T1]{fontenc}

\usepackage{textcomp}

%\usepackage{cleveref}

\usepackage{hyperref}
\hypersetup{
    colorlinks,
    linkcolor={MidnightBlue},
    citecolor={Green}
}
\usepackage[nameinlink]{cleveref}

\usepackage[normalem]{ulem}

\newcommand{\m}{\mathfrak m}
\newtheorem{notation}{Notation}[section]
\newtheorem{defn0}{Definition}[section]
\newtheorem{prop0}[defn0]{Proposition}
\newtheorem{thm0}[defn0]{Theorem}
\newtheorem{lemma0}[defn0]{Lemma}
\newtheorem{corollary0}[defn0]{Corollary}
\newtheorem{example0}[defn0]{Example}
\newtheorem{remark0}[defn0]{Remark}
\newtheorem{conjecture0}[defn0]{Conjecture}
\DeclareFontFamily{U}{MnSymbolC}{}
\DeclareSymbolFont{MnSyC}{U}{MnSymbolC}{m}{n}
\DeclareFontShape{U}{MnSymbolC}{m}{n}{
    <-6>  MnSymbolC5
   <6-7>  MnSymbolC6
   <7-8>  MnSymbolC7
   <8-9>  MnSymbolC8
   <9-10> MnSymbolC9
  <10-12> MnSymbolC10
  <12->   MnSymbolC12}{}
\DeclareMathSymbol{\aprod}{\mathbin}{MnSyC}{'270}

\newenvironment{definition}{ \begin{defn0}}{\end{defn0}}
\newenvironment{proposition}{\bigskip \begin{prop0}}{\end{prop0}}
\newenvironment{theorem}{\bigskip \begin{thm0}}{\end{thm0}}
\newenvironment{lemma}{\bigskip \begin{lemma0}}{\end{lemma0}}
\newenvironment{corollary}{\bigskip \begin{corollary0}}{\end{corollary0}}
\newenvironment{example}{ \begin{example0}\rm}{\end{example0}}
\newenvironment{remark}{ \begin{remark0}\rm}{\end{remark0}}
\newenvironment{conjecture}{\begin{conjecture0}}{\end{conjecture0}}
\newtheorem{alg}{Algorithm}
\newcommand{\K}{\mathbb K}


\newcommand{\ale}[1]{{\color{magenta}#1}}


\title{A refinement on the local cactus rank algorithm}

\author{Alessandra Bernardi, Oriol Reig Fité}

\address{Universit\`a di Trento, Via Sommarive, 14 - 38123 Povo (Trento), Italy}
\email{alessandra.bernardi@unitn.it, oriol.reigfite@unitn.com}


\makeatletter
\@namedef{subjclassname@2020}{\textup{2020} Mathematics Subject Classification}
\makeatother

\subjclass[2020]{14N07, 13H10}
\keywords{Symmetric tensors, Tensor decomposition, Local cactus rank, Gorenstein algebras,  Algorithms.}

\begin{document}
\begin{abstract} We present an algorithm to  recover a minimal local apolar scheme to a homogeneous polynomial $F$.
The scheme’s socle degree determines whether it is evinced by a Generalized Additive Decomposition (GAD) of $F$ or of an extension.
We give constructive procedures for both cases and compute the Hilbert function efficiently via Hankel operators.
\end{abstract}
\maketitle
% \begin{document}
% %\maketitle
% \begin{abstract}
% We use the structure of local Gorenstein algebras to refine the existing algorithm that computes local schemes of minimal length that are apolar to a homogeneous polynomial. We show effective methods to compute, in the setting of the algorithm, the Hilbert function of the resulting algebra as well as its symmetric decomposition. Finally, we describe how to compute the generalized additive decomposition that evinces the scheme that reveals the cactus rank of the polynomial. 
% \end{abstract}
% \begin{keyword}
% Symmetric tensors, Tensor decomposition, Local cactus rank, Gorenstein algebras,  Algorithms.
% \MSC[2020] .
% \end{keyword}
% \begin{frontmatter}
% \title{A refinement on the local cactus rank algorithm}
% \author{Alessandra Bernardi}
% \ead{alessandra.bernardi@unitn.it}
% \author{Oriol Reig}
% \ead{oriol.reigfite@unitn.com}
% \date{\today}
% \end{frontmatter}
% \tableofcontents
\begin{comment}
\section{Introduction}
The symmetric tensor decomposition algorithm, first presented in \cite{Bernard}, computes the Waring rank and the Waring decomposition of a given homogeneous polynomial $F$. It can also be used, under mild modifications, to compute the minimal length of an apolar scheme to a homogeneous polynomial (\cite{Alessandra}) known as \emph{the cactus rank of $F$} \cite{BR13}, and a minimal apolar scheme is also called \emph{a cactus scheme for $F$} \cite{JarekWer}, \cite{IK}. In many cases the scheme revealing the cactus rank is local, and therefore it is an interesting problem to study the local cactus rank of a polynomial. Due to the fact that a scheme apolar to a homogeneous polynomial is Gorenstein (\cite{JarekWer}), this opens the possibility of exploiting the properties of local Artinian Gorenstein algebras, which have been actively studied in the last decades (\cite{Iarrobino}, \textcolor{red}{add some others}).

\bigskip

The starting point of the symmetric tensor decomposition algorithm in \cite{Bernard} and its tangential and cactus  variants in \cite{Alessandra} is to assign, using apolarity, to a homogeneous polynomial $F$ of degree $d$ a linear functional defined on polynomials on degree at most $d$ in the affine ring. This linear form is the truncation of an infinite series, which we can identify with some linear functional of the whole affine ring. The crux of the algorithm is to parameterize the extension of this truncation and impose conditions on the new variables to prove the existence of a linear form satisfying the requirements to define a scheme of minimal length. Roughly speaking, one needs to access the Gorenstein algebra associated with this unknown dual generator that extends the given polynomial through the matrices corresponding by the multiplication by the variables in the algebra. For this, one first needs a basis of the algebra, and then impose that the associated matrices define an Artinian algebra. For a fixed choice of basis, this is translated into a set of closed conditions on the parameters given by the pairwise commutativity of the matrices (\Cref{Thm: ExtensionBernard}). 
\bigskip
These structural matrices are built from the Hankel operator of the extended linear form (see \eqref{MatMult}), which has some numerical part, given by the coefficients of the polynomial we want to decompose, and some parametrized part, corresponding to the higher terms of the infinite series. It may happen that for a candidate basis the matrices commute, automatically terminating the rank computation of the algorithm. These cases correspond to schemes with low regularity (\cite{Alessandra}, \cite{GAD}). In general, the matrices are parametrized and need to satisfy some nonlinear system given by the pairwise commutativity. An additional complication is that not all conditions on the parameters are closed, since the solutions of the polynomial system have to be compatible with the full rank condition of the minor in the Hankel operator associated to the basis (\Cref{Prop: BasisInHankel}). 
\bigskip
Therefore there are two critical aspects of the extension step in the algorithm. One is the polynomial system associated to the commutation of the matrices. The other is the choice of the basis of the algebra. Its elements need not be arbitrary, for we know they can be chosen monomials and forming a set closed under division and containing all the variables of the polynomial (\Cref{Prop: BasisInHankel}). For high ranks and large number of variables, the possibilities are still too large to be tested by hand. In this work we show that, in the local case, an additional criterion for selecting candidate basis is that they reveal the grading structure of the local algebra \Cref{lemma: BasisFromHF}. This means that we can start with a guess on the Hilbert function, which for local Gorenstein algebras satisfy certain properties, and then choose a monomial basis accordingly.
\bigskip
The existence of parameters verifying the conditions described above imply the existence of some zero dimensional scheme whose ideal is apolar to our starting polynomial, with length the size of the matrices, that is, that the cactus rank of the polynomial is at most the size of the matrices. There is still the question of how much we can know about the scheme. This problem was studied in \cite{BJPR}, \cite{GAD}, proving that generalized additive decompositions of a homogeneous polynomial define a direct sum of local Artinian Gorenstein algebra whose geometric ideal is apolar to the polynomial. In \cite{GAD} the authors deal with the direction from decompositions to apolar schemes of minimal length. In \Cref{Section:GAD} we solve the converse for the local case. That is, how to find a generalized additive decomposition evincing an apolar scheme of minimal length.  
\bigskip
In the last part we show that there is a simple way to compute the Hilbert function and its symmetric decomposition of the Gorenstein Artinian algebra associated to a given dual generator from the Hankel operator of the functional.
\end{comment}
\section*{Introduction}

The study of ranks of homogeneous forms has a long history, starting from Sylvester’s classical theorem on binary forms (cf. eg. \cite{BERNARDI201134, Bernard}), and continuing through the modern theory of secant varieties and tensor ranks.  
Among these notions, the \emph{cactus rank} of a form, defined as the minimal length of a zero-dimensional apolar scheme (cf. \cite{IK, BR13}), plays a key role in algebraic geometry and its applications (cf. eg. \cite{MR3729273}).  
When the minimal apolar scheme to a given polynomial $F$ is local, its structure captures subtle information about the form and the possible additive decompositions of either $F$ itself or an extension of $F$.

A classical way to describe such forms is via \emph{Generalized Additive Decompositions} (GADs). The connection between local Gorenstein schemes and GADs has been explored in works such as \cite{IK,MR3250539,GAD, BJPR}, where criteria are given to decide when a local scheme is evinced by a GAD.  
In these settings, the Macaulay correspondence between Gorenstein algebras and divided power series is a standard tool, though the presentation is often given in the standard polynomial ring and the divided power formalism is introduced only when necessary.

Recent algorithmic approaches to cactus rank and GADs have focused on practical computation of apolar schemes \cite{Alessandra, GAD}.  
However, most of these algorithms are formulated in the standard polynomial setting, and the explicit transition to divided powers is typically performed only in the final stages.  
This choice, while standard in the literature, may introduce additional technicalities when performing coordinate changes, homogenizations, or computations involving higher socle degrees.

\medskip
\noindent
\textbf{Our contribution.}  
In this paper, we develop a complete algorithm to determine the local cactus rank of a form $F$ \emph{entirely in the divided power setting}, to recover its minimal local apolar scheme, and to decide whether the scheme is evinced by either a GAD of $F$ or of an extension of $F$.  
The algorithm is constructive: it computes the support of the minimal scheme from multiplication operators, extracts the Hilbert function and its symmetric decomposition, and recovers the GAD when possible.  
Our presentation is self-contained in the divided power setting, so that all coordinate changes and dual actions are explicit.

\medskip
\noindent
\textbf{Structure of the paper.}  
\Cref{sec:prelim} recalls the necessary background on divided powers, the dual space, Gorenstein algebras, and Hilbert functions, setting the notation used throughout the paper.  
Although these notions are classical, we present them in detail because the algorithms are implemented entirely in the divided power context, whereas the standard exposition is usually given in the polynomial ring.   
An expert reader may consult it mainly to verify the notation. 

\Cref{sec:main} recalls the background on apolarity, local Gorenstein algebras, and Generalized Additive Decompositions (GAD), and contains the main results. It includes the description of the local cactus rank algorithm, and the recovery of the local GAD in the two cases where the socle degree of the minimal local apolar scheme is either smaller or bigger than the degree of $F$, and
the computation of the Hilbert function via Hankel operators  with several examples and computational remarks.

Our main contributions are:
\begin{itemize}
    \item An explicit algorithm for computing local cactus rank and minimal local apolar schemes directly in the divided power setting.
    \item A constructive method to determine whether a minimal local apolar scheme is evinced by a GAD of a polynomial $F$, and to recover the decomposition when it exists and generalize it to an extension of $F$.
    \item A unified treatment of background material and algorithmic procedures in the divided power framework, avoiding repeated transitions from standard polynomials.
\end{itemize}


\section*{Acknowledgements}
We are grateful to Jarosław Buczyński, Maciej Gałązka, Joachim Jelisiejew, Daniele Taufer, and Bernard Mourrain for insightful conversations that greatly benefited this work.

\medskip

This work has been supported by European Union’s HORIZON–MSCA-2023-DN-JD programme under the Horizon Europe (HORIZON) Marie Skłodowska-Curie Actions, grant agreement 101120296 (TENORS).

%\section{Duality, cactus rank and Hilbert functions}
\section{Preliminaries and Background}\label{sec:prelim}
In this section we revisit some foundational concepts and fix the notation used throughout the paper. Our aim is to provide a self-contained introduction to Gorenstein algebras and Hilbert functions in the divided power setting, making these essential tools accessible even to non-expert readers. 



%\subsection{Divided powers and the dual space}

\medskip

    Let $\mathbb{K}$ be an algebraically closed field of characteristic zero, and let $V$ be a $\mathbb{K}$-vector space of dimension $n+1$. Denote by $S := %\mathcal{S}(V) \cong 
    \mathbb{K}[x_0, \ldots, x_n]$ the symmetric algebra of $V$, identified with the homogeneous coordinate ring of projective space $\mathbb{P}^n$. Let also $R := \mathbb{K}[x_1, \ldots, x_n]$ be the coordinate ring of the standard affine chart $\{x_0 \neq 0\} \subset \mathbb{P}^n$.

\medskip

Denote by $S_d$ (respectively $R_d$) the space of homogeneous polynomials of degree $d$ in $S$ (respectively in $R$). Similarly, $S_{\leq d}$ and $R_{\leq d}$ denote polynomials of degree at most $d$. The dual space of a vector space $V$ is denoted by $V^*:=\mathrm{Hom}_\mathbb{K}(V, \mathbb{K})$.


\medskip

For multi-indices $\alpha = (\alpha_1, \ldots, \alpha_n)$ and $\beta = (\beta_1, \ldots, \beta_n)$ in $\mathbb{N}^n$, we use standard multi-index notations:
\[
\binom{d}{\alpha} := \frac{d!}{\alpha!(d - |\alpha|)!}, \quad \text{where} \quad \alpha! = \prod_{i=1}^n \alpha_i!, \quad |\alpha| = \sum_{i=1}^n \alpha_i,
\]



and similarly for $\binom{\alpha}{\beta}:= \frac{\alpha!}{\beta!(\alpha - \beta)!}$. 

Given $\zeta \in \mathbb{K}^n$, the evaluation at $\zeta$ is the functional $\mathds{1}_\zeta : R \to \mathbb{K}$, defined as $\mathds{1}_\zeta(p) = p(\zeta)$.



Finally, we abbreviate products of the form $(x_0 - \zeta_0)^{\alpha_0} \cdots (x_n - \zeta_n)^{\alpha_n}$ as $(x - \zeta)^\alpha$.

 
\subsection{Divided powers and the dual space}
The dual space $S^* := \mathrm{Hom}_{\mathbb{K}}(S, \mathbb{K})$ can be identified with the power series ring in the basis of divided powers. Consider the monomial basis ${x^\alpha}$ of $S$ and its dual basis ${Y^{(\alpha)}}$ in $S^*$, defined by $Y^{(\alpha)}(x^\beta)=\delta(\alpha, \beta)$. This yields an isomorphism
%Fix the monomial basis $ \{x^\alpha\}_{\alpha \in \mathbb{N}^{n+1}} $ of $ S $, and consider its dual basis in $ S^* $, denoted $ \{Y^{(\alpha)}\}_{\alpha \in \mathbb{N}^{n+1}} $, where
%\[
%Y^{(\alpha)}(x^\beta) =
%\begin{cases}
%1 & \text{if } \alpha = \beta, \\
%0 & \text{otherwise}.
%\end{cases}
%\]
%We define $ Y_i := Y^{(e_i)} $, with $ e_i $ the standard basis vector in $ \mathbb{N}^{n+1} $. This gives an isomorphism 
of $ \K $-vector spaces:
\[
\Psi: S^* \longrightarrow \K[[Y_0, \ldots, Y_n]], \qquad
\Lambda \mapsto \sum_{\alpha \in \mathbb{N}^{n+1}} \Lambda(x^\alpha)\, Y^{(\alpha)}.
\]
An affine translation 
\begin{equation}\label{affine:tranlsation}
\varphi: R \to R, \qquad \varphi(x_i) = x_i + \zeta_i
\end{equation}
induces a dual map $\varphi^*: R^* \to R^*$ acting as differential operators:
\begin{equation}\label{dual:affine:tranlsation}
    \varphi^*: R^* \longrightarrow R^*, \qquad \varphi^*(\Lambda)(p) := \Lambda(\varphi(p)).
\end{equation}

In particular,
$\varphi^*(1)=\sum_{\alpha\in \mathbb N^n} \varphi^*(1)(x^\alpha) Y^{(\alpha)}=\sum_{\alpha\in \mathbb N^n} \zeta^\alpha Y^{\alpha}= \mathds{1}_\zeta$ where $ \mathds{1}_\zeta $ is the evaluation at the point $ \zeta $.

\begin{lemma}\label{lemma:DualChangeBasis}
Let $ \K $ be a field of characteristic zero. Let $ \zeta \in \K^n $ and $ \varphi$ be the affine change of coordinates of \eqref{affine:tranlsation}. Then, for every $ \alpha \in \mathbb{N}^n $,
\[
\varphi^*\left(Y^{(\alpha)}\right) = \frac{1}{\alpha!} \mathds{1}_\zeta \circ \frac{\partial^\alpha}{\partial x^\alpha}.
\]
\end{lemma}

\begin{proof}
Due to linearity, we evaluate the expression on $x^\beta$:
\[
\varphi^*\left(Y^{(\alpha)}\right)\left(x^\beta\right) = Y^{(\alpha)}\left((x + \zeta)^\beta\right).
\]
By the binomial theorem, $(x + \zeta)^\beta = \sum_{k \le \beta} \binom{\beta}{k} x^k \zeta^{\beta - k}$. Applying $Y^{(\alpha)}$ extracts the coefficient of $x^\alpha$, which is:
\[
\prod_{i=1}^n \binom{\beta_i}{\alpha_i} \zeta_i^{\beta_i - \alpha_i} \quad \text{for } \alpha \leq \beta, \text{ and } 0 \text{ otherwise}.
\]
This is equivalent to $\frac{1}{\alpha!} \frac{\beta!}{(\beta - \alpha)!} \zeta^{\beta - \alpha}$ (when $\alpha \le \beta$), precisely matching $\left(\frac{1}{\alpha!} \mathds{1}_\zeta \circ \frac{\partial^\alpha}{\partial x^\alpha}\right)\left(x^\beta\right)$.
\end{proof}

\begin{comment}
    
\begin{proof}
By linearity, it suffices to evaluate the expression on a monomial $ x^\beta $:
\[
\varphi^*\left(Y^{(\alpha)}\right)\left(x^\beta\right) = Y^{(\alpha)}\left((x + \zeta)^\beta\right).
\]
Expanding each term using the binomial theorem:
\[
(x + \zeta)^\beta = \prod_{i=1}^n \left( \sum_{k_i=0}^{\beta_i} \binom{\beta_i}{k_i} x_i^{k_i} \zeta_i^{\beta_i - k_i} \right).
\]
The coefficient of $ x^\alpha $ in this expansion is:
\[
\prod_{i=1}^n \binom{\beta_i}{\alpha_i} \zeta_i^{\beta_i - \alpha_i} \quad \text{if } \alpha \leq \beta, \text{ and } 0 \text{ otherwise}.
\]
Rewriting using factorials:
\[
=
\begin{cases}
\frac{1}{\alpha!} \frac{\beta!}{(\beta - \alpha)!} \zeta^{\beta - \alpha} & \text{if } \alpha \leq \beta, \\
0 & \text{otherwise},
\end{cases}
\]
which is exactly $ \left(\frac{1}{\alpha!} \mathds{1}_\zeta \circ \frac{\partial^\alpha}{\partial x^\alpha}\right)\left(x^\beta\right) $.
\end{proof}



\end{comment}



%\begin{proof}
%    By linearity it suffices to check it on monomials:
%\[\varphi^*(Y^{(\alpha)})(x^\beta)=Y^{(\alpha)}((x+\zeta)^\beta)=
%Y^{(\alpha)}\left( \prod_{i=1}^n \left( \sum_{k_i=0}^{\beta_i} \binom{\beta_i}{k_i}x_i^{k_i}\zeta_i^{\beta_i-k_i}\right) \right)=
%\]
%\[Y^{(\alpha)}\left(\sum_{k_1,\ldots,k_n=0}^{\beta_1,...,\beta_n} \left( \prod_{i=1}^n \binom{\beta_i}{k_i}\zeta_i^{\beta_i-k_i}x_i^{k_i} \right)\right)=
%\begin{cases}
%     \prod_{i=1}^n \binom{\beta_i}{\alpha_i}\zeta_i^{\beta_i-\alpha_i} \quad \text{if $\alpha\leq \beta$}\\
%    0 \quad \quad \text{if $\alpha>\beta$}
%\end{cases}\]
%\[=
%\begin{cases}
%    \frac{1}{\alpha!}\frac{\beta!}{(\beta-\alpha)!}\zeta^{\beta-\alpha} \text{if $\alpha\leq \beta$}\\
%    0 \quad \quad \text{if $\alpha>\beta$}
%
%\end{cases} =
%\begin{cases}
%    \frac{1}{\alpha!}\mathds{1}_\zeta\circ \frac{\partial^\alpha}{\partial x^\alpha} \quad \text{if $\alpha\leq \beta$}\\
%    0 \quad \quad \text{if $\alpha>\beta$}

%\end{cases}
%.\]
%\end{proof}

From this, we deduce:
\[
\left( \mathds{1}_\zeta \circ \frac{\partial^\alpha}{\partial x^\alpha} \right) \left( \frac{1}{\alpha!} (x - \zeta)^\beta \right) =
\begin{cases}
1 & \text{if } \alpha = \beta, \\
0 & \text{otherwise}.
\end{cases}
\]
This implies that $\mathcal{B}_\zeta := \left\{ \frac{1}{\alpha!}(x - \zeta)^\alpha \right\}_{\alpha \in \mathbb{N}^n}$ forms a basis for $R$. Its dual basis is given by $ \{\partial^\alpha_\zeta\} \subset R^* $, where $ \partial^\alpha_\zeta(p) := \left( \mathds{1}_\zeta \circ \frac{\partial^\alpha}{\partial x^\alpha} \right)(p) $, and $\partial^{\alpha}_{\zeta}$ acts naturally as $\left(\mathds{1}_{\zeta} \circ x^{\alpha}(\partial)\right)(p)$. This establishes an isomorphism $S^* \to \mathbb{K}[[\partial_{1, \zeta},\ldots,\partial_{n, \zeta}]]$ mapping $\Lambda \mapsto \sum_{\alpha\in \mathbb N^{n+1}} \partial^{\alpha}_{\zeta}(x^{\alpha}) \partial^{\alpha}_{\zeta}$. Note that this demonstrates the dual change of basis from $\left\{Y^{(\alpha)}\right\}$ to $\mathcal{B}_\zeta$'s dual basis (see \Cref{lemma:DualChangeBasis}). This formalism also extends to positive characteristic fields $\mathbb{K}$ by replacing $\mathcal{B}_\zeta$ with $\left\{(x-\zeta)^\alpha\right\}_{\alpha \in \mathbb{N}^n}$.

\begin{comment}
    


From this, we deduce:
\[
\left( \mathds{1}_\zeta \circ \frac{\partial^\alpha}{\partial x^\alpha} \right) \left( \frac{1}{\alpha!} (x - \zeta)^\beta \right) =
\begin{cases}
1 & \text{if } \alpha = \beta, \\
0 & \text{otherwise}.
\end{cases}
\]
Thus, the set $ \mathcal{B}_\zeta := \left\{ \frac{1}{\alpha!}(x - \zeta)^\alpha \right\}_{\alpha \in \mathbb{N}^n} $ is a basis for $ R $, with dual basis $ \{\partial^\alpha_\zeta\} \subset R^* $, where:
\[
\partial^\alpha_\zeta(p) := \left( \mathds{1}_\zeta \circ \frac{\partial^\alpha}{\partial x^\alpha} \right)(p),
\]

and monomials in $\partial_{i,\zeta}(p)$ in the natural way: $$\partial^{\alpha}_{\zeta}(p):=\partial^{\alpha_1}_{1,\zeta}\ldots\partial^{\alpha_n}_{n,\zeta}(p)=\left(\mathds{1}_{(\zeta_1,\ldots,\zeta_n)} \circ x^{\alpha}(\partial)\right)(p).$$

Then there is an isomorphism 
\[ S^*\to \K[[\partial_{1, \zeta},\ldots,\partial_{n, \zeta}]]\]
\begin{equation*}
\Lambda \mapsto \sum_{\alpha\in \mathbb N^{n+1}} \partial^{\alpha}_{\zeta}(x^{\alpha}) \partial^{\alpha}_{\zeta}.
\end{equation*}
Note that \Cref{lemma:DualChangeBasis} shows the dual change of basis from $\left\{Y^{(\alpha)}\right\}_{\alpha\in \mathbb N^n}$ to the dual basis of $\mathcal{B}_\zeta$. Moreover, this formalism is essentially still valid when $\K$ has positive characteristic, but replacing the basis $\mathcal{B}_\zeta$ of $R$ by $\left\{(x-\zeta)^\alpha\right\}_{\alpha\in \mathbb N^n}$.


\end{comment}


\subsubsection{Comultiplication and divided powers}

The diagonal map $ v \mapsto (v,v) $ on $ V $ induces a comultiplication on $ S $, defined by:
\[
\Delta(x_i) = x_i \otimes 1 + 1 \otimes x_i, \qquad
\Delta(x^\alpha) = \sum_{\beta + \gamma = \alpha} \binom{\alpha}{\beta} x^\beta \otimes x^\gamma.
\]
The multiplication on $ S^* $ is defined by:
\[
(\Lambda_1 \cdot \Lambda_2)(p) := (\Lambda_1 \otimes \Lambda_2)(\Delta(p)).
\]
In this setting:
\[
Y^{(\alpha)} \cdot Y^{(\beta)} = \binom{\alpha + \beta}{\alpha} Y^{(\alpha + \beta)}, \qquad
Y_0^{\alpha_0} \cdots Y_n^{\alpha_n} = \alpha! Y^{(\alpha)}.
\]
One also defines, for any $ F \in S^* $ and $ k \geq 1 $,
$F^{(k)} := \frac{1}{k!} F^k$.


The ring $\mathbb{K}_{dp}[Y_0,\ldots,Y_n]$ is called the \emph{divided power ring}, and it is isomorphic to a polynomial ring in characteristic zero via $Y^{(\alpha)} \mapsto \frac{1}{\alpha!} Z^\alpha$.


%This endows $ S^* $ with a commutative algebra structure. The subalgebra of polynomials in the variables $ Y_i $ is called the ring of \emph{divided powers}, denoted $ \K_{dp}[Y_0,\ldots,Y_n] $, and the corresponding formal power series ring is $ \K_{dp}[[Y_0,\ldots,Y_n]] \cong S^* $.  
%In characteristic zero, there is an isomorphism of rings
%\[
%\K_{dp}[Y_1,\ldots,Y_n] \longrightarrow \K[Z_1,\ldots,Z_n], \qquad Y^{(\alpha)} \mapsto \frac{1}{\alpha!} Z^\alpha.
%\]

%One can check that this inner product defines a structure of commutative algebra in $S^*$. Its subring of finite polynomials in $Y_0,\ldots, Y_n$ is known as \emph{divided powers}, which are denoted $\K_{dp}[Y_0,\ldots,Y_n]$. See \cite{Eisenbud}, Appendix 2 for more details. By a slight abuse of terminology and notation, we will refer to an infinite power series in $Y_0,\ldots,Y_n$ as being expressed in the divided powers basis, and we use $\K_{dp}[[Y_0,\ldots,Y_n]]$ to denote the power series ring identified with $S^*$, with the multiplication defined above. In characteristic 0, the divided powers ring is isomorphic to the usual polynomial algebra:


%\begin{proposition}[{\cite[Proposition 3.13]{Joachim}}]
%There is an isomorphism of rings:
%\[
%\K_{dp}[Y_1,\ldots,Y_n] \longrightarrow \K[Z_1,\ldots,Z_n], \qquad Y^{(\alpha)} \mapsto \frac{1}{\alpha!} Z^\alpha.
%\]
%\end{proposition}


%\begin{example}
%The polynomial $ Z_1^3 + 2Z_2^2 \in \K[Z_1,Z_2] $ corresponds to $ 6Y_1^{(3)} + 4Y_2^{(2)} \in \K_{dp}[Y_1,Y_2] $.
%\end{example}


\subsubsection{Grading and homogeneous coordinate changes}

The decomposition $ S = \bigoplus_{d \geq 0} S_d $ induces a grading on the divided power ring, where:
\[
\mathrm{Hom}_{\mathbb{K}}(S_d, \mathbb{K})* \cong \K_{dp}[Y_0,\ldots,Y_n]_d, \qquad
\K_{dp}[Y_0,\ldots,Y_n] = \bigoplus_{d \geq 0} \mathrm{Hom}_{\mathbb{K}}(S_d, \mathbb{K}).
\]
Note that $S^*=\prod_d \mathrm{Hom}_{\mathbb{K}}(S_d, \mathbb{K})$. Now take $ \zeta = (\zeta_1, \ldots, \zeta_n) \in \K^n $, and define a homogeneous change of coordinates:

\begin{equation}\label{change:of:coordinates}
\varphi(x_0) = x_0, \qquad \varphi(x_i) = x_i + \zeta_i x_0 \quad \text{for } i = 1, \ldots, n.
\end{equation}

The dual map $ \varphi^*: S^* \to S^* $ preserves degrees, and acts on each graded piece $ S_d^* \cong \K_{dp}[Y_0,\ldots,Y_n]_d $.


\begin{example}\label{example: DualChangeCoord}
Let $ n = 1 $, $ \zeta = 2 $. Then the map $ \varphi: \K[x_0,x_1]_2 \to \K[x_0,x_1]_2 $ has matrix:
\[
\begin{bmatrix}
1 & 2 & 4 \\
0 & 1 & 4 \\
0 & 0 & 1
\end{bmatrix}
\]
in the monomial basis. Its transpose gives the matrix of $ \varphi^*: \K_{dp}[Y_0,Y_1]_2 \to \K_{dp}[Y_0,Y_1]_2 $.
\end{example}


%The map $\varphi^*$ is also an algebra homomorphism in $K_{dp}[Y_0,\ldots,Y_n]$ with the canonical product defined before.


Furthermore, due to the functoriality of the symmetric algebra, a homogeneous change of coordinates $\varphi$ in $S$ induces an algebra homomorphism 
\begin{equation}\label{prop: DualChangeIsAlgebraHom}
    \varphi^*: \mathbb{K}_{dp}[Y_0,\ldots,Y_n] \to \mathbb{K}_{dp}[Y_0,\ldots,Y_n]
    \end{equation}
on the divided power algebra (cf. \cite[Chap. III, \S11.5]{Bourbaki}). In particular, we can use the multiplicity of changes of coordinates to avoid computing the matrices in \Cref{example: DualChangeCoord}.
\begin{proposition}\label{prop: ChangeCoordInDivPowers}
    Let $\zeta=(\zeta_1,\ldots, \zeta_n)\in \mathbb K^n$ be an affine point, $\varphi:S\to S$ the homogeneous change of coordinated sending $[1:0:\ldots :0]$ to $[1:\zeta_1:\ldots : \zeta_n]$. Let $\Psi$ be the algebra ismomorphism converting divided powers expressions into polynomials. Then the algebra homomorphism $\Tilde{\varphi}:S\to S$ induced by $x_0\mapsto x_0+\zeta_1x_1+\ldots + \zeta_n x_n$, $x_i\mapsto x_i$ makes the following diagram commute:

\[\begin{tikzcd}
	{K_{dp}[Y_0,\ldots, Y_n]} & {K_{dp}[Y_0,\ldots, Y_n]} \\
	S & S
	\arrow["{\varphi^*}", from=1-1, to=1-2]
	\arrow["\Psi", from=1-1, to=2-1]
	\arrow["\Psi", from=1-2, to=2-2]
	\arrow["{\Tilde{\varphi}}", from=2-1, to=2-2]
\end{tikzcd}\]    
\end{proposition}


\begin{comment}
    

\begin{proposition}\label{prop: DualChangeIsAlgebraHom}
Let $ \varphi $ be the homogeneous change of coordinates as \eqref{change:of:coordinates}. Then the dual map $ \varphi^*: \K_{dp}[Y_0,\ldots,Y_n] \to \K_{dp}[Y_0,\ldots,Y_n] $ is an algebra homomorphism.
\end{proposition}


\begin{proof}
We must check that for all $ \Lambda, \Upsilon \in S^* $ and all $ p \in S $,
\[
\varphi^*(\Lambda \cdot \Upsilon)(p) = (\varphi^*(\Lambda) \cdot \varphi^*(\Upsilon))(p).
\]
The left-hand side is:
\[
\varphi^*(\Lambda \cdot \Upsilon)(p) = (\Lambda \cdot \Upsilon)(\varphi(p)) = (\Lambda \otimes \Upsilon)(\Delta(\varphi(p))),
\]
while the right-hand side is:
\[
(\Lambda \circ \varphi \cdot \Upsilon \circ \varphi)(p) = (\Lambda \otimes \Upsilon)((\varphi \otimes \varphi)(\Delta(p))).
\]
Thus, we must verify:
\[
\Delta(\varphi(p)) = (\varphi \otimes \varphi)(\Delta(p)),
\]
i.e., that $ \varphi $ is a coalgebra homomorphism. This follows from the functoriality of the symmetric algebra and can be found in \cite[§11.5]{Bourbaki}.
\end{proof}
\end{comment}


%\begin{proof}
%    We have to check that for every $\Lambda, \Upsilon\in K_{dp}[Y_0,\ldots,Y_n]$ and every $p\in S$,
%    \[\varphi^*(\Lambda\cdot \Upsilon)(p)=(\varphi^*(\Lambda)\cdot \varphi^*(\Upsilon))(p).\]
%    The left hand side is 
%    \[\varphi^*(\Lambda\cdot \Upsilon)(p)=(\Lambda\cdot \Upsilon)(\varphi(p)=(\Lambda\otimes \Upsilon)(\Delta(\varphi)),\]
%    while the right hand side is 
%    \[((\Lambda\circ \varphi)\cdot (\Upsilon\circ \varphi))(p)=((\Lambda\circ \varphi)\otimes (\Upsilon\circ \varphi))(\Delta(p))=((\Lambda\otimes \Upsilon)\circ (\varphi\otimes \varphi))(\Delta(p)).\]
 %   So we are left to prove the equality 
 %   \[\Delta(\varphi(p))=(\varphi\otimes \varphi)(\Delta(p),\]
%    that is, $\varphi$ is a coalgebra homomorphism on $S$. This follows from the universal property of the symmetric algebra, and can be found in \cite[$\S$ 11.5]{Bourbaki}. 
%\end{proof}






\subsection{Macaulay correspondence and Gorenstein algebras}
Define the contraction action of $S$ on $S^*$ as the dual operation to multiplication:
\[
(p \aprod \Lambda)(q) = \Lambda(pq), \quad p,q \in S,\; \Lambda \in S^*.
\]





%We define the \emph{contraction action} as the dual operation to multiplication in the polynomial ring. Given $p \in S = \K[x_0, \ldots, x_n]$, the linear map $q \mapsto pq$ induces the dual map 
%\[
%(p \aprod \Lambda)(q) := \Lambda(pq)
%\quad \text{for all } \Lambda \in S^*, \, q \in S.
%\]
With the choice of coordinates and dual basis introduced above, one can verify that on monomials:

\begin{equation}\label{eq: ContractionOnMonomials}
x^{\alpha} \aprod Y^{(\beta)} =
\begin{cases}
Y^{(\beta - \alpha)} & \text{if } \alpha_i \leq \beta_i \text{ for all } i, \\
0 & \text{otherwise}
\end{cases}.
\end{equation}

Thus, on polynomials in $\K_{dp}[Y_0, \ldots, Y_n]$, contraction acts as (scaled) partial derivation. In fact, contraction and derivation induce two isomorphic $S$-module structures on $\K_{dp}[Y_0, \ldots, Y_n]$ (see \cite[Proposition 3.13]{Joachim}).

\bigskip

We also consider the contraction action of $R = \K[x_1, \ldots, x_n]$ on $R^*$. As seen in \eqref{eq: ContractionOnMonomials}, contraction of two monomials lowers the degree. This is not necessarily true for infinite series in $\K_{dp}[[Y_1, \ldots, Y_n]]$. For example, for any $p \in R$:
\[
p \aprod \mathds{1}_{\zeta} = p(\zeta)\mathds{1}_\zeta,
\]
since for any $q \in R$,
$
(p \aprod \mathds{1}_\zeta)(q) = \mathds{1}_\zeta(pq) = \mathds{1}_\zeta(p) \cdot \mathds{1}_\zeta(q) = p(\zeta) \mathds{1}_\zeta(q).
$

\begin{remark}\label{remark: contrVsorth}
    For all $p \in S$ and $\Lambda \in S^*$, if $p \aprod \Lambda = 0$, then $\Lambda(p) = 0$. Indeed,
    \[
    \Lambda(p) = \Lambda(p \cdot 1) = (p \aprod \Lambda)(1) = 0.
    \]
    Clearly the same holds for $\Lambda \in R^*$ and $p \in R$. However, the converse does not hold: for instance, $Y^{(3)}(x^2) = 0$ (since $Y^{(3)}$ is the dual of $x^3$), but $x^2 \aprod Y^{(3)} = Y$.
\end{remark}

\begin{definition}
    Let $\Lambda \in S^*$, where $S^*$ is viewed as an $S$-module via contraction. We define the \emph{annihilator} of $\Lambda$ as
    \[
    \operatorname{Ann}(\Lambda) := \{ p \in S \mid p \aprod \Lambda = 0 \}.
    \]
    An ideal $I \subseteq S$ is said to be \emph{apolar} to a homogeneous polynomial $F \in \K_{dp}[Y_0, \ldots, Y_n]_d$ if $I \subseteq \operatorname{Ann}(F)$. A zero-dimensional scheme $Z$ is apolar to $F$ if $I(Z)$ is apolar to $F$.
\end{definition}

Note that the definition of contraction is basis-independent, so $\operatorname{Ann}(\Lambda)$ does not depend on the specific representation of $\Lambda$.

The classical Apolarity lemma can be stated as follows.
\begin{lemma}[Apolarity Lemma]\label{apolarity:lemma}%[{\cite[Lemma 27]{BJPR}}]
  An ideal $I\subset S$ is apolar to   $F\in S_d^*$ if and only if
%    
 %   Let $F\in S_d^*$ and $I\subseteq S$ an ideal. The following are equivalent:
  %  \begin{enumerate}[label=\roman*)]
   %     \item $I\subseteq \operatorname{Ann}(F)$.
    %    \item 
    $I_d\subseteq \{p\in S_d\; | \; p\aprod F =0\}$.
    %\end{enumerate}
\end{lemma}
After having introduced the standard dehomogenization isomorphism $\pi: S_d\cong R_{\leq d}$ and its dual defined by: $(\pi^{-1})^*:(S_d)^*\to (R_{\leq d})^*$, s.t.  $  f^*(p):=F(\pi^{-1}(p))$,  for $f^*=F(Y_0=1)$,
we can also state
an affine version of the apolarity criterion. %\cite[Lemma 27]{BJPR} 

\begin{lemma}[Affine Apolarity Criterion]\label{affine:apolarity:criterion}
For $F\in S_d^*$ homogeneous and $f=F(Y_0=1)$, we have:
\begin{enumerate}[label=\roman*),leftmargin=*]
\item   A zero-dimensional scheme $Z\subseteq \mathbb{P}^n$ defined locally by $\operatorname{Spec}(R/I)$ is apolar to $F$ if and only if $f(p)=0$ for every $p\in I$ with degree $\leq d$.
\item Any $\Lambda\in R^*$ extending $f$ defines a scheme $Z=\operatorname{Spec}(R/\ker(H_\Lambda))$ apolar to $F$.
\end{enumerate}
\end{lemma}

\begin{proof}
\begin{enumerate}[label=\roman*),leftmargin=*]
    \item\label{rmk: affineApolarity} \leavevmode
  By \Cref{apolarity:lemma}, $I(Z)\subseteq \operatorname{Ann}(F)$ if and only if $I(Z)_d\subseteq \operatorname{Ann}(F)$. Since contraction and the duality pairing agree on $S_d$, this is equivalent to say that $F(I(Z)_d)=0$, i.e. for every $p\in I\subseteq R$ of degree $\leq d$, $F(\pi^{-1}(p))=0$. This is to say that $(\pi^{-1})^*(F)(p)=0$ for every $p\in I$ of degree $\leq d$, or equivalently $F(Y_0=1)(p)=0$ for every $p\in I$ of degree $\leq d$.
\item\label{remark:extensionsAreApolar} If $Z$ is locally defined by $\ker(H_\Lambda)$ then by what we just proved in the previous item, $I(Z)\subseteq \operatorname{Ann}(F)$ if and only if $f(p)=0$ for all $p\in \ker(H_\Lambda)\subseteq R$ with $p\in R_{\leq d}$. And the last condition holds since $p\in \ker(H_\Lambda)$ implies $\Lambda(p)=0$ and $\Lambda$ extends $f$.
\end{enumerate}
\end{proof}

We further generalize this result explicitly for schemes supported at arbitrary points:

\begin{lemma}
Let $F\in S_d^*$ homogeneous and $Z$ a local Gorenstein scheme supported at $[1:\zeta_1:\dots:\zeta_n]$ with $I(Z)\subseteq \operatorname{Ann}(F)$. Then there exists $\Lambda=\mathds{1}_\zeta\circ H(\partial)$, $H\in R$, such that $Z$ is defined by $\operatorname{Ann}(\Lambda)$ and the form $\Lambda$ coincides with $f$ in degree $\leq d$. In particular, $f$ is the truncation of the infinite series $\Lambda$ in degree $d$. 
\end{lemma}

\begin{proof}
    Let $\phi$ be a change of coordinates in $S$ that maps the point $[1 : \zeta_1 : \dots : \zeta_n]$ to the origin $[1 : 0 : \dots : 0]$. This transformation is induced by the linear map on coordinates $x_0 \mapsto x_0$ and $x_i \mapsto x_i - \zeta_i x_0$ for $i=1, \dots, n$.

This projective transformation $\phi$ induces an affine change of coordinates $\phi'$ on the space of polynomials $R_{\le d}$ (of degree at most $d$), given by $x_i \mapsto x_i - \zeta_i$.

Under these transformations:
\begin{itemize}
    \item The form $F \in S_d^*$ is transformed into $\tilde{F} = \phi^*(F) \in S_d^*$.
    \item The associated affine form $f^*$ is transformed into $\tilde{f}^* = (\pi^{-1})^* \circ \phi^*(F) = (\phi')^* \circ (\pi^{-1})^*(F) \in R_{\le d}$.
    \item The support $Z$ of the ideal is now located at the origin $[1 : 0 : \dots : 0]$.
\end{itemize}

With $Z$ supported at the origin, the hypotheses of \cite[Proof of Proposition 4]{BJPR} are now satisfied. Therefore, the ideal $I(Z)$ is locally defined by $\text{Ann}(g)$, where $g \in K_d[Y_1, \dots, Y_n] \subseteq R^*$ is a polynomial that defines the same form as $\tilde{f}^*$ in $R_{\le d}$.

Now, we reverse the change of coordinates. According to the change of basis \Cref{lemma:DualChangeBasis}, the ideal $I(Z)$ is supported at its original point $[1 : \zeta_1 : \dots : \zeta_n]$. Locally, it is defined by $\text{Ann}(\Lambda)$, where $\Lambda$ is of the form $1_\zeta \circ H(\partial)$ for some polynomial $H \in R$ such that $\Lambda$ extends the form $f$.

Finally, we note that the homogenization of $\text{ker}(H_\Lambda)$ is apolar to $F$ by \Cref{rmk: affineApolarity}.
\end{proof}









\subsubsection{Annihilators of infinite series}

Since evaluation maps $\mathds{1}_\zeta$ are defined only on the affine space, we restrict our discussion to the affine coordinate ring $R$ and its dual $R^*$.

\begin{remark}
    For every polynomial $f \in \K_{dp}[Y_1, \ldots, Y_n]$, one has
    \[
    (x_1, \ldots, x_n)^{\deg f + 1} \subseteq \operatorname{Ann}(f) \subseteq (x_1, \ldots, x_n),
    \]
    hence $R/\operatorname{Ann}(f)$ is a local ring supported at $0$. This need not hold for annihilators of infinite series. For instance, let $\zeta_1, \zeta_2 \in \K^n$ and define $\Lambda := \mathds{1}_{\zeta_1} + \mathds{1}_{\zeta_2} \in R^*$. Then $\mathcal{V}(\operatorname{Ann}(\Lambda)) = \{\zeta_1, \zeta_2\}$, where $\mathcal{V}(I)$ denotes the variety defined by the ideal $I$.

    Indeed, if $p$ vanishes at both $\zeta_1$ and $\zeta_2$, then for all $q \in R$:
    \[
    (p \aprod \Lambda)(q) = \mathds{1}_{\zeta_1}(p)\mathds{1}_{\zeta_1}(q) + \mathds{1}_{\zeta_2}(p)\mathds{1}_{\zeta_2}(q) = 0,
    \]
    so $p \in \operatorname{Ann}(\Lambda)$. Conversely, let $p \in \operatorname{Ann}(\Lambda)$ and let $q \in R$ such that $q(\zeta_1) \neq 0$, $q(\zeta_2) = 0$, then
    \[
    0 = (p \aprod \Lambda)(q) = \mathds{1}_{\zeta_1}(p) \cdot \mathds{1}_{\zeta_1}(q),
    \]
    forcing $p(\zeta_1) = 0$. Repeating the argument for $\zeta_2$ shows that $p$ vanishes at both points.
\end{remark}

By linearity of contraction, the annihilator of $\Lambda \in R^*$ is the kernel of the linear map known as the \emph{Hankel operator} associated to $\Lambda$:
\[
H_{\Lambda} : R \to R^*, \quad p \mapsto p \aprod \Lambda.
\]

\begin{lemma}\label{rmk: changeCoordinatesKernel}
    Let $\zeta \in \K^n$ and let $\varphi$ 
    %$: R \to R$
    be the affine change of coordinates \eqref{affine:tranlsation}. 
%    defined by $\varphi(x^\alpha) = (x - \zeta)^\alpha$
 Then for any $\Lambda \in R^*$:
    \[
    \varphi(\ker H_\Lambda) = \ker H_{(\varphi^{-1})^*(\Lambda)},
    \]
    where $\varphi^*$ denotes the dual change of coordinates \eqref{dual:affine:tranlsation} given by the transpose of $\varphi$.
\end{lemma}

\begin{proof}
    We show both inclusions. Let $\varphi(p) \in \varphi(\ker H_\Lambda)$. Then for all $q \in R$:
    \[
    (\varphi(p) \aprod (\varphi^{-1})^* \Lambda)(q) = \Lambda\left( (\varphi^{-1} \circ \varphi)(p) \cdot \varphi^{-1}(q) \right) = (p \aprod \Lambda)(\varphi^{-1}(q)) = 0.
    \]
    Conversely, if $p \aprod (\varphi^{-1})^* \Lambda = 0$, then for all $q \in R$:
    \[
    (\varphi^{-1}(p) \aprod \Lambda)(q) = \Lambda\left( \varphi^{-1}(p) \cdot \varphi^{-1}(\varphi(q)) \right) = ((\varphi^{-1})^* \Lambda)(p \cdot \varphi(q)) = (p \aprod (\varphi^{-1})^* \Lambda)(\varphi(q)) = 0.
    \]
\end{proof}
\begin{definition}
    Let $A$ be an Artinian $\K$-algebra. Its \emph{canonical module} $\omega_A$ is the dual vector space $\operatorname{Hom}_\K(A, \K)$ with $A$-module structure given by contraction:
    \[
    (p \aprod \Lambda)(q) = \Lambda(pq) \quad \forall p,q \in A,\, \Lambda \in \omega_A.
    \]
    We say that $A$ is \emph{Gorenstein} if $\omega_A$ is isomorphic to $A$ as an $A$-module, i.e., it is generated by a single element (called the \emph{dual generator}).
\end{definition}


Gorenstein Artinian algebras are quotients of Henkel operators with finite rank.

\begin{lemma}[{\cite[\S 3.1]{Polyexp}}]\label{lemma: GorensteinHankel}
    Let $\Lambda \in R^*$ be such that the algebra $R / \ker(H_\Lambda)$ has finite dimension. Then $R / \ker(H_\Lambda)$ is Gorenstein, with dual generator $\Lambda$. Conversely, every Gorenstein Artinian algebra is isomorphic to $R / \ker(H_\Lambda)$ for some $\Lambda \in R^*$.
\end{lemma}

Gorenstein Artinian algebras are typically studied in the local case (supported at $0$), since any Artinian algebra decomposes as a sum of local Artinian algebras and Gorensteinness is a local property. Recall that for an Artinian local ring $(A, \mathfrak{m})$, the largest integer $d$ such that $\mathfrak{m}^d \neq 0$ is called the \emph{socle degree} of $A$.

\begin{theorem}[{\cite[Lemma 1.2]{Iarrobino}}]
    Let $A = R/I$ be a Gorenstein Artinian algebra supported at $0 \in \K^n$ with socle degree $d$. Then $A = R/\ker(H_\Lambda)$ for some polynomial $\Lambda \in \K_{dp}[Y_1, \ldots, Y_n]$ of degree $d$. Conversely, if $\Lambda \in \K_{dp}[Y_1, \ldots, Y_n]$ is a polynomial of degree $d$, then $R / \ker(H_\Lambda)$ is a local Gorenstein algebra supported at $0$, with socle degree $d$. If $A$ is graded, $\Lambda$ can be chosen homogeneous.
\end{theorem}

It is possible to establish a version of this result for arbitrary coordinates: apply  \Cref{rmk: changeCoordinatesKernel}, \Cref{lemma:DualChangeBasis} and obtain the corresponding change of basis in $R^*$.

\begin{corollary}[Generalized Macaulay Correspondence]\label{Thm: GeneralisedMacaulayCorrespondence}
    Let $A$ be a Gorenstein Artinian algebra of socle degree $d$ supported at $\zeta \in \K^n$. Then $A \cong R / \ker(H_\Lambda)$ for some $\Lambda \in R^*$ of the form $\mathds{1}_\zeta \circ H(\partial)$, where $H \in R$ is a polynomial of degree $d$ and $H(\partial)$ denotes the differential operator obtained by replacing each $x_i$ with $\partial / \partial x_i$. Conversely, if $\Lambda = \mathds{1}_\zeta \circ H(\partial)$ with $\deg H = d$, then $R / \ker(H_\Lambda)$ is a local Gorenstein algebra supported at $\zeta$, of socle degree $d$.
\end{corollary}

\subsection{Hilbert Functions of Gorenstein algebras}

The Hilbert function of the associated graded algebra of a local Gorenstein algebra must satisfy certain structural properties, and admits a symmetric decomposition induced by the interplay of two natural filtrations. We recall the basic constructions from \cite{Iarrobino}.

\begin{definition}
    Let $A$ be a local Artinian $\K$-algebra with maximal ideal $\mathfrak{m}$. Its \emph{associated graded algebra} is
    \[
    gr(A) := \bigoplus_{i \geq 0} \mathfrak{m}^{i} / \mathfrak{m}^{i+1}.
    \]
    The \emph{Hilbert function} of $A$ is defined as
    \[
    H_A(i) := \dim_{\K} \left( \mathfrak{m}^{i} / \mathfrak{m}^{i+1} \right).
    \]
\end{definition}

Since $A / \mathfrak{m} \cong \K$, we have $H_A(0) = 1$. The \emph{socle degree} of $A$ is the largest integer $i$ for which $H_A(i) \neq 0$.


As explained in \cite{Iarrobino}, if $A$ is graded, the dual generator of the Gorenstein algebra induces nondegenerate bilinear pairings between graded components $
A_i \times A_{d-i} \to \K$,
where $d$ is the socle degree. This implies that the Hilbert function is symmetric 
$H_A(i) = H_A(d - i)$, 
and in particular, $H_A(d) = 1$. Moreover, in the graded case one has
\[
(0 : \mathfrak{m}^{d - i}) = \mathfrak{m}^{i + 1}.
\]
However, this identity does not hold in general if $A$ is not graded (see e.g., \cite[Example 3.28]{Joachim}). In such cases, two canonical filtrations naturally arise:

\begin{itemize}
    \item the \emph{powers filtration}:
    \[
    \{0\} \subseteq \mathfrak{m}^d \subseteq \mathfrak{m}^{d-1} \subseteq \cdots \subseteq \mathfrak{m} \subseteq A;
    \]
    \item the \emph{Löwy filtration} (see \cite{Iarrobino}):
    \[
    \{0\} \subseteq (0 : \mathfrak{m}) \subseteq (0 : \mathfrak{m}^2) \subseteq \cdots \subseteq (0 : \mathfrak{m}^{d+1}) = A.
    \]
\end{itemize}

These filtrations are related through duality.

\begin{lemma}[{\cite[Lemma 2.31]{Joachim}}] 
    Let $A$ be a local Artinian $\K$-algebra, and let $I \subseteq A$ be an ideal. Let $\Lambda$ be a generator of the canonical module $\omega_A$, and consider the pairing
$    A \times A \to \K, \quad (p,q) \mapsto \Lambda(pq)$.
    Then:
    \[
    (0 : I) = \{ a \in A \mid \Lambda(ai) = 0 \text{ for all } i \in I \}.
    \]
\end{lemma}

In particular, there is a natural duality between the graded pieces of the two filtrations:
\[
\left(\frac{\mathfrak{m}^{i}}{\mathfrak{m}^{i+1}}\right)^* \cong \frac{(0 : \mathfrak{m}^{i+1})}{(0 : \mathfrak{m}^i)},
\]
and therefore:
\[
H_A(i) = \dim_\K \left( \frac{(0 : \mathfrak{m}^{i+1})}{(0 : \mathfrak{m}^i)} \right).
\]

This leads to the following family of subquotients: for each $a = 0, \ldots, d$, define
\[
C_{a,i} := \frac{(0 : \mathfrak{m}^{d+1 - a - i}) \cap \mathfrak{m}^i}{(0 : \mathfrak{m}^{d+1 - a - i}) \cap \mathfrak{m}^{i+1}},
\]
\[
C_a := \bigoplus_{i = 0}^{d - a} C_{a,i}, \quad Q_a := C_a / C_{a+1}.
\]
Let $\Delta_{Q_a}(i) := \dim_\K (Q_a)_i$ denote the Hilbert function of $Q_a$. From the definitions, one sees that:
\[
\Delta_{Q_0}(0) = 1, \quad \Delta_{Q_a}(0) = 0 \text{ for all } a > 0.
\]

The modules $Q_a$ give a symmetric decomposition of the Hilbert function of $A$:
\begin{theorem}[{\cite{Iarrobino}}]\label{thm:HF:decomposition}
    The Hilbert functions $\Delta_{Q_a}$ satisfy the following properties:
    \begin{enumerate}
        \item \textbf{Decomposition:}
        \[
        H_A(i) = \sum_{a = 0}^{d - i} \Delta_{Q_a}(i).
        \]
        
        \item \textbf{Symmetry:} Each $\Delta_{Q_a}$ is symmetric around $(d - a)/2$:
        \[
        \Delta_{Q_a}(i) = \Delta_{Q_a}(d - a - i).
        \]
        
        \item \label{HF: PartialSums} \textbf{Partial sums:} For each $k \geq 0$, the partial sum $\sum_{a = 0}^{k} \Delta_{Q_a}$ is the Hilbert function of a graded $\K$-algebra generated in degree 1.
    \end{enumerate}
\end{theorem}

\noindent
Note that \Cref{HF: PartialSums} in \Cref{thm:HF:decomposition} implies that each partial sum satisfies the Macaulay bound for Hilbert functions (see \cite{Joachim}).

\section{Refining the Cactus Rank Algorithm in Local Settings}\label{sec:main}
This section delves into the core contribution of our work: refining the classic algorithm for computing the cactus rank of a homogeneous polynomial (cf. \cite{Alessandra}) in a local setting. By leveraging the unique properties of local Gorenstein algebras, we develop effective methods to compute key algebraic invariants such as  the Hilbert function and its symmetric decomposition, within the algorithm's framework. Ultimately, these advancements pave the way for a more efficient and insightful computation of the generalized additive decomposition of a polynomial or of its extension, directly revealing the polynomial's cactus rank.
\subsection{The cactus algorithm}\label{sec:algorithms}

In this section, we recall how to compute a minimal apolar schemes to homogeneous polynomials in the global setting (cf. \cite{Alessandra}).

%We recall briefly the basic setting and key definitions from \cite{BJPR, GAD, Bernard, Alessandra} highlighting explicitly the relationship and differences with the cactus algorithm discussed in \cite{Alessandra}.

%Let $S=\K[x_0,\ldots,x_n]$ be the coordinate ring of $\mathbb{P}^n$ and $R=\K[x_1,\ldots,x_n]$ the coordinate ring of the affine chart defined by $x_0\neq 0$. We denote homogeneous polynomials in $S$ or $\K_{dp}[Y_0,\ldots,Y_n]\subseteq S^*$ by uppercase letters, while lowercase letters denote polynomials in $R$ or $\K_{dp}[Y_1,\ldots,Y_n]\subseteq R^*$.

%Consider a homogeneous polynomial $F\in \K_{dp}[Y_0,\ldots,Y_n]_d$ expressed in divided powers. A local Gorenstein scheme $Z=\text{Spec}(A)$, where $A=R/I$ is a local Gorenstein algebra supported at a point $\zeta\in \mathbb{A}^n$, is said to be \emph{apolar} to $F$ if $I(Z)\subseteq \operatorname{Ann}(F)$. By duality, we represent $I$ as the annihilator ideal $\operatorname{Ann}(\Lambda)$ for some $\Lambda\in R^*$. Our primary interest lies in the explicit relationship between $\Lambda$ and $F$.

\begin{comment}
    
% MOVED BELOW

A fundamental result, originally established for local apolarity at $[1:0:\dots:0]$ in \cite[Proof of Proposition 4]{BJPR}, states that if $F$ is homogeneous and $\Gamma=\text{Spec}(R/\operatorname{Ann}(g))$ is minimal among apolar schemes supported at $[1:0:\dots:0]$, then the degree-$d$ tails of $g$ and the dehomogenization $f=F(Y_0=1)$ coincide. This result naturally extends to general points by affine changes of coordinates.
\end{comment}

\begin{comment}
    
% MOVED IN THE PREVIOUS SECTION

\smallskip

The standard dehomogenization isomorphism $\pi: S_d\cong R_{\leq d}$ and its dual $(\pi^{-1})^*$ defined by:

\[(\pi^{-1})^*:(S_d)^*\to (R_{\leq d})^* \quad \quad f^*(p):=F(\pi^{-1}(p)), \]

play a central role.

\end{comment}
\begin{comment}
 %   MOVED IN THE PREVIOUS SECTION
Let us recall the classical \emph{Apolarity Lemma}. %(\cite{BJPR}).

\begin{lemma}[Apolarity Lemma]%[{\cite[Lemma 27]{BJPR}}]
    Let $F\in S_d^*$ and $I\subseteq S$ an ideal. The following are equivalent:
    \begin{enumerate}[label=\roman*)]
        \item $I\subseteq \operatorname{Ann}(F)$.
        \item $I_d\subseteq \{p\in S_d\; | \; \langle p, F\rangle_\aprod=0\}$.
    \end{enumerate}
\end{lemma}
\end{comment}

\begin{comment}

% MOVED IN THE PREVIOSU SECTION

An affine apolarity criterion %\cite[Lemma 27]{BJPR} 
can thus be explicitly stated as follows:

\begin{lemma}[Affine Apolarity Criterion]
For $F\in S_d^*$ homogeneous and $f=F(Y_0=1)$, we have:
\begin{enumerate}[label=\roman*)]
\item A zero-dimensional scheme $Z\subseteq \mathbb{P}^n$ defined locally by $\text{Spec}(R/I)$ is apolar to $F$ if and only if $f(p)=0$ for every $p\in I$ with degree $\leq d$.
\item Any $\Lambda\in R^*$ extending $f$ defines a scheme $Z=\text{Spec}(R/\ker(H_\Lambda))$ apolar to $F$.
\end{enumerate}
\end{lemma}

\begin{proof}
\begin{enumerate}[label=\roman*)]
    \item\label{rmk: affineApolarity} \leavevmode
  By the Apolarity Lemma $I(Z)\subseteq \operatorname{Ann}(F)$ if and only if $I(Z)_d\subseteq \operatorname{Ann}(F)$. Since contraction and the duality pairing agree on $S_d$, this is equivalent to say that $F(I(Z)_d)=0$, i.e. for every $p\in I\subseteq R$ of degree $\leq d$, $F(\pi^{-1}(p))=0$. This is to say that $(\pi^{-1})^*(F)(p)=0$ for every $p\in I$ of degree $\leq d$, or equivalently $F(Y_0=1)(p)=0$ for every $p\in I$ of degree $\leq d$.
\item\label{remark:extensionsAreApolar} If $Z$ is locally defined by $\ker(H_\Lambda)$ then by what we justproved in the previous item, $I(Z)\subseteq \operatorname{Ann}(F)$ if and only if $f^*(p)=0$ for all $p\in \ker(H_\Lambda)\subseteq R$ with $p\in R_{\leq d}$. And the last condition holds since $p\in \ker(H_\Lambda)$ implies $\Lambda(p)=0$ and $\Lambda$ extends $f^*$.
\end{enumerate}
\end{proof}

We further generalize this result explicitly for schemes supported at arbitrary points:

\begin{lemma}
Let $F\in S_d^*$ homogeneous and $Z$ a local Gorenstein scheme supported at $[1:\zeta_1:\dots:\zeta_n]$ with $I(Z)\subseteq \operatorname{Ann}(F)$. Then there exists $\Lambda=\mathds{1}_\zeta\circ H(\partial)$, $H\in R$, such that $I(Z)=\operatorname{Ann}(\Lambda)$ and the form defined by $\Lambda$ coincides with $f$ in degree $\leq d$. In particular, $f$ is the truncation of the infinite series $\Lambda$ in degree $d$. 
\end{lemma}

\begin{proof}
    Let $\phi$ be a change of coordinates in $S$ that maps the point $[1 : \zeta_1 : \dots : \zeta_n]$ to the origin $[1 : 0 : \dots : 0]$. This transformation is induced by the linear map on coordinates $x_0 \mapsto x_0$ and $x_i \mapsto x_i - \zeta_i x_0$ for $i=1, \dots, n$.

This projective transformation $\phi$ induces an affine change of coordinates $\phi'$ on the space of polynomials $R_{\le d}$ (of degree at most $d$), given by $x_i \mapsto x_i - \zeta_i$.

Under these transformations:
\begin{itemize}
    \item The form $F \in S_d^*$ is transformed into $\tilde{F} = \phi^*(F) \in S_d^*$.
    \item The associated affine form $f^*$ is transformed into $\tilde{f}^* = (\pi^{-1})^* \circ \phi^*(F) = (\phi')^* \circ (\pi^{-1})^*(F) \in R_{\le d}$.
    \item The support $Z$ of the ideal is now located at the origin $[1 : 0 : \dots : 0]$.
\end{itemize}

With $Z$ supported at the origin, the hypotheses of \cite[Proof of Proposition 4]{BJPR} are now satisfied. Therefore, the ideal $I(Z)$ is locally defined by $\text{Ann}(g)$, where $g \in K_d[Y_1, \dots, Y_n] \subseteq R^*$ is a polynomial that defines the same form as $\tilde{f}^*$ in $R_{\le d}$.

Now, we reverse the change of coordinates. According to the change of basis \Cref{lemma:DualChangeBasis}, the ideal $I(Z)$ is supported at its original point $[1 : \zeta_1 : \dots : \zeta_n]$. Locally, it is defined by $\text{Ann}(\Lambda)$, where $\Lambda$ is of the form $1_\zeta \circ H(\partial)$ for some polynomial $H \in R$ such that $\Lambda$ extends the form $f$.

Finally, we note that the homogenization of $\text{ker}(H\Lambda)$ is apolar to $F$ by \Cref{rmk: affineApolarity}.
\end{proof}
\end{comment}
%\subsection{On Gorenstein local apolar schemes}

%We recall that $S=\K[x_0,\ldots,x_n]$ is the coordinate ring of $\mathbb P^n$ and $R=\K[x_1,\ldots,x_n]$ is the coordinate ring of the affine chart $x_0\neq 0$. For the reader's convenience, we will use capital letters to denote homogeneous polynomials in $S$ or $\K_{dp}[Y_0,\ldots, Y_n]\subseteq S^*$, and lower case letters for polynomials in $R$ or $\K_{dp}[Y_1,\ldots, Y_n]\subseteq R^*$.

%\bigskip

%Consider $F\in \K_{dp}[Y_0,\ldots, Y_n]_d$ a homogeneous polynomial in divided powers of degree $d$ and $Z$ a local Gorenstein scheme apolar to $F$, i.e., $I(Z)\subseteq Ann(F)=\ker(H_F)$, and $Z=\text{Spec}(A)$ for some $A=R/I$ local Gorenstein algebra supported at $\zeta$. By the correspondence seen in the previous section, $I=Ann(\Lambda)$ for some $\Lambda\in R^*$. We are interested in the relationship between $\Lambda$ and $F$. 

%\bigskip

%The case where $\zeta=0$ was explored in \cite{BJPR}:

%\begin{proposition}[{\cite[Proof of Proposition 4]{BJPR}}]\label{Prop: ExtensionAt0}
%    Let $F$ be a homogeneous polynomial of degree $d$, and let $f= F(Y_0=1)$. Let $\Gamma$ be a zero-dimensional scheme apolar to $F$ minimal among those supported at $[1 : 0 :\ldots: 0]$. Then $\Gamma=\text{Spec}(R/Ann(g))$ for some polynomial $g$ and the degree $d$ tails of $g$ and $f$ agree.     
%\end{proposition}

%Here we will prove this result for a fixed system of coordinates, that is, for $\Gamma$ supported at a general $[1:\zeta_1:\ldots:\zeta_n]$. A key observation is that in \Cref{Prop: ExtensionAt0} above, $g$ is a form in $R^*$ that agrees with $f\in R^*$ in $R_{\leq d}$.  

%\bigskip

%For any $d\geq 1$, contraction defines a natural perfect pairing 
%\[S_d \times (S_d)^*\to \K \quad \quad (p, \Lambda)\mapsto \langle p, \Lambda\rangle_{\aprod}:= p\aprod \Lambda.\]

%The following result, known as \emph{apolarity lemma} (\cite{BJPR}), shows that the annihilator of a homogeneous polynomial of degree $d$ is in some sense determined by this pairing:

%\begin{proposition}[{\cite[Lemma 27]{BJPR}}]
%    Let $F\in S_d^*$ and $I\subseteq S$ an ideal. The following are equivalent:
%    \begin{itemize}
%        \item[i)] $I\subseteq Ann(F)$.
%        \item[ii)] $I_d\subseteq \{p\in S_d\; | \; \langle p, F\rangle_\aprod=0\}$.
%    \end{itemize}
%\end{proposition}
%
%Dehomogenization gives us an isomorphism of vector spaces $\pi:S_d\cong R_{\leq d}$, with inverse map 
%\[\pi^{-1}:R_{\leq d}\to S_d \quad \quad p\mapsto x_0^d p(\frac{x_1}{x_0}, \ldots, \frac{x_n}{x_0}).\]
%
%The induced dual isomorphism is 
%\[(\pi^{-1})^*:(S_d)^*\to (R_{\leq d})^* \quad \quad f^*(p):=F(\pi^{-1}(p)). \]
 
%\begin{remark}\label{rmk:ProductAndDehom}
%    A direct check shows that if $F\in S_d^*$, then $(\pi^{-1})^*(F)=F(Y_0=1)$. Note however that this dual dehomogenization is not an algebra homomorphism. For example, take $F=Y_0^{(2)}$ and $G=Y_0Y_1$:
%    \[(\pi^{-1})^*(F\cdot G)=(\pi^{-1})^*(3Y_0^{3}Y_1)=3Y_1\neq Y_1=(\pi^{-1})^*(F)\cdot (\pi^{-1})^*(G).\]
%\end{remark}

%The apolarity lemma and dehomogenization gives us a powerful tool to produce affine schemes that are apolar to a given homogeneous polynomial:
%\begin{lemma}
%    Let $F\in S_d^*$ be a homogeneous polynomial and let $f=(\pi^{-1})^*(F)=F(Y_0=1)\in R_{\leq d}^*$.
%    \begin{itemize}
%        \item[i)]\label{rmk: affineApolarity} Let $Z\subseteq \mathbb P^n$ be a zero dimensional scheme defined locally by $\text{Spec}(R/I)$.Then $Z$ is apolar to $F$ if and only if $f(p)=0$ for every $p\in I$ of degree $\leq d$. 
 %       \item[ii)]\label{remark:extensionsAreApolar} Any $\Lambda\in R^*$ extending $f$ (that is, for every $p\in R_{\leq d}$ $\Lambda(p)=f(p)$) defines a zero dimensional scheme $Z=\text{Spec}(R/\ker(H_\Lambda))$ such that $I(Z)$ is apolar to $F$.
%    \end{itemize}
%\end{lemma}
%\begin{proof}
%     For the first statement, we have that $I(Z)\subseteq Ann(F)$ if and only if $I(Z)_d\subseteq Ann(F)$ (by the apolarity lemma), if and only if $F(I(Z)_d)=0$ (since contraction and the duality pairing agree on $S_d$), if and only if, for every $p\in I\subseteq R$ of degree $\leq d$, $F(\pi^{-1}(p))=0$, if and only if $(\pi^{-1})^*(F)(p)=0$ for every $p\in I$ of degree $\leq d$, if and only if $F(Y_0=1)(p)=0$ for every $p\in I$ of degree $\leq d$.
%
%    \bigskip
%
%    For the second, if $Z$ is locally defined by $\ker(H_\Lambda)$ then by \Cref{rmk: affineApolarity} $I(Z)\subseteq Ann(F)$ if and only if $f^*(p)=0$ for all $p\in \ker(H_\Lambda)\subseteq R$ with $p\in R_{\leq d}$. And the last condition holds since $p\in \ker(H_\Lambda)$ implies $\Lambda(p)=0$ and $\Lambda$ extends $f^*$.
%
% \end{proof}

%Now we can prove a slightly more general result than \Cref{Prop: ExtensionAt0}:
%\begin{lemma}
 %   Let $F\in S_d^*$ be a homogeneous polynomial, $f=(\pi^{-1})^*(F)=F(Y_0=1)\in R_{\leq d}^*$ and let $Z$ be a local Gorenstein scheme supported at $[1:\zeta_1:\ldots : \zeta_n]$ such that $I(Z)\subseteq Ann(F)$. Then $I(Z)$ is defined locally by $Ann(\Lambda)$ for some $\Lambda\in R^*$ of the type $\Lambda=\mathds{1}_\zeta \circ H(\partial)$, with $H\in R$, and such that $f$ and $\Lambda$ define the same form in $R_{\leq d}$. In particular, $f$ is the truncation of the infinite series $\Lambda$ in degree $d$.  
%\end{lemma}

%\begin{proof}
%    Let $\varphi$ be the change of coordinates in $S$ sending $[1:\zeta_1:\ldots : \zeta_n]$ to $[1:0:\ldots : 0]$, induced by $x_0\mapsto x_0$, $x_i\mapsto x_i-\zeta_ix_0$ for $i=1,\ldots, x_n$. The induced affine change of coordinates $\varphi'$ on $R_{\leq d}$ is given by $x_i\mapsto x_i-\zeta_i$. Then $F\in S_d^*$ is transformed into $\Tilde{F}=\varphi^*(F)\in S_d^*$, $f^*$ is transformed into $\Tilde{f^*}=(\pi^{-1})^*\circ \varphi^*(F)=(\varphi')^*\circ (\pi^{-1})^*(F)\in R_{\leq d}$, and $Z$ is now supported at $[1:0:\ldots:0]$.
 %   
 %   The hypothesis of \Cref{Prop: ExtensionAt0} now hold, therefore $I(Z)$ is defined locally by $Ann(g)$ with $g\in \K_{dp}[Y_1,\ldots,Y_n]\subseteq R^*$ such that
   %  it defines the same form as $\Tilde{f^*}$ in $R_{\leq d}$. We now undo the change of coordinates, so that according to the change of basis lemma  \Cref{lemma:DualChangeBasis} $I(Z)$ is supported at $[1:\zeta_1:\ldots:\zeta_n]$ and it is defined locally by $Ann(\Lambda)$ with $\Lambda$ of the type $\mathds{1}_{\zeta}\circ H(\partial)$, for some $H\in R$ such that $\Lambda$ extends $f$. Note that (the homogenization of) $\ker(H_\Lambda)$ is apolar to $F$ by \Cref{rmk: affineApolarity}.  
    % 
    %\end{proof}

%-----------------
%Motivated by these considerations, we 
%Recall the following  definitions from \cite{GAD, IK}.

\begin{definition}[\cite{BR13}]
For $F\in S_d^*$, the \emph{cactus rank} of $F$ is the minimal length of a scheme apolar to $F$. The \emph{local cactus rank} is the minimal length of a local apolar scheme.
\end{definition}


%Among schemes that are apolar to a given $F\in S_d^*$, we are interested in those having minimum length. 
%\begin{definition}
%    $F\in S_d^*$, we define the \emph{cactus rank of} $F$ as the minimal length of a scheme apolar to $F$, called the \emph{cactus scheme of} $F$. We call the \emph{local cactus rank of} $F$ as the minimal length of a local scheme apolar to $F$.      
%\end{definition}

\begin{comment}
% THIS PART IS LOCAL. I MOVED IT IN THE LOCAL SECTION 
The construction of local apolar schemes explicitly considered in \cite{GAD} for polynomials of type \begin{equation}\label{GADLocal}
    F=L^{(k)}\cdot G,
\end{equation} $L\in S_1^*$, $G\in S_{d-k}^*$, through dehomogenization with respect to $L$, will be our starting point. Such schemes, called \emph{evinced}, allow an explicit and constructive description of minimal local apolar schemes.


%As it was shown in \cite{BJPR}, \cite{GAD}, there is a natural way to define a local scheme that is apolar to a homogeneous polynomial $F\in S^*_d$ of the form 

%\begin{equation}\label{GADLocal}
%    F=L^{(k)}\cdot G
%\end{equation}
%where $L\in S_1^*$ and $G\in S_{d-k}^*$. For this construction, first we dehomogenize $F$ with respect to $L$, that is, we consider $G_L$ as the image of the projection $S^*\to S^*/(L-1)$, and then define $Z_{F,L}$ as the scheme that is defined locally by the annihilator of $G_L$. The resulting scheme is said to be \emph{evinced} by Eq.\eqref{GADLocal}. 
A fundamental result, originally established for local apolarity at $[1:0:\dots:0]$ in \cite[Proof of Proposition 4]{BJPR}, states that if $F$ is homogeneous and $\Gamma=\text{Spec}(R/\operatorname{Ann}(g))$ is minimal among apolar schemes supported at $[1:0:\dots:0]$, then the degree-$d$ tails of $g$ and the dehomogenization $f=F(Y_0=1)$ coincide. This result naturally extends to general points by affine changes of coordinates.



Another way of stating {\cite[Proof of Proposition 4]{BJPR}} is that a zero-dimensional scheme of minimal length among local schemes supported at
$[1 : 0 : \cdots: 0]$ that are apolar to $F$ is evinced by 
\[G=Y_0^{(d'-k)}\cdot H\in S_{d'}^*\]
where $d'\geq d$, $k\in \{0,\ldots,d'\}$ and $(\pi^{-1})^*(G)=G(Y_0=1)$ has the same degree $d$-tail as $f$. We will return to this construction of the scheme in \Cref{Section:GAD}.

\end{comment}

The strategy developed in \cite{Alessandra} for the global cactus rank algorithm relies on the Hankel operator $H_\Lambda: R \to R^*$ defined as $H_\Lambda(p)(q)=\Lambda(pq)$. We recall here the crucial algebraic conditions from \cite{Bernard} for extending a truncated linear form $f\in R^*_{\leq d}$ to a form $\Lambda\in R^*$, necessary for defining finite-dimensional local algebras.


\begin{definition}
Let $V^* \subset R^*$ be two spaces equipped with a restriction map
\[
\mathrm{res} : R^* \longrightarrow V^*.
\]
Given $\Lambda \in V^*$, we say that $\tilde \Lambda \in R^*$ \emph{extends} $\Lambda$ if
\[
\mathrm{res}(\tilde \Lambda) = \tilde \Lambda\big|_{V^*} = \Lambda.
\]
\end{definition}
%In other words, $\tilde \Lambda$ is an extension of $\Lambda$ if it coincides with $\Lambda$ when restricted to the corresponding subdomain (or substructure).

\begin{theorem}[\cite{Bernard}]
Given a monomial basis $B$ connected to $1$ and a truncated linear form $\Lambda\in \langle B\cdot B^+\rangle_{\leq d}^*$, an extension $\Tilde{\Lambda}\in R^*$ exists with $rank(H_{\Tilde{\Lambda}})=|B|$ and $B$ basis of $R/\ker(H_{\Tilde{\Lambda}})$ if and only if the associated matrices of multiplication by the variables commute and $\det(H_{\Lambda}^B)\neq 0$. Such extension is unique.
\end{theorem}

%These known results form the foundation for our main contribution: an explicit algorithm to compute local cactus ranks, extending and refining the global approach developed in \cite{Alessandra}. 

The essential ingredients used in \cite{Alessandra} for the cactus algorithm based on the previous result are:

\begin{itemize}
    \item (Cf.\cite{Bernard}) The fact that the multiplication operators $H_{a\thinstar \Lambda}$ for any $\Lambda\in R^*$ such that $\operatorname{rank} \; H_{\Lambda}<\infty$ and any $a\in A_{\Lambda}$ can be computed as
    \begin{equation}\label{MatMult}
    H_{a\thinstar \Lambda}=M_a^t \circ H_{\Lambda};\end{equation} 

    \item 
 (Cf. {\cite[Theorem 4.23]{BernardBook}}) \label{Thm: eigenvaluesMultOperatos}
        For every $a\in A_{\Lambda}$, the eigenvalues of the operators $M_a$ and $M_a^t$ are $\{a(\zeta_i), \ldots, a(\zeta_r)\}$ and the common eigenvectors of $M_{x_i}^t$, $i=1,\ldots,n$ are, up to scalar, $\mathds{1}_{\zeta_i}$.
    \item (Cf. {\cite[Proposition 3.2]{Alessandra}})\label{Prop: BasisInHankel}
    If $B\subseteq R/\ker H_\Lambda$ with $|B|=r$, then $B$ is a basis of the vector space $A_{\Lambda}$ if and only if $H_{\Lambda}^B$ is invertible and $\operatorname{rank} \; H_{\Lambda}=r$. Moreover, there exists a monomial basis $B$ that is a complete staircase.
\end{itemize}

    








%In \cite{GAD}, the authors explored the relation between minimal apolar schemes to a polynomial $F\in S_d^*$ and schemes evinced by sum of terms as in Eq.\eqref{GADLocal},  
%\[\sum_i^r L_i^{k}G_i,\]
%called generalized additive decomposition (GAD) of $F$ (\cite{IK}). 
%A necessary and sufficient condition for the minimal apolar scheme to be evinced by a GAD of $F$ is that it is contained in the fat points of the correct size (see \cite[Proposition 4.3]{GAD}, \cite[Example 4.3]{GAD} and for the local case).


\begin{comment}
\subsection{Extensions and finite Hankel algebras}
So we want to find a form in $R^*$ that extends $f^*$ and that verifies two conditions: the dimension of $R/\ker(H_{\Lambda})$ is minimal and it defines a local algebra. In this Section we will review the results related to these requirements.

\bigskip

First, having $\dim_{\K}(R/\ker(H_{\Lambda}))<\infty$ is not enough to guarantee that it is local, as the next result shows: 

\begin{proposition}\label{EigenValues}
    Let $\Lambda\in R^*$ such that $R/\ker(H_{\Lambda})$ is a finite dimensional vector space. Then \begin{equation}\label{SumDifs}
        \Lambda= \sum_{i=1}^r \mathds{1}_{\zeta_i} \circ p_i(\partial)
    \end{equation} for some $\zeta_i\in \mathbb K^n$ and $p_i(\partial)\in \mathbb K[\partial_1,\ldots,\partial_n]$ 
\end{proposition}


We are interested in a computational way to read the points of \eqref{SumDifs}. Let $A=R/I$ be an Artinian algebra and let $a\in A$. We define the multiplication operator $M_a:A\to A$ as $p\mapsto ap$. 

\begin{theorem}[{\cite[Theorem 4.23]{BernardBook}}]\label{Thm: eigenvaluesMultOperatos}
        \begin{itemize}
        \item[i)] For every $a\in A_{\Lambda}$, the eigenvalues of the operators $M_a$ and $M_a^t$ are $\{a(\zeta_i), \ldots, a(\zeta_r)\}$.
        \item[ii)] The common eigenvectors of $M_{x_i}^t$, $i=1,\ldots,n$ are, up to scalar, $\mathds{1}_{\zeta_i}$.
    \end{itemize}
    
\end{theorem}

We see that if we have the multiplication operators of the coordinate polynomials, $M_{x_i}$ for $i=1,\ldots,n$, the $i$-th components of $\zeta_1,\ldots, \zeta_r$ are the eigenvalues of $M_{x_i}$. These multiplication operators can be obtained from the Hankel of $\Lambda$:

\begin{lemma}[{\cite{Bernard}}]
    For any $\Lambda\in R^*$ such that $rank \; H_{\Lambda}<\infty$ and any $a\in A_{\Lambda}$, 
    \begin{equation}\label{MatMult}
    H_{a\thinstar \Lambda}=M_a^t \circ H_{\Lambda}\end{equation}
\end{lemma}

\bigskip

Since the rank of the Hankel of $\Lambda$ is finite, we should be able to find a basis of the quotient algebra $R/\ker(H_{\Lambda})$ and associate a matrix to $M_{x_i}$. The first step is to restrict the Hankel to a finite dimensional subspace. Given $B=\{b_1,\ldots,b_r\}, B'=\{b_1',\ldots,b'_{r'}\}\subseteq R$, we define $$H_{\Lambda}^{B,B'}: \langle B \rangle_{\mathbb K}\to \langle B' \rangle_{\mathbb K}^*$$
the restriction of $H_{\Lambda}$ to $\langle B \rangle_{\mathbb K}$ and projection of $R^*$ into $\langle B' \rangle_{\mathbb K}^*$. In this case, the matrix of $H_{\Lambda}^{B,B'}$ is 
$$(H_{\Lambda}^{B,B'})_{i,j}=\Lambda(b_i\cdot b_j') \quad \quad i=1,\ldots,r, \;j=1,\ldots,r'$$ If $B=B'$, we will use $H_{\Lambda}^{B,B'}=H_{\Lambda}^{B}$.

\bigskip

Next, let us see the conditions that a basis needs to verify. 

\begin{definition}
    Let $B\subseteq R$ be a set of monomials. We say $B$ is connected to 1 if $\forall m\in B$ either $m=1$ or there is $1\leq i\leq n$ and $m'\in B$ such that $m=x_im'$. We say that $B$ is a staircase if for all $i=1,\ldots,n$ $x_ix^{\beta}\in B$ implies $x^{\beta}\in B$. If in addition $B$ contains all the degree one monomials, we say that $B$ is a complete staircase.
\end{definition}



\begin{proposition}[{\cite[Proposition 3.2]{Alessandra}}]\label{Prop: BasisInHankel}
    Let $B\subseteq R/\ker H_\Lambda$ with $|B|=r$. Then $B$ is a basis of the vector space $A_{\Lambda}$ if and only if $H_{\Lambda}^B$ is invertible and $rank \; H_{\Lambda}=r$. Moreover, there exists a monomial basis $B$ that is a complete staircase.
\end{proposition}





\begin{remark}\label{rmk:computeSupports}
    If $\Lambda\in R^*$ with $R/\ker(H_\Lambda)$ an Artinian algebra of $\K$-dimension $r$, then as we have seen $\Lambda$ is of the type \eqref{SumDifs}, and we can calculate the supports $\zeta_1,\ldots,\zeta_n$ by first finding a complete staircase of monomials of $r$ elements such that the matrix $H_\Lambda^B$ is invertible, then use \eqref{MatMult} to obtain the matrices corresponding to the multiplication by $x_i$ in the algebra $R/\ker H_\Lambda$, and from their eigenvalues read the coordinates of $\zeta_1,\ldots,\zeta_r$.
\end{remark}


    
As we have seen, we look for an extension of $f^*\in R_{\leq d}^*$. We consider the "infinite matrix"  in the monomial basis of the Hankel operator of a generic extension $\Lambda\in R^*$:

\[(H_{\Lambda})_{\alpha, \beta}=\begin{cases}
    f^*(x^{\alpha+\beta}) \quad \quad \text{if } |\alpha + \beta|\leq d \\
    h_{\alpha + \beta} \quad \quad \text{if } |\alpha + \beta|> d
    \end{cases}\]

The next result describes the conditions that the parameters $h$'s above need to satisfy. 

\begin{theorem}[{\cite[Theorem 4.2]{Bernard}}]\label{Thm: ExtensionBernard}
    Let $B=\{x^{\beta_1}, \ldots, x^{\beta_r}\}$ be a set of monomials of degree at most $d$ connected to 1 and $\Lambda\in \langle B\cdot B^+\rangle^*_{\leq d}$. Let $\Lambda(h)$ be a form in $\langle B\cdot B^+\rangle^*$ defined by $\Lambda(h)(x^{\alpha})=\Lambda(x^{\alpha})$ if $|\alpha|\leq d$ and $h_{\alpha}$ otherwise. Then $\Lambda(h)$ admits an extension $\Tilde{\Lambda}\in R^*$ such that $H_{\Tilde{\Lambda}}$ is of rank $r$ with $B$ a basis of $A_{\Tilde{\Lambda}}$ if and only if 
    \[M_{x_i}^B(h) \circ M_{x_j}^B(h) - M_{x_j}^B(h) \circ M_{x_i}^B(h)=0\]
    for all $i,j=1,\ldots,n$ and the determinant of $H_{\Lambda}^B(h)$ is non-zero. Moreover, $\Tilde{\Lambda}$ is unique. 
\end{theorem}


\bigskip 

\end{comment}
%\subsection{The cactus algorithm}

%In \cite{Alessandra}, the authors used these results to find an algorithm to compute the cactus rank of a homogeneous polynomial, based on the algorithm in \cite{Bernard}.



\begin{remark}
The initial setup in \cite{Alessandra} for connecting a homogeneous polynomial $F \in S_d$ to a dual form in $S_d^*$ utilizes the apolar product introduced in \cite{Bernard}. Our approach, however, is intended to show that this apolar product can be entirely avoided by adopting the divided power setting, yielding the same dual form in $S_d^*$ in a more coordinate-free manner.

The apolar product of two homogeneous polynomials $F=\sum_{|\alpha|=d} f_{\alpha}x^{\alpha}$ and $G=\sum_{|\alpha|=d} g_{\alpha}x^{\alpha}\in S_d$, is: %\cite{Bernard}:
\[\langle G,F \rangle:=\sum_{|\alpha|=d} \frac{f_\alpha g_\alpha}{\binom{d}{\alpha}}=\frac{1}{d!}\sum_{|\alpha|=d} \alpha! f_\alpha g_\alpha. \]
In our divided power formalism, for $P\in S_d$ and $F\in S_d^*$, we use the apolarity pairing $\langle p, F\rangle_{\aprod}:= p\aprod F$ instead. If we identify the divided powers form of $F$ with a polynomial in $\K_{dp}[Y_0,...,Y_n]$, $F=\sum_{|\alpha|=d}\alpha! f_\alpha Y^{(\alpha)}$, then the pairing with $G=\sum_{|\alpha|=d} g_{\alpha}Y^{(\alpha)}$ is:
\[\langle G,F\rangle_\aprod=\langle G,\sum_{|\alpha|=d} \alpha! f_\alpha Y^{(\alpha)}\rangle_\aprod =\sum_{|\alpha|=d} \alpha! f_\alpha g_\alpha.\]
Evidently, these two expressions agree up to a factor of $d!$.
\end{remark}


%Given $F \in S_d^*$, by  \Cref{remark:extensionsAreApolar} in \Cref{affine:apolarity:criterion}, the kernel of the Hankel operator of an extension $\Lambda$ of $f=(\pi^{-1})^*(F)=F(Y_0=1)$ is apolar to $F$. Thus, by iterating on the rank $r$ of $R/\ker(H_\Lambda)$, one can use all the above results to find minimal apolar schemes for $F$ and their supports. 
%To initiate this process, one needs to determine the first rank to test.

\begin{notation}\label{notation: InitialRank}
Let $H_f^{\ostar}$ be a largest numerical submatrix of $H_{\Lambda(h)}$  with full rank.
\end{notation}
We can present the algorithm for the cactus rank described in \cite{Alessandra}.
  
%whose size is $\binom{n+\lceil d/2 \rceil}{n}\times \binom{n+\lfloor d/2\rfloor}{n} $. 
%If the rank of $H_{\Lambda}$ is $r$, then $r\geq \text{rank }H_f^{\ostar}$, since $H_f^{\ostar}$ is a submatrix of $H_{\Lambda(h)}$, which has rank $r$.


\begin{comment}
    
    
The initial setting of \cite{Alessandra}  slightly differs from our approach in the sense that there is first an assignment from a homogeneous polynomial $F\in S_d$ to a dual form in $S_d^*$, using the \emph{apolar product} introduced in \cite{Bernard}. 

We  show that this apolar product can be completely avoided if one takes the divided power setting, obtaining the same form in $S_d^*$ in a more coordinate-free way:   

\begin{definition}
We define the \emph{apolar product} of $F=\sum_{|\alpha|=d} f_{\alpha}x^{\alpha}$ and $G=\sum_{|\alpha|=d} g_{\alpha}x^{\alpha}\in S_d$ as
\[\langle F,G \rangle:=\sum_{|\alpha|=d} \frac{f_\alpha g_\alpha}{\binom{d}{\alpha}}=\frac{1}{d!}\sum_{|\alpha|=d} \alpha! f_\alpha g_\alpha \]    
\end{definition}

\bigskip

Recall that given $d\geq 1$, we have defined the apolarity pairing $ \langle p, F\rangle_{\aprod}:= p\aprod F$ for $P\in S_d$, $F\in S_d^*$. We first show that this pairing (in the divided powers formalism), is essentially the same as the \emph{apolar product} defined in  \cite{Bernard}.


\bigskip

If we identify the divided powers form of $F$ with a polynomial in $\K_{dp}[Y_0,...,Y_n]$, $F=\sum_{|\alpha|=d}\alpha! f_\alpha Y^{(\alpha)}$ then
we have 
\[\langle G,F\rangle_\aprod=\langle G,\sum_{|\alpha|=d} \alpha! f_\alpha Y^{(\alpha)}\rangle_\aprod =\sum_{|\alpha|=d} \alpha! f_\alpha g_\alpha\]

We can see that the two expressions agree up to multiplication by $d!$.


\bigskip
Now given $F\in S_d^*$, by Remark \Cref{remark:extensionsAreApolar} the kernel of the Hankel of an extension of $f=(\pi^{-1})^*(F)=F(Y_0=1)$ is apolar to $F$. Hence, iterating the rank $r$ of $R/\ker(H_\Lambda)$, where $\Lambda$ is an extension of $f$, we can use the previous results to find minimal apolar schemes to $F$ and its supports.

\bigskip
We still need to know what the first rank to test is. We define $H_f^{\ostar}$ as the largest numerical submatrix of $H_{\Lambda(h)}$, which is checked to be of size $\binom{n+\lceil d/2 \rceil}{n}\times \binom{n+\lfloor d/2\rfloor}{n} $. If the rank of $H_{\Lambda}$ is $r$, then $r\geq \text{rank }H_f^{\ostar}$. Indeed, $H_f^{\ostar}$ is a submatrix of $H_{\Lambda(h)}$, which is of rank $r$, so $H_f^{\ostar}$ cannot have bigger rank than $r$. 


\end{comment}
%All this discussion produced the following algorithm, which is essentially the one presented in \cite{Alessandra}: 

\vspace{0.5cm}
\begin{mdframed}[]
\begin{alg}[Cactus rank and decomposition]\label{alg:cactus}\end{alg}
%\begin{center}
%\underline{\textbf{Algorithm:} Cactus rank and decomposition}
%\end{center}
\vspace{0.1cm}
\noindent\textbf{Input:} A degree $d \geq 2$ polynomial $F \in S_d$ \\
    \textbf{Output:} Cactus rank of $F$ and its support $\zeta_1, \ldots ,\zeta_s\in \K^n$.

\begin{enumerate}
    \item Construct the matrix $H_{\Lambda(h)}$ with parameters $\{h_{\alpha}\}_{\alpha\in \mathbb N^n}$, $|\alpha|>d$.
    \item\label{alg:H:star} Set $r=\text{rank } H_f^{\ostar}$.
    \item\label{alg:basis} Take $B\subseteq R$ a complete staircase of monomials with $|B|=r$, do: 
    \begin{itemize}
        \item\label{alg:h's} Find $h$'s such that:
        \begin{itemize}
            \item $H_{\Lambda(h)}^B$ has nonzero determinant
            \item The multiplication operators $(M_{x_i})^t$ commute for all $i=1,\ldots,n$.
        \end{itemize}
        \item If found, go to \Cref{solutions}. If not, go to \Cref{alg:basis} with another choice of bases $B$. If all choices of $B$ with $|B|=r$ have been already performed, go to \Cref{increase:r}.
    \end{itemize}
    \item\label{increase:r} Set $r\to r+1$ and go to \Cref{alg:basis}.
    \item\label{solutions} The supports $\zeta_j$ of a minimal apolar scheme to $F$ are the common eigenvectors of $M_{x_i}$. The cactus rank of $F$ is the sum of the dimensions of the intersections over $j$ of the generalized egenspaces of $(M_{x_j}^B)^t$ relatives to $(\zeta_i)_j$. %(see \Cref{rmk:computeSupports}).
\end{enumerate}
\end{mdframed}
\vspace{0.5cm}


\subsection{The local cactus algorithm}\label{sec:GAD}


We now explore how the global cactus algorithm in \cite{Alessandra} can be adapted explicitly to the local setting. We start by understanding how to use the basic facts on Hilbert functions of local Artinian algebras described in \cite{Iarrobino}.




\subsubsection{Basis from Hilbert function and local conditions on parameters}

The Hilbert function (HF) of a graded Artinian Gorenstein algebra $R/I$, with $I$ homogeneous, uniquely determines the degrees of monomials in a monomial basis. However, for local Artinian Gorenstein algebras, if $I$ is not homogeneous, the Hilbert function alone does not determine a unique monomial basis. For example, in the algebra $\K[x,y]/(x-y^2,y^3)$, the Hilbert function is $(1,1,1)$, but distinct monomial bases  are possible: $\{1,x,y\}$  can be found via grevlex order, but since $x=y^2 \mod I$,  also $\{1,y,y^2\}$, is a basis for  $\K[x,y]/I$, precisely the one highlighting the structure of the Hilbert function.

%The Hilbert function of a graded local Artinian Gorenstein algebra, $R/I$ with $I$ homogeneous, gives us a lot of information on the structure of a monomial basis of the vector space. Since $R/I$ is graded with $(R/I)_i=R_i/I_i$, $R_i$ is spanned by monomials of degree $i$, and the class of monomials of different degree are linearly independent, any monomial basis of $A$ has $HF(i)$ of monomials of degree $i$. 

%\bigskip

%This is false if $I$ is not homogeneous and $R/I$ is a local algebra. For example, take $\K[x,y]/(x-y^2, y^3)$. Its Hilbert function is $(1,1,1)$. One basis is the representatives of $\{1,x,y\}$ (the one given by the grevlex order), but since $x=y^2$ mod $I$, we can also choose the representatives of $\{1,y,y^2\}$, which is the basis "prescribed" by the Hilbert function.  


%\bigskip

The following results, which we prove explicitly, show that despite this ambiguity, it is always possible to choose a suitable basis compatible with the local structure:

\begin{lemma}\label{base:not:change}
Let $A=R/I$ be a local Artinian ring supported at $\zeta$ with a complete staircase basis $B$. After the affine change of coordinates $\varphi$ of \eqref{affine:tranlsation}, 
%sending $\zeta$ to $0$
 the set $B$ remains a basis of $\varphi(A)$.
\end{lemma}
\begin{proof}
    Let $x^{\alpha}\in B$. Then 
    \[\varphi(x^\alpha)=(x_1-\zeta_1)^{\alpha_1}\cdots (x_n-\zeta_n)^{\alpha_n}=\prod_{i=1}^n \left( \sum_{k_i=0}^{\alpha_i} (-1)^{\alpha_i - k_i}\binom{\alpha_i}{k_i}x_i^{k_i}\zeta_i^{\alpha_i-k_i}\right)=\]
    \begin{equation}\label{summand:expansion}
      =  \sum_{k_1,\ldots,k_n=0}^{\alpha_1,...,\alpha_n} \left( \prod_{i=1}^n (-1)^{\alpha_i -k_i}\binom{\alpha_i}{k_i}\zeta_i^{\alpha_i-k_i}x_i^{k_i} \right).
    \end{equation}
    Remark that the linear combination \eqref{summand:expansion} above is made by $x^{\alpha}$ and lower degree terms, and all other monomials divide $x^{\alpha}$. 

Hence, since $B$ is a complete staircase, we have that \eqref{summand:expansion} is a linear combination of elements in $B$. Thus, the span of $\varphi(B)$ is the span of $B$, and since $\varphi$ is an isomorphism of rings and of vector spaces, we have that $B$ is also a basis of $\varphi(B)$. \end{proof}
This allows to assume the algebra to be supported at 0.




%The next two results state that there is always a choice of basis satisfying this property. 

%\bigskip

%First we prove that we can assume the algebra to be supported at 0:

%\begin{lemma}
%    Let $A=R/I$ be an Artinian local ring supported at $\zeta$ with a basis $B$ that is a complete staircase. Let $\varphi$ be the linear change of variables sending $\zeta$ to 0, that is, $x_i\mapsto x_i-\zeta_i$. Then $B$ is also a basis of $\varphi(A)$. 
%\begin{proof}
%    Let $x^{\alpha}\in B$. Then 
%    \[\varphi(x^\alpha)=(x_1-\zeta_1)^{\alpha_1}\cdots (x_n-\zeta_n)^{\alpha_n}=\prod_{i=1}^n \left( \sum_{k_i=0}^{\alpha_i} (-1)^{\alpha_i - k_i}\binom{\alpha_i}{k_i}x_i^{k_i}\zeta_i^{\alpha_i-k_i}\right)=\]
 %   \begin{equation}\label{summand:expansion}
%      =  \sum_{k_1,\ldots,k_n=0}^{\alpha_1,...,\alpha_n} \left( \prod_{i=1}^n (-1)^{\alpha_i -k_i}\binom{\alpha_i}{k_i}\zeta_i^{\alpha_i-k_i}x_i^{k_i} \right).
%    \end{equation}
%    Remark that the linear combination \eqref{summand:expansion} above is made by $x^{\alpha}$ and lower degree terms and all other monomials divide $x^{\alpha}$. 
%
%Hence, since $B$ is a complete staircase, we have that \eqref{summand:expansion} is a linear combination of elements in $B$. Thus, the span of $\varphi(B)$ is the span of $B$, and since $\varphi$ is an isomorphism of rings and of vector spaces, we have that $B$ is also a basis of $\varphi(B)$. 
%    \end{proof}

%\end{lemma}


\paragraph{\textbf{Restriction on the number of bases for the local case}}
\begin{lemma}\label{lemma: BasisFromHF}
    Let $A= \K[x_1 , \ldots , x_n]/I$ be an Artinian local ring of codimension $n$. Then there exists a complete staircase of monomials $B$ whose reduction modulo $I$ is a $\K$-basis, and the number of elements in $B$ of degree $i$ is $\operatorname{HF}_A(i)$.
\end{lemma}

\begin{proof}
    By \Cref{base:not:change} and the invariance of the Hilbert function of a local Artinian algebra under affine change of coordinates, we may assume that $A$ is supported at the origin.  
    We will prove the existence of a basis $B$ of $A$ consisting of monomials forming a complete staircase, with the degree distribution prescribed by the Hilbert function. The proof is in three steps:

    \begin{itemize}
        \item Existence of a monomial basis respecting the degree distribution.  
        Let $\mathfrak{m}$ be the maximal ideal of $A$. For each $i \ge 0$, the quotient $\mathfrak{m}^i / \mathfrak{m}^{i+1}$ is spanned by the residue classes of monomials of total degree $i$.  
        Choosing a $\K$-basis of $\mathfrak{m}^i / \mathfrak{m}^{i+1}$ made of monomials for each $i$, and collecting them for all $i$, we obtain a monomial basis $B$ of $A$. The number of basis elements of degree $i$ is then $\dim_{\K} \mathfrak{m}^i / \mathfrak{m}^{i+1} = \operatorname{HF}_A(i)$.

        \item Staircase property.  
        Suppose $b \in B$ has degree $i$ and $x_j \in \{x_1,\ldots,x_n\}$. If $b x_j$ is nonzero in $\mathfrak{m}^{i+1}/\mathfrak{m}^{i+2}$, then $b \notin \mathfrak{m}^{i+1}$ (otherwise $b x_j \in \mathfrak{m}^{i+2}$, contradicting $b x_j \neq 0$).  
        This shows that whenever $b \in B$ and $b x_j \neq 0$ in $A$, the product $b x_j$ lies in the basis of degree $i+1$. Hence $B$ can be chosen to be closed under multiplication by variables, i.e. a complete staircase.

        \item Inclusion of $1$ and the variables.
        Since $\operatorname{HF}_A(0) = 1$ and $\operatorname{HF}_A(1) = n$, the constant $1$ and the classes of $x_1,\ldots,x_n$ must appear in $B$.
    \end{itemize}
\end{proof}



\begin{comment}
\begin{proof} 
    By \Cref{base:not:change} and the fact that the Hilbert function of a local Artinian algebra is invariant under linear change of coordinates, we can assume that $A$ is supported at $0$. In order to show that there exists a basis of $A$ that is a complete staircase satisfying the property of the statement, we need to show that there exists a monomial basis of $A$ with that property, that this basis can be chosen so that it is a staircase, and that 1 and all the variables $x_1,\ldots,x_n$ are in the basis. We do these three parts separately:
    
    \begin{enumerate}[label=\roman*)]
        
        \item If $\m$ is the maximal ideal of $A$,  we have that $\m^{i}\subseteq A$ is linearly spanned by the residue classes of all monomials of degree greater or equal than $i$. Hence, there is a monomial basis of $\m^{i}/\m^{i+1}$ of elements of degree $i$. Moreover, elements in $\m^{i}/\m^{i+1}$ are linearly independent to elements in $\m^j$, $j>i$. This allows to choose monomial basis $B$ of $A$ such that, if $b$ has degree $i$, then it is an element of $\m^{i}/\m^{i+1}$.
        \item Let $x_j \in A$ be a variable and $b \in B$, then if $b  x_j$ is a non zero element of $\m^{i+1}/\m^{i+2}$  for some $j=1,\ldots, n$, then $b$ is non-zero in $\m^{i}/\m^{i+1}$ (simply because if $b \in \m^{i+1}$ then $bx_j \in \m^{i+2}$ contradicting $bx_j \neq 0$ in $\m^{i+1}/\m^{i+2}$).
        \item Since $H_A(0)=1$ and $H_A(1)=n$, we can take $1,x_1,\ldots , x_n$ as elements of the basis $B$.
    \end{enumerate}     
    %, as otherwise there is a polynomial $p$ of degree at least $i+1$ such that $b-p\in I$, therefore $x_jb-x_jp\in I$ and $x_jb\in \m^{i+2}$, contradiction. 
\end{proof}

\end{comment}

In our algorithmic setting, \Cref{lemma: BasisFromHF} is key to selecting admissible bases when searching for local apolar schemes. This refines \Cref{alg:basis} of the global algorithm in \cite{Alessandra} by  shrinking the search space: instead of testing all possible bases of length $r = \dim_\K A$, we can restrict to staircase bases consistent with the possible Hilbert functions.

We compare the growth of the search space in two settings:
\begin{itemize}
    \item 
\emph{Unconstrained search.}
Let $r = \dim_\K A$ and suppose we consider all monomial subsets of length $r$ from the set of all monomials of degree $\le d$. The number of such monomials is 
$N_d = \binom{n+d}{d}
$.

For fixed $d$ ($n$ resp.), $N_d$ grows polynomially in $n$ ($d$ resp.); if both grow, the growth is superpolynomial.  
The number of possible bases is therefore $
\binom{N_d}{r} = \binom{\binom{n+d}{d}}{r}$,
which becomes rapidly intractable as $n$ and $d$ increase.

\item \emph{Search constrained by the Hilbert function.}
If $\operatorname{HF}_A = (h_0,\ldots,h_d)$, then any staircase basis $B$ has exactly $h_k$ elements of degree $k$.  
The number of sets of monomials that merely satisfy the degree distribution is bounded above by $
\prod_{k=0}^d \binom{\binom{n+k-1}{k}}{h_k}$, 
where $\binom{n+k-1}{k}$ is the number of monomials of exact degree $k$.  
However, \Cref{lemma: BasisFromHF} ensures that $B$ must also be a complete staircase, which reduces again the number of admissible choices: for a given $\operatorname{HF}_A$, often only very few staircase bases exist.  

\end{itemize}
Thus, instead of choosing $r$ monomials arbitrarily from a large and rapidly growing set, we are choosing a highly structured set dictated by the Hilbert function, yielding a far smaller and more tractable search space.

\begin{example}
\begin{comment}
    Consider $n=2$ variables with Hilbert function 
$
\operatorname{HF}_A = (1,2,2,1)$,

corresponding, for instance, to $A \cong \K[x,y]/(x^3,y^3)$, which has $\dim_\K A = r = 6$.  
Here $d = 3$ is the socle degree.


\emph{Unconstrained search.}
%All monomials in $x,y$ of degree $\le 3$ form the set
%\[
%\{ 1, x, y, x^2, xy, y^2, x^3, x^2y, xy^2, y^3 \}.
%\]
There are $N_3 = \binom{2+3}{3} = 10$ monomials of degree $ 3$ in 2 variables.  
An unconstrained search would choose any $r=6$ monomials from these $10$, giving
$\binom{10}{6} = 210$ possible candidate bases.

\emph{Search constrained by $\operatorname{HF}_A$.}
The Hilbert function prescribes the degree distribution, i.e.:  $h_0=1$ element of degree 0, $h_1=2$ elements of degree 1,  
 $h_2=2$ elements of degree 2 and  
 $h_3=1$ element of degree 3.  
%The number of monomials of exact degree $k$ is $\binom{2+k-1}{k} = k+1$.  
An upper bound for the number of bases respecting this distribution is
$
\binom{1}{1} \cdot \binom{2}{2} \cdot \binom{3}{2} \cdot \binom{4}{1} 
= 1 \cdot 1 \cdot 3 \cdot 4 = 12$.
So at most $12$ sets of monomials satisfy the degree distribution.

\emph{Search restricted to staircase bases.}
\Cref{lemma: BasisFromHF} guarantees that the basis can be chosen as a \emph{complete staircase}, i.e.\ closed under divisibility.  
For $\operatorname{HF}_A=(1,2,2,1)$ there are only two such staircases:$
B_1 = \{1, x, y, x^2, xy, x^2y \}, \quad 
B_2 = \{1, x, y, x^2, xy, xy^2 \}$.
Thus, from $210$ possible bases in the unconstrained setting, the search space collapses to just $2$ admissible candidates.

\end{comment}

Consider $n=2$ variables with Hilbert function 
$\operatorname{HF}_A = (1,2,2,1,1)$,

corresponding, for instance, to $A \cong \K[x,y]/\operatorname{Ann}(x^{(4)}+y^{(3)})$, which has $\dim_\K A = r = 7$.  
Here $d = 4$ is the socle degree.

\emph{Unconstrained search.}
%All monomials in $x,y$ of degree $\le 3$ form the set
%\[
%\{ 1, x, y, x^2, xy, y^2, x^3, x^2y, xy^2, y^3 \}.
%\]
There are $N_4 = \binom{2+4}{4} = 15$ monomials of degree at most 4 in 2 variables.  
An unconstrained search would choose any $r=7$ monomials from these $15$, giving
$\binom{15}{7} = 6435$ possible candidate bases.

\sloppy\emph{Search constrained by $\operatorname{HF}_A$.}
The Hilbert function prescribes the degree distribution, i.e.:  $h_0=1$ element of degree 0, $h_1=2$ elements of degree 1,  
 $h_2=2$ elements of degree 2, 
 $h_3=1$ element of degree 3 and $h_4=1$ element of degree 4  
%The number of monomials of exact degree $k$ is $\binom{2+k-1}{k} = k+1$.  
An upper bound for the number of bases respecting this distribution is
$
\binom{1}{1} \cdot \binom{2}{2} \cdot \binom{3}{2} \cdot \binom{4}{1} \cdot \binom{5}{1} 
= 1 \cdot 1 \cdot 3 \cdot 4\cdot 5 = 60$.
So at most $60$ sets of monomials satisfy the degree distribution.

\emph{Search restricted to staircase bases.}
\Cref{lemma: BasisFromHF} guarantees that the basis can be chosen as a \emph{complete staircase}, i.e.\ closed under divisibility.  
For $\operatorname{HF}_A=(1,2,2,1,1)$ there are only four such staircases: $
B_1 = \{1, x, y, x^2, xy, x^3, x^4\}$, $ 
B_2 = \{1, x, y, x^2, y^2, x^3, x^4\}$, $ 
B_3 = \{1, x, y, y^2, xy, y^3, y^4\}$ and $ 
B_4 = \{1, x, y, y^2, xy, y^3, y^4 \}$.
Thus, from $6435$ possible bases in the unconstrained setting, the search space collapses to just $4$ admissible candidates.


\end{example}
\begin{remark}
For $n=2$, staircase bases of an Artinian monomial algebra correspond to \emph{Young diagrams} determined by the Hilbert function: the monomials in the basis can be represented as lattice points in $\mathbb{N}^2$ below a monotone path.  

For $n=3$, staircase bases correspond to \emph{3-dimensional Ferrers diagrams}, 
also known as \emph{plane partitions}: the exponent set of the basis can be visualized as a stack of unit cubes in $\mathbb{N}^3$ with nonincreasing heights along each axis.  
A fixed Hilbert function imposes strong constraints on the shape of such a diagram. 

While for $n=2$ the number of staircases can be enumerated combinatorically via lattice paths, for $n=3$ this becomes a much harder problem: it amounts to counting plane partitions with a prescribed profile, for which no general closed formula is known. 
Only special families of Hilbert functions can be handled explicitly.
\end{remark}
\paragraph{\textbf{Restriction on the parameters for the local case}}

\begin{remark}\label{remark:local}    
By exploiting the local condition, we can make another  restriction. The operators $M_{x_i}$ (multiplication by variables) have a single eigenvalue at a local point  $\zeta=(1, \frac{1}{r}tr(M_{x_1}),\ldots,\frac{1}{r}tr(M_{x_n}))$, so one can enforce the algebraic constraints on the characteristic polynomials of these operators:
%This accounts for \Cref{alg:basis} in \Cref{alg:cactus}, the search of possible basis. As for the conditions on the parameters $h$ (\cref{alg:basis} in \Cref{alg:cactus}), by \cite[Theorem 4.23]{BernardBook}, the eigenvalues of $M_{x_i}$ are $\{\zeta_{1,i},\ldots,\zeta_{m,i}\}$, where $\{\zeta_1,\ldots,\zeta_m\}$ is the variety of points that define $\ker(H_\Lambda)$. Hence, we want the operators $M_{x_i}$ to have a single eigenvalue. As before, we will have the support \[\zeta=(1, \frac{1}{r}tr(M_{x_1}),\ldots,\frac{1}{r}tr(M_{x_n}))\]
%which now depends on the parameters $h's$.
%\bigskip
%We can impose that $M_{x_i}$ has a single eigenvalue by the closed condition 
$$\operatorname{Char}_{M_{x_i}}(t)=\left(t-\frac{1}{r}tr(M_{x_i})\right)^r.$$
\end{remark}
These two restrictions give effective algebraic constraints to the algorithm for the local setting and we can refine it.

%Summarizing, the refinement of the local cactus rank algorithm explicitly includes:
%\begin{itemize}
%\item The explicit use of Hilbert functions to select admissible complete staircase monomial bases.
%\item Additional algebraic constraints enforcing locality conditions on the multiplication operators.
%\end{itemize}
\begin{comment}
    
\vspace{0.5cm}
\begin{mdframed}[]
\begin{alg}[Local Cactus rank]\label{alg:LocalCactus}\end{alg}
%\begin{center}
%\underline{\textbf{Algorithm:} Cactus rank and decomposition}
%\end{center}
\vspace{0.1cm}
\noindent\textbf{Input:} Degree $d \geq 2$ polynomial $F \in S_d$ \\
    \textbf{Output:} Local cactus rank of $F$, cactus scheme of $F$ and its support $\zeta \in \K^n$.

\begin{enumerate}
    \item Construct the matrix $H_{\Lambda(h)}$ with parameters $\{h_{\alpha}\}_{\alpha\in \mathbb N^n}$, $|\alpha|>d$.
    \item\label{alg:local:CHStar} Set $r=\text{rank } H_f^{\ostar}$.

    \item\label{Step3} For each admissible Hilbert function $ h = (1,h_1,\dots,h_{d'}) $ of length $ r $:
    
    \item\label{alg:local:basis} Select $B\subseteq R$ a complete staircase of monomials with $|B|=r$ with $h_i$ monomials of degree $i$ and do: 
    \begin{itemize}
        \item\label{alg:localC:h's} Find $h$'s such that:
        \begin{itemize}
            \item $H_{\Lambda(h)}^B$ has nonzero determinant
            \item The multiplication operators $(M_{x_i})^t$ commute for all $i=1,\ldots,n$.
            \item The multiplication operators $M_{x_i}$ have a unique eigenvalue: 
            \[
            \text{char}_{M_{x_i}}(t) = (t - \tfrac{1}{r} \mathrm{tr}(M_{x_i}))^r \]
            
        \end{itemize}
    
        \item If found, go to \Cref{alg:localC:increase:r}. If not, go to \Cref{alg:basis} with another choice of $B$. 
    \end{itemize}
    \item\label{alg:localC:increase:r} Set $r\to r+1$ and go to \Cref{Step3}.
    \item Compute the matrices $M_{x_i}$ and the support $\zeta$
    .
\end{enumerate}
\end{mdframed}
\vspace{0.5cm}
\end{comment}

%This explicit adaptation is illustrated concretely in the following example.

\begin{example}\label{Example:LocalCactus}
Consider the following polynomial:    
{\small{
\[
\begin{aligned}
F =\;& 6x_0^{(3)}x_1x_2 + 12x_0^{(2)}x_1^{(2)}x_2 + 18x_0x_1^{(3)}x_2 + 24x_1^{(4)}x_2 \\
&+ 12x_0^{(3)}x_2^{(2)} + 12x_0^{(2)}x_1x_2^{(2)} + 12x_0x_1^{(2)}x_2^{(2)} + 12x_1^{(3)}x_2^{(2)} \\
&+ 6x_0^{(2)}x_1x_2x_3 + 12x_0x_1^{(2)}x_2x_3 + 18x_1^{(3)}x_2x_3 \\
&+ 12x_0^{(2)}x_2^{(2)}x_3 + 12x_0x_1x_2^{(2)}x_3 + 12x_1^{(2)}x_2^{(2)}x_3 \\
&+ 6x_0x_1x_2x_3^{(2)} + 12x_1^{(2)}x_2x_3^{(2)} + 12x_0x_2^{(2)}x_3^{(2)} + 12x_1x_2^{(2)}x_3^{(2)} \\
&+ 12x_0^{(2)}x_3^{(3)} + 12x_0x_1x_3^{(3)} + 12x_1^{(2)}x_3^{(3)} + 6x_1x_2x_3^{(3)} + 12x_2^{(2)}x_3^{(3)} \\
&+ 48x_0x_3^{(4)} + 48x_1x_3^{(4)} + 120x_3^{(5)}
\end{aligned}
\]
}}

After building its Hankel matrix, start the iteration on the rank. The lowest possible rank is 4, since it is the number of essential variables of the polynomial. The only possible complete staircase basis is $\{1,x_1,x_2,x_3\}$, which has the following minor in the Hankel:
\[
\begin{bmatrix}
0 & 0 & 0 & 0 \\
0 & 0 & 1 & 0 \\
0 & 1 & 2 & 0 \\
0 & 0 & 0 & 0 \\
\end{bmatrix}.
\]
Since it has zero determinant, the local cactus rank cannot be 4. One checks that all the 6 possible complete staircase monomial basis of 5 elements also have minors in the Hankel with zero determinant. Thus, the rank is not 5. For $r=6$, there are 18 possible basis. We can use the Hilbert function criterion in the following way. For codimension 3, the first possible Hilbert function of a local Artinian Gorenstein algebra of length 6 is $(1,3,1,1)$. So for this $h$-vector there are 3 candidate basis $B$:
\[\{1,x_1,x_2,x_3,x_1^2,x_1^3\}, \quad \{1,x_1,x_2,x_3,x_2^2,x_2^3\}, \quad \{1,x_1,x_2,x_3,x_3^2,x_3^3\}.\]
The first two have vanishing minor in the Hankel, but the last has the associated matrix (up to scaling)
\[H_{\Lambda}^B=
\begin{bmatrix}
0 & 0 & 0 & 0 & 0 & 2 \\
0 & 0 & 1 & 0 & 0 & 2 \\
0 & 1 & 2 & 0 & 0 & 0 \\
0 & 0 & 0 & 0 & 2 & 8 \\
0 & 0 & 0 & 2 & 8 & 20 \\
2 & 2 & 0 & 8 & 20 & h_{(3,55)} \\
\end{bmatrix}
\]
which has determinant non-zero for every value of the parameter. We now build the multiplication operators with respect to $B$ and impose commutation and locality. We have a zero dimensional polynomial system with a unique solution. Thus, the rank is 6.  
\end{example}

These considerations set the stage for subsequent explicit constructions and algorithmic procedures, such as recovering local Gorenstein algebra decompositions (GAD) and Hilbert functions from commuting multiplication operators, as detailed in the following subsections.

%This is as much as we can say for the search of general Artinian Gorenstein. If we happen to know some information a priori on the local Hilbert function of the scheme, we can obtain more equations on the parameters by computing the Hilbert function that corresponds to the matrices $M_{x_i}$. 

%%%%%%%%%%%%%%%%%%%%%%%%%
%\subsection{Decomposition of the Hilbert function from multiplication operators}

\subsubsection{Hilbert function via multiplication operators}
In this section, we compute the Hilbert function and its symmetric decomposition from an explicit presentation of the apolar algebra $A=R/I$ via multiplication matrices in a complete staircase basis. This method applies when the local algebra is already determined (for instance, as the output of an extension algorithm for $F$), and leverages the full linear structure encoded in the multiplication operators $M_{x_i}$ to recover both the function and its decomposition.

As already observed in \Cref{remark:local}, when $A=R/I$ is a local Artinian algebra the support of the algebra $A$ is explicitly determined from the multiplication operators:
%Assume we have a local Artinian algebra $A=R/I$ such that $I$ is $\m_\zeta$-primary. This can be given in different ways. 
%In our setting, it is given by the set of $n$ pairwise commuting matrices  $M_{x_i}$, $i=1,\ldots n$ such that $M_{x_i}$ is the linear operator of multiplication by $x_i$ in $A$, for some monomial basis of $A$ that is a complete staircase. 
%Here we show how to recover the Hilbert function of the algebra and its symmetric decomposition. 
%\bigskip
%As explained before, we can know which is the support $\zeta=(\zeta_1,\ldots,\zeta_n)\in \K^n$ of the algebra using Theorem \Cref{EigenValues}, for each matrix $M_{x_i}$ has a unique eigenvalue $\zeta_i$ which is equal to the trace divided by the size of the matrix: 
$\zeta=\left(\frac{1}{r}tr(M_{x_1}),\ldots,\frac{1}{r}tr(M_{x_n})\right).$
%\bigskip
The multiplication matrices contain all the multiplicative structure of the algebra. For example, one can also  obtain the coordinates in the basis $B$ of the residue class of any $p\in R$.
\begin{definition}
Given an Artinian algebra $A=R/I$ and a basis $B$ of $A$, the \emph{normal form} of a polynomial $p\in R$ with respect to $I$ and $B$ is the representation of the residue class of $p$ in $A$ as a linear combination of elements in $B$.
\end{definition}

%\begin{definition}
%    Let $A=R/I$ be an Artinian $\K$-algebra and $B$ a $\K$-basis of $A$. For any $p\in R$, the class of $p$ in $A$ expressed as a linear combination of the elements in $B$ is called the \emph{normal form } of $p$ with respect to $I$ and $B$.  
%\end{definition}

\begin{lemma}[{\cite[Proposition 3.2]{NormalForm}}]\label{lemma: NormalForm}
    Let $A=R/I$ be an Artinian $\K$-algebra and $B$ a $\K$-basis of $A$. Let $M_{x_i}$ be the matrix corresponding to the multiplication by $x_i$ in $A$ in the basis $B$. Then for every $p\in R$, the normal for of $p$ with respect to $I$ and $B$ is given by 
    $$p(M_x)(1)\in \langle B\rangle_\K,$$ where $p(M_x)$ is the linear operator in $A$ obtained by substituting each variable $x_i$ in $p$ by $M_{x_i}$.
    \end{lemma}

Note that, since by assumption $1\in B$, the coordinates of $p(M_x)(1)$ is just the first column of $p(M_x)$.


This characterization immediately provides a method to generate powers of the maximal ideal $\m_\zeta$. Indeed, $\m_\zeta^i$ is generated as an ideal by the normal forms of $(x-\zeta)^\alpha$ with $|\alpha|=i$, computed explicitly as $(M_x-\zeta I)^\alpha(1)$. We abbreviate this as $x_{\zeta}^\alpha$. Consequently, the socle degree of $A$ is the smallest integer $d$ such that $x_{\zeta}^\alpha=0$ for all $|\alpha|=d$.
%This result gives us an effective way to compute the generators of powers of the maximal ideal: we know that $\m^{i}$ is generated as an ideal by the class of $\{(x-\zeta)^{\alpha}\}_{|\alpha|= i}$, and therefore by the previous lemma by $(M_x-\zeta Id)^{\alpha}(1)$, $|\alpha|=i$. We will abbreviate this notation with $x_{\zeta}^{\alpha}$. Hence, the socle degree of $A$ is the minimum positive integer $d$ such that $x_{\zeta}^\alpha=0$ for all $|\alpha|=d$.
%\bigskip
Therefore, as a vector subspace, $\m^{i}$ is generated by the class of $\{x_{\zeta}^{\alpha}\}_{i\leq |\alpha|\leq d}$. Thus, we have

\[H_A(i)=\dim_{\K} \frac{\m^{i}}{\m^{i+1}}=\dim_{\K} \langle x_{\zeta}^{\alpha}\rangle_{i\leq |\alpha|\leq d} - \dim_{\K} \langle x_{\zeta}^{\alpha}\rangle_{i+1\leq |\alpha|\leq d}\]
From these observations, we obtain an effective algorithm to compute the Hilbert function $H_A$:
\vspace{0.5cm}
\begin{mdframed}[]
\begin{alg}[Hilbert Function from commuting matrices]\label{alg:HF}\end{alg}
%\begin{center}
%\underline{\textbf{Algorithm:} Cactus rank and decomposition}
%\end{center}
\vspace{0.1cm}
\noindent\textbf{Input:} A list $M$ of $n$ pairwise commuting matrices  of size $r$ with a unique eigenvalue. \\
    \textbf{Output:} The Hilbert function of the Artinian and local algebra $A$, with $M_i$ the matrix corresponding to the multiplication by $x_i$ in $A$.

 \begin{enumerate}
     \item Compute the support of the algebra: $$\zeta=\frac{1}{r}(tr(M_1),\ldots,tr(M_n)).$$
     \item Define $M'_i = M_i - \zeta_i\cdot I$, $i=1,\ldots, n$.
     \item Compute the generators of the powers of the maximal ideal: set $k=0$ and $\text{GensMax}_0= \{1\}$. While $\text{GensMax}_k \neq 0$ do:

     % \[\text{GensMax}_1 = \{M'_i(1) \; | \; i= 1, \ldots , n\}.\]  and set $k:=1$.
    
     
     \begin{itemize}
         \item Compute the generators of the $k+1$-th power of the maximal ideal $$\text{GensMax}_{k+1}= \{M^L \text{for $L$ in compositions$(k+1,n)$}\},$$ where \text{compositions$(k,n)$} is the set of partitions of the integer $k$ with $n$ summands. 
         \item Set $k \; \rightarrow k+1$
     \end{itemize}

     \item Define $d:=k$, the socle degree of the algebra.
     \item For $i$ from 1 to $d$ do:
     \begin{itemize}
         \item Define $\text{Max}_i$ as the matrix of vectors in $\text{GensMax}_j$ for $j$ from $i$ to $d$. 
     \end{itemize}
     \item Print Hilbert function: $H_A(i)=\text{rank}(\text{Max}_i)-\text{rank}(\text{Max}_{i+1})$ for $i$ from $1$ to $d$.
     \end{enumerate}
\end{mdframed}
\vspace{0.5cm}


\begin{example}\label{Example: HF}
    Consider the algebra
    $$A=\K[x,y,z]/\operatorname{Ann}\left(2z^{(3)}+xy+2y^{(2)}\right).$$
    It is Artinian and of length 6, with $\{1,x,y,z,z^2,z^3\}$ as a possible basis. In this basis, the matrices corresponding to the multiplication by the variables are:
{\scriptsize

\[
\textstyle
M_x =
\begin{bmatrix}
0 & 0 & 0 & 0 & 0 & 0 \\
1 & 0 & 0 & 0 & 0 & 0 \\
0 & 0 & 0 & 0 & 0 & 0 \\
0 & 0 & 0 & 0 & 0 & 0 \\
0 & 0 & 0 & 0 & 0 & 0 \\
0 & 0 & \tfrac{1}{2} & 0 & 0 & 0 \\
\end{bmatrix}
\quad
M_y =
\begin{bmatrix}
0 & 0 & 0 & 0 & 0 & 0 \\
0 & 0 & 0 & 0 & 0 & 0 \\
1 & 0 & 0 & 0 & 0 & 0 \\
0 & 0 & 0 & 0 & 0 & 0 \\
0 & 0 & 0 & 0 & 0 & 0 \\
0 & \tfrac{1}{2} & 1 & 0 & 0 & 0 \\
\end{bmatrix}
\quad
M_z =
\begin{bmatrix}
0 & 0 & 0 & 0 & 0 & 0 \\
0 & 0 & 0 & 0 & 0 & 0 \\
0 & 0 & 0 & 0 & 0 & 0 \\
1 & 0 & 0 & 0 & 0 & 0 \\
0 & 0 & 0 & 1 & 0 & 0 \\
0 & 0 & 0 & 0 & 1 & 1 \\
\end{bmatrix}.
\]
}

We clearly see that the three matrices have the unique eigenvalue 0 (as expected since it is the quotient of the annihilator of a polynomial). Then, for instance, the ideal $\m^3$ is generated as a $\K$-vector space by the column vector
{\small{
\[M_z^3(1) =
\begin{bmatrix}
0 \\
0 \\
0 \\
0 \\
0 \\
1 \\
\end{bmatrix}.\]
}}
In our basis, it corresponds to $z^3$.
All the other combinations $M_x^{\alpha_1}\cdot M_y^{\alpha_2}\cdot M_z^{\alpha_3}$ with $\alpha_1+\alpha_2+\alpha_3=3$ give the zero matrix. 
Further computations show that there are no non-zero matrices corresponding to elements in $\m^4$. Hence, the socle degree of $A$ is $3$. Similarly, $\m^2$ is generated as a $\K$-vector space by the column vectors corresponding to $\{z^2,z^3\}$:
{\small{
\[
\begin{bmatrix}
0 & 0\\
0 & 0\\
0 & 0\\
0 & 0\\
0 & 1\\
1 & 0\\
\end{bmatrix}.\]}}


Computing the different powers of the matrices and multiplying them accordingly, we obtain the Hilbert function $(1,3,1,1)$. 
\end{example}

The symmetric decomposition of the Hilbert function is similarly derived by explicitly computing annihilators $(0:\m^{i})$. Specifically, for each $i$, one solves the linear system:


%The symmetric decomposition of the Hilbert function can be obtained through a similar procedure. The ideal $(0:\m^{i})$ is the set of elements $p\in A$ such that $p\cdot x_\zeta^{\alpha}=0$ for all $i\leq |\alpha|\leq d$. Hence, if $B=\{b_1,\ldots,b_r\}$, we have the linear system
\begin{equation}\label{eq: annOfMax}
\sum_{i}\lambda_i b_i x_\zeta^{\alpha}=0, \quad i\leq |\alpha|\leq d,
\end{equation}
where $B={b_j}$ is the chosen monomial basis. This system's dimension equals $\dim_\K(0:\m^{i})$, leading to the full symmetric decomposition.

%The dimension of the linear system above is $\dim_\K (0:\m^{i})$.

\begin{remark}
Each monomial $x_\zeta^{\alpha}$ can be expressed as a linear combination of the basis elements $b_i$, and products are computed via 
$
b_i b_j = (b_i b_j)(M_x)(1)
$.
The basis $B$ itself need not be known explicitly: it can be reconstructed (up to equivalence modulo $I$) from the multiplication matrices.  
Specifically, if for some multi-index $\alpha \in \mathbb{N}^n$ the matrix $M_x^{\alpha}$ has a single nonzero entry in its first column, then $x^{\alpha}$ corresponds to a basis element.  
See \Cref{Example: SymmetricDecompFromMatrices} for an explicit instance.
\end{remark}
\begin{comment}
    \begin{remark}
    Note that $x_\zeta^{\alpha}$ is a linear combination of $b_i$'s, and the product $b_ib_j$ is computed by $b_ib_j(M_x)(1)$. Remark also that we need not know the basis $B$, that is, we can recover it (up to equivalence modulo the ideal $I$) with just the matrices. Essentially, if there is some tuple $\alpha\in \mathbb N^n$ such that $M_x^{\alpha}$ has a unique non-zero entry in the fist column, then $x^\alpha$ is in the basis.  See \Cref{Example: SymmetricDecompFromMatrices} for a concrete example.
\end{remark}

\end{comment}



Using the definitions of the symmetric filtrations, we can recover the full symmetric decomposition. 

\begin{example}\label{Example: SymmetricDecompFromMatrices}
    Let us compute the symmetric decomposition of the Hilbert function obtained in \Cref{Example: HF}.
\bigskip

    Let us first see how to recover the elements of the basis from the matrices $\{M_x,M_y,M_z\}$. Since the size of the matrices is $6$ and $\{1,x,y,z\}$ is always part of a complete monomial staircase basis, we look for two other elements of the basis. We see that 
    {\scriptsize
\[
M_z^2 =
\begin{bmatrix}
0 & 0 & 0 & 0 & 0 & 0 \\
0 & 0 & 0 & 0 & 0 & 0 \\
0 & 0 & 0 & 0 & 0 & 0 \\
0 & 0 & 0 & 0 & 0 & 0 \\
1 & 0 & 0 & 0 & 0 & 0 \\
0 & 0 & 0 & 1 & 0 & 0 \\
\end{bmatrix}
\quad
M_z^3 =
\begin{bmatrix}
0 & 0 & 0 & 0 & 0 & 0 \\
0 & 0 & 0 & 0 & 0 & 0 \\
0 & 0 & 0 & 0 & 0 & 0 \\
0 & 0 & 0 & 0 & 0 & 0 \\
0 & 0 & 0 & 0 & 0 & 0 \\
1 & 0 & 0 & 0 & 0 & 0 \\
\end{bmatrix}.
\]
}
Looking at the first column of the first matrix we see that $z^2$ is in the basis, while the second matrix shows $z^3$ is in the basis. One can check that we also have $M_y^2=M_z^3$, which is consistent with the fact that $y^2$ is equal to $z^3$ modulo $I$. 

We have seen the computation of the powers of the maximal ideal. For the annihilators of the powers of the maximal ideal, applying \eqref{eq: annOfMax} we have that for instance $(0: \m^3)$ is the space of polynomials in $A$ whose coordinates in the basis $\{1,x,y,z,z^2,z^3\}$ are in the kernel of this matrix: 
\[
\begin{bmatrix}
1 & 0 & 0 & 0 & 0 & 0 \\
0 & 0 & 0 & 1 & 0 & 0 
\end{bmatrix}\]
That is, if a polynomial in $A$ is $\lambda_1+\lambda_2x+\ldots + \lambda6z_3$, it must verify $\lambda_1=\lambda_3=0$. Note that this is consistent with the fact that $\m^2$ is generated by $z^2,z^3$.

We can compute the rest of the annihilators and then the intersections with powers of maximal ideals, obtaining the decomposition 

\begin{table}[h]
\centering
\begin{tabular}{|l|l|l|l|l|}
\hline
Degree         & 0 & 1 & 2 & 3  \\ \hline
$\Delta_{Q_0}$ & 1 & 1 & 1 & 1  \\ \hline
$\Delta_{Q_1}$ & 0 & 2 & 0 & 0  \\ \hline
Total          & 1 & 3 & 1 & 1  \\ \hline
\end{tabular} \,.
\end{table}

\end{example}
%%%%%%%%%%%%%%%%%%%%%%%%%%%%%
\subsubsection{Recovery of Local Generalized Additive Decomposition}\label{Section:GAD}

In this subsection we describe how to recover a local Generalized Additive Decomposition (GAD) of  given $F$ (or of one of its extensions) once the support and socle degree of the desired local GAD are fixed.
This viewpoint provides a constructive procedure, starting from the prescribed data.

After a suitable change of coordinates, a homogeneous polynomial $F$ admitting a local GAD supported at a point can be expressed in the form
\begin{equation}\label{GADLocal}
    F = L^{(k)} \cdot G, \quad L \in S_1^*, \; G \in S_{d-k}^*,
\end{equation}
where $L$ is a linear form vanishing at the support point.  
Such an expression naturally arises from dehomogenizing $F$ with respect to $L$; the resulting local apolar schemes were explicitly analyzed in \cite{GAD}, where the scheme is said to be \emph{evinced by the GAD}.


%The construction of local apolar schemes for polynomials of the form 
%\begin{equation}\label{GADLocal}
%    F = L^{(k)} \cdot G, \quad L \in S_1^*, \; G \in S_{d-k}^*,
%\end{equation}
%via dehomogenization with respect to $L$, was explicitly considered in \cite{GAD}.  
%Such schemes is called \emph{evinced by th GAD}.


The relationship between $F$ and its minimal local apolar schemes is clarified in \cite[Proposition~4]{BJPR}: for schemes supported at $[1:0:\dots:0]$, the dehomogenization $f = F(Y_0=1)$ and the form $g$ defining the apolar scheme $\Gamma = \operatorname{Spec}(R/\operatorname{Ann}(g))$  share the same degree-$d$ tail. By an affine change of coordinates, this applies to any support point.

%A key result from \cite[Proposition 4]{BJPR}, proved for schemes supported at $[1:0:\dots:0]$, states that if $F$ is homogeneous and $\Gamma = \operatorname{Spec}(R/\operatorname{Ann}(g))$ is a minimal local apolar scheme to $F$, then the dehomogenization $f = F(Y_0=1)$ and $g$ share the same degree-$d$ tail (by an affine change of coordinates, this description applies to schemes supported at any point).  



Equivalently, a minimal local apolar scheme supported at $[1:0:\dots:0]$ is \emph{evinced} by a form 
\[
G = Y_0^{(d'-k)} \cdot H \in S_{d'}^*, \quad d' \ge d,
\]
such that $G(Y_0=1)$ has the same degree-$d$ tail as $f=F(Y_0=1)$.

\begin{definition}\label{def:extension}
Let $F \in S_d^*$ be a homogeneous polynomial.  
An \emph{extension} of $F$ is a form $G \in S_{d'}^*$ with $d' \ge d$ such that there exists a linear form $L \in S_1^*$ and an integer $m = d' - d \ge 0$ for which 
\[
\partial_L^{\,m} G = F.
\]
When $m=0$ (equivalently $d'=d$), $G$ coincides with $F$; for $m>0$, $G$ is a genuine higher-degree extension of $F$.
\end{definition}
Given $F\in S_d^*$, the local cactus algorithm finds a functional $\Lambda\in R^*$ that extends $F(Y_0=1)\in R_{\leq d}$, so that the algebra $A=R/\ker H_\Lambda$ is local and has minimal length. The minimal apolar scheme is then $Z=\operatorname{Spec}(A)$. The support $\zeta\in \K^n$ is recovered from the multiplication matrices, while the remaining structure is obtained from the dual generator $G$ associated to $\Lambda$.


\begin{remark}
By the Generalized Macaulay Correspondence \Cref{Thm: GeneralisedMacaulayCorrespondence}, a local Gorenstein algebra $A = R/\ker H_\Lambda$ of socle degree $d'$ supported at $\zeta$ corresponds to a functional 
\[
\Lambda = \mathds{1}_\zeta \circ G(\partial),
\]
where $G \in S_{d'}^*$ is a dual generator (possibly an extension of $F$).

As shown in \cite[Proposition~4.3]{GAD}, $Z$ is evinced by a GAD of $F$ with $G=F$ and $d'=d$ if and only if
\[
\mathfrak{p}_L^{d+1} \subseteq I(Z),
\]
where $\mathfrak{p}_L$ is the homogeneous maximal ideal of the support point $[L] \in \mathbb{P}^n$. For a local scheme supported at $[L]$, the socle degree of $R/\ker H_\Lambda$ is the smallest $m$ such that $\mathfrak{p}_L^{m+1} \subseteq I(Z)$. Hence:
\begin{itemize}
    \item If $\operatorname{socdeg}(R/\ker H_\Lambda) \le d$, then $G=F$ (no genuine extension is needed).  
    \item If $\operatorname{socdeg}(R/\ker H_\Lambda) > d$, then $G$ has degree $d'>d$ and is a genuine extension of $F$.
\end{itemize}
\end{remark}


%We will use this terminology when distinguishing between the two cases: socle degree $\le d$, where the extension is $F$ itself, and socle degree $> d$, where a genuine higher-degree extension is needed.


\begin{comment}
---FROM ABOVE, BE FIXED


The construction of local apolar schemes explicitly considered in \cite{GAD} for polynomials of type \begin{equation}\label{GADLocal}
    F=L^{(k)}\cdot G,
\end{equation} $L\in S_1^*$, $G\in S_{d-k}^*$, through dehomogenization with respect to $L$, will be our starting point. Such schemes, called \emph{evinced}, allow an explicit and constructive description of minimal local apolar schemes.


%As it was shown in \cite{BJPR}, \cite{GAD}, there is a natural way to define a local scheme that is apolar to a homogeneous polynomial $F\in S^*_d$ of the form 

%\begin{equation}\label{GADLocal}
%    F=L^{(k)}\cdot G
%\end{equation}
%where $L\in S_1^*$ and $G\in S_{d-k}^*$. For this construction, first we dehomogenize $F$ with respect to $L$, that is, we consider $G_L$ as the image of the projection $S^*\to S^*/(L-1)$, and then define $Z_{F,L}$ as the scheme that is defined locally by the annihilator of $G_L$. The resulting scheme is said to be \emph{evinced} by Eq.\eqref{GADLocal}. 
A fundamental result, originally established for local apolarity at $[1:0:\dots:0]$ in \cite[Proof of Proposition 4]{BJPR}, states that if $F$ is homogeneous and $\Gamma=\text{Spec}(R/\operatorname{Ann}(g))$ is minimal among apolar schemes supported at $[1:0:\dots:0]$, then the degree-$d$ tails of $g$ and the dehomogenization $f=F(Y_0=1)$ coincide. This result naturally extends to general points by affine changes of coordinates.



Another way of stating {\cite[Proof of Proposition 4]{BJPR}} is that a zero-dimensional scheme of minimal length among local schemes supported at
$[1 : 0 : \cdots: 0]$ that are apolar to $F$ is evinced by 
\[G=Y_0^{(d'-k)}\cdot H\in S_{d'}^*\]
where $d'\geq d$, $k\in \{0,\ldots,d'\}$ and $(\pi^{-1})^*(G)=G(Y_0=1)$ has the same degree $d$-tail as $f$. We will return to this construction of the scheme in \Cref{Section:GAD}.

-----
\end{comment}
\begin{comment}
    
Given $F\in S_d^*$, the algorithm described in the previous section for the local cactus rank, finds a form $\Lambda\in R^*$ such that it extends $F(Y_0=1)\in R_{\leq d}$ and the algebra $A=R/\ker H_\Lambda$ is local and has minimal length. Then $Z=\text{Spec}(A)$ is the minimal apolar scheme to $F$. We have seen how to recover the support of $\zeta\in \K^n$ of $A$ using the matrices of the multiplication by the variables. We now show how to recover the remaining information of $Z$.

By \cite[Proposition 4]{BJPR}, there exists a form $G \in S_{d'}^*$ with $d'\geq d$ such that

where $L\in S_1^*$, $H\in S_k^*$, and the scheme $Z$ can be explicitly represented as:
%We recall that by \cite[Proposition 4]{BJPR}, there is a polynomial $G\in S_{d'}^*$ with $d\leq d'$ such that
%\[G=L^{(d'-k)}\cdot H\]
%where $L\in S_1^*$, $k\geq 0$, $H\in S_k^*$ and such that if $G_L$ is the image of $G$ in the projection $S\to S/(L-1)$, then $G_L=H_L$ and $F_L$ agree up to degree $d$ and 
\[Z=\text{Spec}(S/(L-1),\operatorname{Ann}(G_L)),\]
with $G_L$ and $F_L$ agreeing up to degree $d$. After a suitable coordinate change (typically assuming $L=Y_0$), this simplifies to
%For example, if $L=Y_0$ (which we can assume after an appropriate change of coordinates), then 
    $$Z=\text{Spec}(R/\operatorname{Ann}(G(Y_0=1)).$$
Thus, the scheme $Z$ is determined once  such a $G$ is found. 

\begin{remark}
Let $F \in S_d^*$ and let $Z$ be the minimal local scheme apolar to $F$.  
Assume $Z = \operatorname{Spec}(R/\ker H_\Lambda)$ for some $\Lambda \in R^*$.  
By the Generalized Macaulay Correspondence (Corollary~\ref{Thm: GeneralisedMacaulayCorrespondence}), $\Lambda$ can be represented in the form
\[
\Lambda = \mathds{1}_\zeta \circ G(\partial),
\]
where $G \in S_{d'}^*$ is a dual generator (possibly an extension of $F$).  

As shown in \cite[Proposition~4.3]{GAD}, $Z$ is evinced by a GAD of $F$ with $G = F$ and $d' = d$ if and only if
\[
\mathfrak{p}_L^{d+1} \subseteq I(Z),
\]
where $\mathfrak{p}_L$ is the homogeneous maximal ideal defining the support point $[L] \in \mathbb{P}^n$.  

For a local scheme $Z$ supported at $[L]$, the socle degree of $R/\ker H_\Lambda$ coincides with the smallest integer $m$ such that $Z$ is contained in the $(m+1)$-fat point at $[L]$, i.e.
\[
\mathfrak{p}_L^{m+1} \subseteq I(Z).
\]
Thus, the condition $\mathfrak{p}_L^{d+1} \subseteq I(Z)$ is equivalent to requiring that the socle degree of $R/\ker H_\Lambda$ is at most $d$.  

In particular:
\begin{itemize}
    \item  If $\operatorname{socdeg}(R/\ker H_\Lambda) \le d$, then $\Lambda$ corresponds to a dual generator $G$ of degree $d$ (no genuine extension is needed, $G = F$).  
    \item If $\operatorname{socdeg}(R/\ker H_\Lambda) > d$, then $\Lambda$ corresponds to a dual generator $G$ of degree $d' > d$ (a genuine extension of $F$ is required to evince $Z$).
\end{itemize}
\end{remark}

\end{comment}

\begin{comment}

This can be precisley sated.
\begin{lemma}\label{lemma:SocleAndExtension}
Let $F \in S_d^*$ and let $Z$ be the minimal local scheme apolar to $F$.  
Assume $Z = \operatorname{Spec}(R/\ker H_\Lambda)$ for some $\Lambda \in R^*$.  
Then $Z$ is evinced by a GAD of $F$ if and only if the socle degree of $R/\ker H_\Lambda$ satisfies $\operatorname{socdeg}(R/\ker H_\Lambda) \le d$.
\end{lemma}


\begin{proof}
    After a change of coordinates we may assume that
    $R/\ker H_\Lambda$ is supported at $0$, and therefore $\Lambda=g$ for some polynomial $g\in \K_{dp}[Y_1,\ldots, Y_n]$ of degree $k$. Then the socle degree of $R/\ker H_g$ is $k$. Again we know by \cite[Proposition 4]{BJPR} that 
    %\[F=Y_0^(d-k)\cdot H\] for some $k\geq 0$ and 
    that $g$ and
    $F(Y_0=1)$ agree up to degree $d$. Moreover, for every $p\in R$ we denote the homogenization of $p$ of degree $u$ with respect to $x_0$ by  
    \[Homog(p, x_0, u):=x_0^{u}p(\frac{x_1}{x_0}, \ldots, \frac{x_n}{x_0}) \quad \text{if }u\geq \deg p\]
    and $0$ if $u<\deg p$. Then we have 
    \[I(Z)=\left(\{Homog(h,x_0,u) \; | \; h\in Ann(g), \; u\in \mathbb N, \; \deg h\leq k\}\right) + (x_0, \ldots, x_n)^{k+1}\]
    Therefore if $k\leq d$ then $(x_0,\dots,x_n)^{d+1}\subseteq I(Z)$ and $Z$ is evinced by a GAD of $F$. Conversely, assume $(x_0,\dots,x_n)^{d+1}\subseteq I(Z)$ and $k>d$. Then  
    %Therefore $Z$ is evinced by a GAD of $F$ if and only if $(x_0,\ldots, x_n)^{d+1}\subseteq I(Z)$, if and only if 
    for every $p\in (x_0,\ldots, x_n)^{d+1}\backslash (x_0,\ldots, x_n)^{k+1}$, 
    \[p\in \left(\{Homog(h,x_0,u) \; | \; h\in Ann(g), \; u\in \mathbb N, \; \deg h\leq k\}\right)\]
    In particular, we can take $p=x_0^{d+1}$ and we have that $p=\{Homog(h,x_0,u)$ for some $h\in Ann(g)$ and some $u\in \mathbb N$. The only possibility is $h=1$ and $u=k$, which implies $1\in Ann(g)$ and therefore $g=0$, contradiction.
\end{proof}

\bigskip
\end{comment}


The approach to recover a GAD of either $F$ or of an extension of $F$  depends on whether the socle degree does not exceed, or exceeds, the degree of $F$.


\medskip

\paragraph{\textbf{Socle degree $\leq  d$:}}  In this case, a GAD of $F$ exists explicitly as:
\[F=(x_0+\zeta_1x_1+\cdots+\zeta_nx_n)^{(d-k)}H\]
where $H$ is recovered by performing a suitable affine coordinate shift to move $\zeta$ to the origin. After the shift, one can easily factor out the maximal possible power of $x_0$ from $F$.

%Let us first deal with the case where the minimal local apolar scheme $Z$ to $F\in S_d^*$ has socle degree $\leq d$. 
%Indeed $Z$ is defined locally by $\text{Spec}(R/\ker H_\Lambda)$ and it is  supported at $\zeta\in \K^n$. Unless $\zeta=0$, $\Lambda\in R^*$ is an infinite series in the divided power basis.
%Moreover, by the previous result, $Z$ is evinced by a GAD of $F$, i.e., there is a decomposition of $F$ evincing $Z$, 
%\[F=(x_0+\zeta_1x_1+\ldots+\zeta_nx_n)^{(d-k)}H\]
%for some $H\in S_k^*$. Let us show how to recover this expression of $F$ from the support $\zeta$ of $Z$.
%
%\bigskip
%Note that if $F=Y_0^{(d-k)}\cdot H$, then $Z=\text
%{Spec}(R/\operatorname{Ann}(f))$, where $f=F(Y_0=1)$ (but in general $f\neq H(Y_0=1)$, see \Cref{rmk:ProductAndDehom}). Hence, in this case knowing $f$ gives  the decomposition and so $Z$.  Therefore the strategy is to make a change of coordinates on the general case to reduce to this one.
%
%\bigskip
So assume we have the minimal apolar scheme $Z$ to $F\in S_d^*$, and we know the support $\zeta=(\zeta_1,\ldots,\zeta_n)\neq 0$. Therefore $Z$ is defined locally by $\text{Spec}(R/\ker(H_\Lambda))$ for $\Lambda$ an infinite series. Consider the change of coordinates $\varphi$ in $S$ induced by 
\[x_0\mapsto x_0 \quad \quad x_i=x_i-\zeta_ix_0 \quad \forall i=1,\ldots, n\]
The induced affine change of coordinates $\varphi$ on $R$ is given by $x_i\mapsto x_i-\zeta_i$. Now $Z'=\varphi(Z)$ is defined locally by $\text{Spec}\left(R/\ker H_{(\varphi')^*(\Lambda)}\right)$. By \Cref{lemma:DualChangeBasis}, $\Lambda$ is now a polynomial $\Tilde{f}\in \K_{dp}[Y_1,\ldots, Y_n]$, and by hypothesis $f$ has degree $k\leq d$. Therefore we have that $Z'$ is the minimal local apolar scheme to
\[\varphi^*(F)= \operatorname{Homog}(\Tilde{f},x_0, d).\]
Dealing carefully with divided powers we obtain a decomposition 
\[\varphi^*(F)=Y_0^{d-k}\cdot \Tilde{H}.\]
%We have seen in the first section how to compute the dual map $\varphi^*$ in every degree of the dual, and also by \Cref{prop: DualChangeIsAlgebraHom} $\varphi^*$ is a ring homomorphism.
Hence,
\[F=(\varphi^{-1})^*(Y_0^{(d-k)})\cdot (\varphi^{-1})^*(\Tilde{H})=
(x_0+\zeta_1x_1+\ldots+\zeta_nx_n)^{(d-k)}\cdot (\varphi^{-1})^*(\Tilde{H}),\]
which gives the desired decomposition of $F$.



\medskip

In terms of an algorithm, once we have the matrices corresponding to the multiplication by the variables in the algebra $R/\ker H_\Lambda$, for some $\Lambda\in R^*$, we can easily compute the socle degree and the support of the scheme $\zeta$. If the socle degree is not greater than the degree of $F\in S_d^*$, then we can apply the method described above, illustrated in the following example.

\begin{example}
    Consider the following homogeneous polynomial in divided powers: 

\[
\begin{aligned}
 F=&
3x_0^{(3)} x_1 + 8x_0^{(2)} x_1^{(2)} + 12x_0 x_1^{(3)} - 2x_0^{(2)} x_1 x_2 - 4x_0 x_1^{(2)} x_2 + x_0^{(2)} x_2^{(2)} +  \\
&3x_0 x_1 x_2^{(2)} + 4x_1^{(2)} x_2^{(2)} - 3x_0 x_2^{(3)} - 6x_1 x_2^{(3)} + 6x_2^{(4)}.
\end{aligned}
\]
Assume we know that its minimal local apolar scheme is supported at $[1:2:-1]$, and it is evinced by a decomposition of $F$. Following the discussion above, we move the support to the affine 0. By \Cref{prop: ChangeCoordInDivPowers}, first we convert it to a standard polynomial, $\Psi(F)$, and then we make the change of variables $\Tilde{\varphi}$, $x_0 \mapsto x_0-2x_1+x_2$, so that it reduces to a quartic in which $x_0^2$ factors out:
\[
(\Tilde{\varphi} \circ \Psi )(F)=x_0^2 \left(\frac{1}{2}x_0x_1-x_1^2+\frac{1}{2}x_1x_2+\frac{1}{4}x_2^2\right).
\]
Substituting back $x_0 \mapsto x_0+2x_1-x_2$ gives
\[
\Psi(F)=(x_0+2x_1-x_2)^2\left(\frac{1}{2}x_0x_1 + \frac{1}{4}x_2\right).
\]
Converting back to divided powers yields the GAD decomposition of $F$.
\end{example}
\begin{comment}
    
\begin{example}\label{Ex: recov. local GAD}
     Consider the following homogeneous polynomial in divided powers: 

\[
\begin{aligned}
 F=&
3x_0^{(3)} x_1 + 8x_0^{(2)} x_1^{(2)} + 12x_0 x_1^{(3)} - 2x_0^{(2)} x_1 x_2 - 4x_0 x_1^{(2)} x_2 + x_0^{(2)} x_2^{(2)} +  \\
&3x_0 x_1 x_2^{(2)} + 4x_1^{(2)} x_2^{(2)} - 3x_0 x_2^{(3)} - 6x_1 x_2^{(3)} + 6x_2^{(4)}.
\end{aligned}
\]
Assume we know that its minimal local apolar scheme is supported at $[1:2:-1]$, and it is evinced by a decomposition of $F$. 

Following the previous algorithm, we construct the matrix $\varphi_4^*$, which is a $15\times 15$ matrix.
We compute the image of $F$ by this matrix and we obtain 
%\[\varphi^*(F)=3x_0^{(3)} x_1 - 4x_0^{(2)} x_1^{(2)} + x_0^{(2)} x_1 x_2 + x_0^{(2)} x_2^{(2)},\]
%which decomposes
\[\varphi^*(F)=x_0^{(2)}\cdot \left(x_0x_1-4x_1^{(2)}+x_1x_2+x_2^{(2)}\right).\]
Compute the inverse change of coordinates in degree 2:
\[
\begin{bmatrix}
1 & 0 & 0 & 0 & 0 & 0 \\
2 & 1 & 0 & 0 & 0 & 0 \\
-1 & 0 & 1 & 0 & 0 & 0 \\
4 & 4 & 0 & 1 & 0 & 0 \\
-2 & -1 & 2 & 0 & 1 & 0 \\
1 & 0 & -2 & 0 & 0 & 1
\end{bmatrix}.
\]
The image of $x_0x_1-4x_1^{(2)}+x_1x_2+x_2^{(2)}$ by this matrix is $x_2^{(2)} + x_0x_1$, hence 
\[F=(x_0+ 2x_1-x_2)^{(2)}\cdot \left(x_2^{(2)} + x_0x_1\right).\]
\end{example}

There is a simpler way to obtain the same result that does not involve computing these matrices. It is based on the following result.
\end{comment}

\begin{comment}
    
\begin{proposition}\label{Lemma:NonDivPowersChangeCoord}
    Let $\zeta=(\zeta_1,\ldots, \zeta_n)\in \mathbb K^n$ be an affine point, $\varphi:S\to S$ the homogeneous change of coordinates induced by $x_0\mapsto x_0$, $x_i\mapsto x_i+\zeta_i x_0$.
    Consider $\varphi^*:S^*\to S^*$ the dual map, and let $\Psi$ be the isomorphism of rings $$\Psi:\K_{dp}[Y_0,\ldots,Y_0]\to S \quad Y^{(\alpha)}\mapsto \frac{1}{\alpha!}x^{\alpha}$$ converting divided powers expressions into polynomials. Then the algebra homomorphism $\Tilde{\varphi}:S\to S$ induced by $x_0\mapsto x_0+\zeta_1x_1+\ldots + \zeta_n x_n$, $x_i\mapsto x_i$ makes the following diagram commute:

\[\begin{tikzcd}
	{K_{dp}[Y_0,\ldots, Y_n]} & {K_{dp}[Y_0,\ldots, Y_n]} \\
	S & S
	\arrow["{\varphi^*}"{description}, from=1-1, to=1-2]
	\arrow["\Psi", from=1-1, to=2-1]
	\arrow["\Psi", from=1-2, to=2-2]
	\arrow["{\Tilde{\varphi}}"{description}, from=2-1, to=2-2]
\end{tikzcd}\]    
\end{proposition}

\begin{proof}
    Using the identity
    \[Y^{(\beta)}=\frac{1}{\beta!}Y^{\beta}=Y_0^{\beta_0}\cdots Y_n^{\beta_n}\]
    and the fact that, by \eqref{prop: DualChangeIsAlgebraHom}, the map $\varphi^*$ is multiplicative, we have that for every $\beta\in \mathbb N^{n+1}$, 
    \[\Psi(\varphi^*(Y^{(\beta)}))=\Psi(\varphi^*(\frac{1}{\beta!}Y^{\beta}))=\frac{1}{\beta!}\Psi(\prod_{i=0}^n(\varphi^*(Y_i))^{\beta_i}) = \]\[\frac{1}{\beta!}\Psi((Y_0 + \zeta_1 Y_1 +\ldots \zeta_n Y_n)^{\beta_0}\cdot  \prod_{i=1}^n(Y_i)^{\beta_i})=\frac{1}{\beta!}(Y_0 + \zeta_1 Y_1 +\ldots \zeta_n Y_n)^{\beta_0}\prod_{i=1}^n(Y_i)^{\beta_i}\]
which clearly agrees with $(\Tilde{\varphi}\circ \Psi)(Y^{(\beta)})$.
\end{proof}


We can then redo the previous \Cref{Ex: recov. local GAD}:
\begin{example}
Starting from the divided power polynomial in \Cref{Ex: recov. local GAD}, we apply $\Psi$ to obtain its standard polynomial form.  
After the change of variables $x_0 \mapsto x_0-2x_1+x_2$, $\Psi(F)$ reduces to a cubic in which $x_0^2$ factors out:
\[
(\Tilde{\varphi} \circ \Psi )(F)=x_0^2 \left(\frac{1}{2}x_0x_1-x_1^2+\frac{1}{2}x_1x_2+\frac{1}{4}x_2^2\right).
\]
Substituting back $x_0 \mapsto x_0+2x_1-x_2$ gives
\[
\Psi(F)=(x_0+2x_1-x_2)^2\left(\frac{1}{2}x_0x_1 + \frac{1}{4}x_2\right).
\]
Converting back to divided powers yields the decomposition in \Cref{Ex: recov. local GAD}.
\end{example}

\end{comment}

\begin{comment}
    

\begin{example}
    We transform the divided power polynomial \Cref{Ex: recov. local GAD} into a standard polynomial: 

    \[
\begin{aligned}
 &
\Psi(F)=\frac{1}{2}x_0^3 x_1 + 2x_0^2 x_1^2 + 2x_0 x_1^3 - x_0^2 x_1 x_2 - 2x_0 x_1^2 x_2 +\\
& \frac{1}{4} x_0^2 x_2^2 + \frac{3}{2}x_0 x_1 x_2^2 + x_1^2 x_2^2 - \frac{1}{2}x_0 x_2^3 - x_1 x_2^3 + \frac{1}{4}x_2^4.
\end{aligned}
\]

We make the substitution $x_0\mapsto x_0-2x_1+x_2$ and we have 
\[(\Tilde{\varphi} \circ \Psi )(F)=\frac{1}{2}x_0^3x_1-x_0^2x_1^2+\frac{1}{2}x_0^2x_1x_2+\frac{1}{4}x_0^2x_2^2.\]

We can factor $x_0^2$ in the usual way, 
\[(\Tilde{\varphi} \circ \Psi )(F)=x_0^2 \left(\frac{1}{2}x_0x_1-x_1^2+\frac{1}{2}x_1x_2+\frac{1}{4}x_2^2\right)\]
and then substitute back $x_0\mapsto x_0+2x_1-x_2$ so that we have

\[\Psi(F)=(x_0+2x_1-x_2)^2(\frac{1}{2}x_0x_1 + \frac{1}{4}x_2).\]
Converting back to divided powers we have the decomposition in \Cref{Ex: recov. local GAD}
\end{example}

\end{comment}

\begin{remark}
    In summary, when the socle degree of the minimal local apolar scheme $Z$ satisfies $\operatorname{socdeg}(Z) \le d$, the recovery of the GAD proceeds is extremely simple:
    \begin{itemize}
        \item Translate the support $\zeta$ to the origin by an affine change of coordinates, \item Factor out the maximal possible power of $x_0$ from $F$, 
        \item Revert the change of coordinates.
    \end{itemize}The resulting expression gives the desired GAD of $F$.
\end{remark}
\begin{comment}
    
\vspace{0.5cm}
\begin{mdframed}[]
\begin{alg}[Local GAD of $F$ if $\operatorname{socdeg}$ of the local apolar scheme $\leq \deg F$]\label{alg:GADEvinced}\end{alg}
%\begin{center}
%\underline{\textbf{Algorithm:} Cactus rank and decomposition}
%\end{center}
%\vspace{0.1cm}
\noindent\textbf{Input:} Homogeneous divided power polynomial $F$ of deg. $d$,  support $\zeta=(\zeta_1,\ldots, \zeta_n)$ 
of a minimal apolar scheme $Z=\text{Spec}(R/I)$ to $F$, and $\operatorname{socdeg}$ $k \leq \deg F$ of $R/I$.%, not bigger than to $\deg F$.
\\
    \textbf{Output:} GAD decomposition of $F$ in divided powers supported at $(1:\zeta)$
    \[F=(x_0+\zeta_1x_1+\ldots+\zeta_nx_n)^{(d-k)}\cdot H\]
    %Remark that the product is in divided powers.
 \begin{enumerate}
     \item\label{step: fromDividedPowers} Convert the divided powers polynomial $F$ into a standrard polynomial $\Psi(F)$ via $Y^{(\alpha)}\mapsto \frac{1}{\alpha!}x^{\alpha}$. 

    \item Make the substitution $x_0\mapsto x_0-\zeta_1x_1-\ldots - \zeta_n x_n$ in $\Psi(F)$. Let $(\Tilde{\varphi} \circ \Psi )(F)$ be the resulting polynomial.

    \item Divide $(\Tilde{\varphi} \circ \Psi )(F)$ by the highest possible power of $x_0$. Let $d-k$ such maximum exponent and $G$ be the resulting polynomial of the division.

    \item Make the change of coordinates $x_0\mapsto x_0+\zeta_1x_1+\ldots + \zeta_n x_n$ on $G$. Let $\Tilde{G}$ be the resulting polynomial.

    \item The scheme is evinced by the GAD of $F$
    \[F=(d-k)!(x_0+\zeta_1x_1+\ldots + \zeta_n x_n)^{(d-k)}\Psi^{-1}\Tilde{G}).\]
    \end{enumerate}
     \end{mdframed}
\vspace{0.5cm}
\end{comment}


\paragraph{\textbf{Socle degree $> d$:}} 
In this case, the minimal local apolar scheme $Z$ is evinced by a polynomial of degree higher than $F$, referred to as an \emph{extension of $F$}.  Recall from \Cref{def:extension} that in this case the minimal local scheme is evinced by a genuine extension $G$ of $F$, with $d'>d$ and $\partial_L^{d'-d} G = F$.

After an affine change of coordinates sending $\zeta$ to the origin, the functional becomes $
\varphi^*(\Lambda) = \mathds{1}_0 \circ g(\partial),
$ 
where $g \in K_{dp}[Y_1,\dots,Y_n]$ represents the dehomogenization of $G$ at $Y_0=1$.  


The key observation from \cite[Proposition~4]{BJPR} is that $g$ agrees with $F(Y_0=1)$ \emph{up to degree $d$}: the coefficients in degrees $\le d$ are completely determined by $F$.  
However, for degrees $> d$, then $g$ must contain additional terms that encode the full structure of $R/\ker H_\Lambda$.  
These higher-degree coefficients are determined uniquely by the values of $\Lambda$ on monomials of degree $>d$, which can be computed as
\[
\varphi^*(\Lambda(x^\beta)) = \Lambda(\operatorname{NF}_B((x+\zeta)^\beta)),
\]
where $\operatorname{NF}_B((x+\zeta)^\beta)$ denotes the normal form of $(x+\zeta)^\beta$ in the local algebra $R/\ker H_\Lambda$ with respect to a basis $B$.  

Once all coefficients of $g$ are determined (degrees $\le d$ from $F$, degrees $>d$ from $\Lambda$), the extension $G$ is obtained by homogenizing $g$ back to degree $d'=\operatorname{socdeg}(Z)$.  
Finally, reverting the coordinate change $\varphi$ expresses $G$ in the original coordinates, yielding a GAD that evinces $Z$.



%by computing the normal forms of suitable polynomials modulo the basis of $R/\ker H_\Lambda$ and evaluating them via $\Lambda$.



%Now we turn to the case where we have $F\in S_d^*$ a homogeneous polynomial and its minimal local apolar scheme $Z$ has socle degree greater than $d$. That is, if $Z$ is defined locally by Spec $(R/\ker H_\Lambda)$, then $R/\ker H_\Lambda$ has socle degree $k>d$. Then by \Cref{lemma:SocleAndExtesion} $Z$ is not evinced by a GAD of $F$, but of a polynomial of higher degree. Again the strategy is to move the support to the affine zero.


\begin{comment}
    
With the notation of the paragraph above, we consider the affine change of coordinates $\varphi:R\to R$ sending $\zeta$ to zero. Then $\varphi(R/\ker H_\Lambda) = R/\ker H_{\varphi^*(\Lambda)}$ is supported at zero and therefore $\varphi^*(\Lambda)$ is a polynomial $g=\varphi^*(\Lambda)\in K_{dp}[Y_1,\ldots, Y_n]$. The polynomials  $g$ and $F(Y_0=1)$ agree up to degree $d$. We show how to compute the higher degree coefficients. 


By definition, the coefficient in $g$ of $Y^{(\alpha)}$ for $|\alpha|>d$ is $$\varphi^*(\Lambda)(x^\alpha)=\Lambda((x-\zeta)^\alpha).$$

Note that if $|\beta|\leq d$, $\Lambda(x^{\beta})=F(Y_0=1)(x^\beta)$, and therefore it is already known from $F$. So we still need to show how to compute $\Lambda(x^{\beta})$ for $x^\beta$ with $|\beta|>d$. If $B$ is a basis of $R/\ker H_\Lambda$ and $NF_B(p)$ is the residue class of $p$ in $R/\ker H_\Lambda$ as a linear combination of the elements in $B$, then as shown in the proof of \cite[Theorem 4.2]{Bernard}:
\[\Lambda(x^\beta)=\Lambda(\operatorname{NF}(x^\beta).\]

Let $b$ be an element of a basis of $R/\ker H_\Lambda$. Since $\Lambda(b)$ is an entry of the matrix $H_\Lambda^B$, we can compute the image by the dual generator of the normal form of any polynomial.


\end{comment}


\begin{example}
    We can test this method with \cite[Example 4.6]{GAD}, with a previous change of coordinates.
{\small{
\[
\begin{aligned}
F= &144x_0^{(3)} + 284x_0^{(2)}x_1 + 444x_0x_1^{(2)} + 624x_1^{(3)} + 150x_0^{(2)}x_2 \\
&+ 150x_0x_1x_2 + 150x_1^{(2)}x_2 + 140x_0x_2^{(2)} + 140x_1x_2^{(2)} + 360x_2^{(3)} \\
&+ 428x_0^{(2)}x_3 + 708x_0x_1x_3 + 1028x_1^{(2)}x_3 + 660x_0x_2x_3 + 660x_1x_2x_3 \\
&+ 400x_2^{(2)}x_3 + 1256x_0x_3^{(2)} + 1816x_1x_3^{(2)} + 2040x_2x_3^{(2)} + 3552x_3^{(3)} \\
&+ 216x_0^{(2)}x_4 + 76x_0x_1x_4 + (-84)x_1^{(2)}x_4 + (-30)x_0x_2x_4 + (-30)x_1x_2x_4 \\
&+ (-140)x_2^{(2)}x_4 + 292x_0x_3x_4 + 12x_1x_3x_4 + (-420)x_2x_3x_4 + 184x_3^{(2)}x_4 \\
&+ (-576)x_0x_4^{(2)} + (-436)x_1x_4^{(2)} + (-90)x_2x_4^{(2)} + (-1012)x_3x_4^{(2)} + 936x_4^{(3)} \\
&+ 452x_0^{(2)}x_5 + 872x_0x_1x_5 + 1352x_1^{(2)}x_5 + 450x_0x_2x_5 + 450x_1x_2x_5 \\
&+ 420x_2^{(2)}x_5 + 1324x_0x_3x_5 + 2164x_1x_3x_5 + 1980x_2x_3x_5 + 3848x_3^{(2)}x_5 \\
&+ 628x_0x_4x_5 + 208x_1x_4x_5 + (-90)x_2x_4x_5 + 836x_3x_4x_5 + (-1708)x_4^{(2)}x_5 \\
&+ 1416x_0x_5^{(2)} + 2676x_1x_5^{(2)} + 1350x_2x_5^{(2)} + 4092x_3x_5^{(2)} + 1824x_4x_5^{(2)} + 4428x_5^{(3)}
\end{aligned}
\]
}}
    When we run the algorithm for the we find that for degree 6, which the number of essential variables of $F$ and thus the minimum possible cactus rank of $F$, the only possible complete staircase is $\{1,x_1,\ldots, x_5\}$, for which the matrices $\{M_{x_1}, \ldots , M_{x_5}\}$ commute and have unique eigenvalues. These unique eigenvalues are $1,0,2,-1,3$ respectively.  Hence, the minimal (local) apolar scheme is supported at the affine point $(1,0,2,-1,3)$. Via the methods described above, using the matrices $M_{x_i}$ one finds that the Hilbert of the local defining algebra $R/\ker H_\Lambda$ is $(1,2,1,1,1)$. Hence, the socle deree is $4>\deg F$. 

    \bigskip
    
    We consider every monomial $x^\alpha$ of degree 4, and compute the normal form of $(x+\zeta)^\alpha$ in the basis $\{1,x_1,\ldots,x_5\}$. It turns out that the only nonzero result is 
\[\operatorname{NF}_B((x_2-0)^4)=6\overline{x_5}-18.\]
    Looking at the dehomogenization of $F$, we see that $$\Lambda(x_2^4)=\Lambda(6x_5-18)=6\cdot 452-18\cdot 144=120.$$
    
    Thus, if $\varphi$ is the homogeneous change of coordinates sending $\zeta$ to 0, the scheme is defined locally by $\operatorname{Ann}(g)$, where 
    \[g=((\varphi^*(F))(Y_0=1)+ 120 x_2^{(4)}=
    120x_2^{(4)} + 360x_2^{(3)} + 120x_2^{(2)}x_3 + 20x_1^{(2)} + 140x_2^{(2)} + \]\[ 360x_2x_3 + 120x_3^{(2)} + 120x_2x_4 + 140x_1 + 150x_2 + 140x_3 + 360x_4 + 20x_5 + 144.
\]
    Then 
    $Z$ is evinced by $Y_0^{(0)}\cdot \operatorname{Homog}(g, Y_0, 4)$ (which agrees with \cite[Example 4.6]{GAD}), or in the original system of coordinates by 
    \[G=(Y_0+Y_1+2Y_3-Y_4 + 3Y_5)^{(0)}\cdot (\varphi^{-1})^*)(\operatorname{Homog}(g, Y_0, 4)).\]
\end{example}
In practice, the algorithm proceeds by
\begin{itemize}
    \item Translating the support $\zeta$ to the origin,
    \item Determining the higher-degree terms of the dual generator $g$ from the normal forms in $R/\ker H_\Lambda$,
    \item  Homogenizing $g$ to obtain the extension $G$, 
    \item Reverting the coordinate change to recover the GAD in the original variables.
\end{itemize}

\begin{comment}
    
\vspace{0.5cm}
\begin{mdframed}[]
\begin{alg}[Local GAD of $F$ when the socle degree of the local apolar scheme  exceed the degree of $F$]\label{alg:GADNotEvinced}\end{alg}
%\begin{center}
%\underline{\textbf{Algorithm:} Cactus rank and decomposition}
%\end{center}
\vspace{0.1cm}
\noindent\textbf{Input:} A homogeneous divided power polynomial $F$ of degree $d$, the support $(\zeta_1,\ldots, \zeta_n)$ of the minimal scheme $Z=\text{Spec}(R/\ker H_\Lambda)$ apolar to $F$, the socle degree $k$ of $R/I$, strictly bigger than $\deg F$, a basis $B$ of the algebra $R/\ker H_\Lambda$, the matrices corresponding to multiplication by $x_i$ in the basis $B$, and the values $\Lambda(b)$ for every $b\in B$. \\
    \textbf{Output:} A GAD decomposition of an extension$G$ of $F$ evincing the scheme $Z$.
 \begin{enumerate}

    \item Convert the divided powers polynomial $F$ into a standard polynomial $\Psi(F)$ via $Y^{(\alpha)}\mapsto \frac{1}{\alpha!}x^{\alpha}$. 

    \item Make the substitution $x_0\mapsto x_0-\zeta_1x_1-\ldots - \zeta_n x_n$ in $\Psi(F)$. Convert the result back to divided powers. Let $\varphi^*(F)$ be the resulting polynomial.

    \item Compute the normal form of $(x+\zeta)^\alpha$, for $d+1\leq |\alpha|\leq k$ (see \Cref{lemma: NormalForm}).

    \item Let 
    \[g'= \sum_{d+1\leq |\alpha|\leq k} \Lambda(\operatorname{NF}_B((x+\zeta)^\alpha))Y^{(\alpha)}.\]

    \item Define
    \[g:=(\varphi^*(F))(Y_0=1) + g'.\]

    \item Homogenize $g$, 
    \[G:=\operatorname{Homog}(g, Y_0).\]
    
    \item Compute $(\varphi^{-1})^*(G)$, that is, convert $G$ into non divided powers, make the substitution $x_0\mapsto x_0+\zeta_1x_1+\ldots + \zeta_n x_n$ and convert back to divided powers. 
    
    \item The scheme is evinced by the GAD
    \[(Y_0+\zeta_1Y_1+\ldots \zeta_nY_nY_n)^{(0)}\cdot (\varphi^{-1})^*(G).\]
    \end{enumerate}
     \end{mdframed}
\vspace{0.5cm}
\end{comment}

%\subsection{Hilbert function decomposition from Hankel operators}

\subsubsection{Global procedure for a minimal local apolar scheme and GAD}\label{Section:GlobalGAD}

The two cases considered above in the previous section whether the socle degree of an apolar scheme of a given $F$ is either $\le d$ or $> d$, can be unified into a single constructive procedure.  
Given a homogeneous polynomial $F \in S_d^*$, the goal is to determine a minimal \emph{local} apolar scheme $Z$ to $F$ and recover the corresponding GAD of $F$ (if $\operatorname{socdeg}(Z) \le d$) or of a genuine extension $G$ of $F$ (if $\operatorname{socdeg}(Z) > d$).

The standard cactus rank algorithm looks for a minimal apolar scheme to $F$, which, in general, may have support at multiple points.  
Here we are interested only in \emph{local} schemes, i.e., zero-dimensional schemes supported at a single point.  
As discussed in the local cactus rank algorithm, the search for a minimal apolar scheme can be constrained to the local setting by incorporating the locality conditions directly into the cactus rank procedure.  
This is achieved by imposing that all multiplication matrices $M_{x_i}$ not only commute but also have a \emph{single eigenvalue}, which coincides with the $i$-th coordinate of the support point.  

Once a minimal local scheme is found, its support $\zeta$ and socle degree $k$ are determined from the multiplication matrices and the Hilbert function.  
The recovery of the GAD then follows the two cases described earlier.

This leads to the following unified algorithm.

\vspace{0.5cm}
\begin{mdframed}[]
\begin{alg}[Minimal local apolar scheme and GAD from $F$]\label{alg:GlobalGAD}
\end{alg}
\vspace{0.1cm}
\noindent\textbf{Input:}  
Homogeneous divided power polynomial $F \in S_d^*$ of degree $d$.

\noindent\textbf{Output:}  
Minimal local apolar scheme $Z = \operatorname{Spec}(R/\ker H_\Lambda)$ to $F$ and a GAD decomposition of $F$ (if $\operatorname{socdeg}(Z) \le d$) or of an extension $G$ of $F$ (if $\operatorname{socdeg}(Z) > d$).

\begin{enumerate}
    \item \textbf{Construct the apolar algebra.}  
    Dehomogenize $F$ as $f = F(Y_0=1)$ and construct the Hankel matrix $H_\Lambda(h)$ with parameters $h_\alpha$ for $|\alpha|>d$.

    \item \textbf{Find the minimal local scheme.}

        \begin{enumerate}
            
        \item\label{alg:local:CHStar} Set $r=\text{rank } H_f^{\ostar}$ (see \Cref{notation: InitialRank}).

    \item\label{Step3} For each admissible Hilbert function $ h = (1,h_1,\dots,h_{d'}) $ of length $ r $, select $B\subseteq R$ a complete staircase of monomials with $|B|=r$ with $h_i$ monomials of degree $i$ and do: 
    \begin{enumerate}
     
        \item\label{alg:localC:h's} Find $h$'s such that:
        \begin{enumerate}
            
            \item $H_{\Lambda(h)}^B$ has nonzero determinant
            \item The multiplication operators $(M_{x_i})^t$ commute for all $i=1,\ldots,n$.
            \item The multiplication operators $M_{x_i}$ have a unique eigenvalue: 
            \[
            \text{char}_{M_{x_i}}(t) = (t - \tfrac{1}{r} \mathrm{tr}(M_{x_i}))^r \]
            
        \end{enumerate}
    
        \item If found, go to \Cref{Alg:Step:SupportAndSocle}. If not, go to \Cref{Step3} with another choice of $B$. If all choices of bases $B$ of length $r$ are performed go to \Cref{alg:localC:increase:r}.
   
    \end{enumerate}
    \item\label{alg:localC:increase:r} Set $r\to r+1$ and go to \Cref{Step3}.
        \end{enumerate}

    % Search for a basis $B$ corresponding to a complete staircase compatible with a local Hilbert function of length $r$, such that:
    % \begin{enumerate}
    %     \item $\det H_\Lambda^B(h) \neq 0$,
    %     \item all multiplication matrices $M_{x_i}$ commute,
    %     \item each $M_{x_i}$ has a single eigenvalue $\frac{1}{r}\mathrm{tr}(M_{x_i})$.
    % \end{enumerate}
    % Among all admissible $B$, choose the one with minimal length $r$.  
    This yields a local scheme $Z = \operatorname{Spec}(R/\ker H_\Lambda)$ supported at a single point.

    \item\label{Alg:Step:SupportAndSocle} \textbf{Determine the support and socle degree.}  
    \begin{enumerate}
        \item 
    Compute the support
    \[
    \zeta = \left(\frac{1}{r}\mathrm{tr}(M_{x_1}), \dots, \frac{1}{r}\mathrm{tr}(M_{x_n})\right).
    \]  
\item     Translate $\zeta$ to the origin via a homogeneous coordinate change $\varphi$.  
   \item 
    From the Hilbert function $(1, h_1, \ldots, h_k)$ of $R/\ker H_\Lambda$, determine the socle degree
    \[
    k = \operatorname{socdeg}(R/\ker H_\Lambda),
    \] i.e. the largest $j$ such that $h_j \neq 0$.

    \end{enumerate}
    \item \textbf{Recover the GAD (case distinction based on $k$).}
    \begin{itemize}
        \item \(\mathbf{k \le d}\) (no extension needed):  
        \begin{enumerate}
            \item Write $\varphi^*(F)$ in the form $Y_0^{(d-k)} \cdot H$ by factoring out the maximal possible power of $Y_0$.  
            \item Apply $\varphi^{-1}$ to recover the GAD of $F$ in the original coordinates.
        \end{enumerate}
        
        \item \(\mathbf{k > d}\) (extension needed):  
        \begin{enumerate}
            \item Dehomogenize $\varphi^*(F)$: $f = \varphi^*(F)(Y_0=1)$ (coefficients for degrees $\le d$ of $g$).  
            \item For each $k>|\alpha|>d$, compute $\Lambda(\operatorname{NF}_B((x+\zeta)^\alpha))$ to determine the higher-degree coefficients of $g$.  
            \item Homogenize $g$ to degree $k$ to obtain the extension $G$.  
            \item Apply $\varphi^{-1}$ to recover the GAD of $G$ in the original coordinates.
        \end{enumerate}
    \end{itemize}

    \item \textbf{Output.}  
    Return $Z = \operatorname{Spec}(R/\ker H_\Lambda)$ and the GAD of $F$ (if $k \le d$) or of $G$ (if $k > d$).
\end{enumerate}
\end{mdframed}

\vspace{0.5cm}


\begin{example}\label{Ex:GlobalAlg_k>d}
Consider the homogeneous divided power polynomial
$
F = x_0^{(2)} + x_0x_1 + x_1^{(2)} + x_0x_2 \in S_2^*$.
We apply \Cref{alg:GlobalGAD} to determine the minimal local apolar scheme $Z$ to $F$ and the corresponding GAD.

Setting $Y_0=1$, we obtain $
f = 1 + y_1 + y_1^{(2)} + y_2 \in K_{dp}[y_1,y_2]$. The coefficients of $f$ determine $\Lambda(y^\alpha)$ for all $|\alpha|\le d$.

We construct the Hankel matrix $H_\Lambda(h)$.
Imposing the local condition, we restrict to complete staircases compatible with Hilbert functions of local Gorenstein algebras, and require that all multiplication matrices $M_{x_i}$ commute and have a unique eigenvalue.  
A suitable basis is $
B = \{1, x_1, x_2, x_1^2, x_1x_2, x_2^2, x_1^3\}$,
with Hilbert function
$
H = (1,3,2,1)$ which already shows that the socle degree is $k=3 > d=2$. Thus $Z$ is not evinced by a GAD of $F$, but by a GAD of a degree-$k$ extension $G$ of $F$.

From the eigenvalues of the multiplication matrices, we find $\zeta = (0,0,0)$. No affine translation is needed.

Since $k>d$, we compute the higher degree terms of $g=\varphi^*(\Lambda)$ by evaluating $\Lambda$ on the normal forms of monomials of degree $3$.  
For $y_1^3$, the normal form is $x_1^3\in B$, and we find $\Lambda(x_1^3)=1$.  
All other monomials of degree $3$ reduce to $0$ in $B$.  
Thus
$
g = 1 + y_1 + y_2 + y_1^{(2)} + y_1^{(3)}$.

 
Homogenizing $g$ to degree $k=3$ gives
$
G = x_0^{(3)} + x_0^{(2)}x_1 + x_0^{(2)}x_2 + x_0x_1^{(2)} + x_1^{(3)}.$

Since $\zeta=0$, the linear form is $L=x_0$.  
Factoring $L^{(k-d)}=x_0^{(1)}$ we obtain
$G = x_0^{(1)} \cdot \left(x_0^{(2)} + x_0x_1 + x_0x_2 + x_1^{(2)} + x_1^{(3)}\right)$.

This is a valid local GAD evincing $Z$.
\end{example}



\subsubsection{Hilbert function via Hankel operators}
In this section, we recover the Hilbert function and its symmetric decomposition directly from a polynomial $F\in R^*$.
This step is essential for determining the socle degree of the apolar algebra $A=R/\ker H_\Lambda$, which in turn governs whether a GAD of $F$ exists or whether a genuine extension is needed. 
Our method, following \cite[Section~1]{BJPR}, exploits the filtrations $\mathfrak{m}^i$ (which corresponds to partial derivatives of $F$ of order at least $i$) and $(0:\mathfrak{m}^i)$ (which corresponds to partial derivatives of degree at most $i$).  
These can be accessed directly via Hankel operators, without explicitly computing a basis or multiplication matrices for $A$. 

%In this section we show how to obtain the Hilbert function of the algebra $\K[x_1,\ldots, x_n]/Ann(f)$ for any $f\in R$. It is based on an interpretation of the filtrations $(0:\m^{i})$ and $\m^{i}$ of the previous section in terms of order and degree of the partials, respectively, of $f$. 
%We follow closely \cite[Section 1]{BJPR}.


%We start with some definitions. 

\begin{definition}
    \begin{itemize}
        \item \leavevmode Let $\m=(x_1,\ldots, x_n)$ and $p\in R$. The \emph{order} of $p$ is the integer $k\geq 0$ such that $p\in \m^k$ but $p\not \in \m^{k+1}$.
        \item Let $F\in \K_{dp}[Y_1,\ldots, Y_n]$ and $p\in \operatorname{Im}(H_f)$. We call the order of $p$ the maximum order of a $q\in R$ such that $p=q\aprod f$.
    \end{itemize}
     
\end{definition}

Under the isomorphism $R/\ker H_F \cong \operatorname{Im}(H_F)$, the ideal $\m^{i}$ corresponds to partial derivatives of $F$ of order at least $i$, while $(0:\m^{i})$ corresponds to partial derivatives of degree at most $i$. The symmetric decomposition $\Delta_{Q_a}(i)$ counts the partials of $F$ of degree $i$ and order $d-a-i$, modulo partials of smaller degree and higher order.

From a computational viewpoint, the space of partials of order at least $k$ is spanned by $\{H_F(x^\alpha)\}_{|\alpha|\ge k}$, i.e., by the columns of the Hankel matrix indexed by monomials of degree at least $k$. The degree of each partial is read from the corresponding row indexing.

%The relationship with the Hankel operator comes from the simple observation that the space of partials of order at least $k$ is spanned by $\{H_f(x^\alpha)\}_{|\alpha|\geq k}$, i.e, by the columns of the Hankel matrix of $f$ that are indexed by monomials of degree at least $k$, and the degree of a partial is directly read from the rows of the Hankel. 

\begin{example}
    We compute the Hilbert function of the apolar algebra $R/\operatorname{Ann}(f)$ for $f=x_1^{(3)}x_2+x_3^{(3)} + x_2^{(2)}$. Let's start by computing its Hankel matrix. Since $f\in R^*$ is 0 on $R_{>4}$, only the matrix indexed by monomials of degree at most $3$ is needed\footnote{The first column corresponds to $1\aprod f=f$, so we can also eliminate it.}.

\[
\resizebox{\textwidth}{!}{$
\begin{array}{c|ccccccccccccccccccc}
                 & x_1 & x_2 & x_3 & x_1^2 & x_1x_2 & x_1x_3 & x_2^2 & x_2x_3 & x_3^2 & x_1^3 & x_1^2x_2 & x_1^2x_3 & x_1x_2^2 & x_1x_2x_3 & x_1x_3^2 & x_2^3 & x_2^2x_3 & x_3^3 \\
\hline 
0               & 0 & 0 & 0 & 0 & 0 & 0 & 1 & 0 & 0 & 0 & 0 & 0 & 0 & 0 & 0 & 0 & 0 & 1 \\
x_1             & 0 & 0 & 0 & 0 & 0 & 0 & 0 & 0 & 0 & 0 & 1 & 0 & 0 & 0 & 0 & 0 & 0 & 0 \\
x_2             & 0 & 1 & 0 & 0 & 0 & 0 & 0 & 0 & 0 & 1 & 0 & 0 & 0 & 0 & 0 & 0 & 0 & 0 \\
x_3             & 0 & 0 & 0 & 0 & 0 & 0 & 0 & 0 & 1 & 0 & 0 & 0 & 0 & 0 & 0 & 0 & 0 & 0 \\
x_1^2           & 0 & 0 & 0 & 0 & 1 & 0 & 0 & 0 & 0 & 0 & 0 & 0 & 0 & 0 & 0 & 0 & 0 & 0 \\
x_1x_2          & 0 & 0 & 0 & 1 & 0 & 0 & 0 & 0 & 0 & 0 & 0 & 0 & 0 & 0 & 0 & 0 & 0 & 0 \\
x_1x_3          & 0 & 0 & 0 & 0 & 0 & 0 & 0 & 0 & 0 & 0 & 0 & 0 & 0 & 0 & 0 & 0 & 0 & 0 \\
x_2^2           & 0 & 0 & 0 & 0 & 0 & 0 & 0 & 0 & 0 & 0 & 0 & 0 & 0 & 0 & 0 & 0 & 0 & 0 \\
x_2x_3          & 0 & 0 & 0 & 0 & 0 & 0 & 0 & 0 & 0 & 0 & 0 & 0 & 0 & 0 & 0 & 0 & 0 & 0 \\
x_3^2           & 0 & 0 & 1 & 0 & 0 & 0 & 0 & 0 & 0 & 0 & 0 & 0 & 0 & 0 & 0 & 0 & 0 & 0 \\
x_1^3           & 0 & 1 & 0 & 0 & 0 & 0 & 0 & 0 & 0 & 0 & 0 & 0 & 0 & 0 & 0 & 0 & 0 & 0 \\
x_1^2x_2        & 1 & 0 & 0 & 0 & 0 & 0 & 0 & 0 & 0 & 0 & 0 & 0 & 0 & 0 & 0 & 0 & 0 & 0 \\
x_1^2x_3        & 0 & 0 & 0 & 0 & 0 & 0 & 0 & 0 & 0 & 0 & 0 & 0 & 0 & 0 & 0 & 0 & 0 & 0 \\
x_1x_2^2        & 0 & 0 & 0 & 0 & 0 & 0 & 0 & 0 & 0 & 0 & 0 & 0 & 0 & 0 & 0 & 0 & 0 & 0 \\
x_1x_2x_3       & 0 & 0 & 0 & 0 & 0 & 0 & 0 & 0 & 0 & 0 & 0 & 0 & 0 & 0 & 0 & 0 & 0 & 0 \\
x_1x_3^2        & 0 & 0 & 0 & 0 & 0 & 0 & 0 & 0 & 0 & 0 & 0 & 0 & 0 & 0 & 0 & 0 & 0 & 0 \\
x_2^3           & 0 & 0 & 0 & 0 & 0 & 0 & 0 & 0 & 0 & 0 & 0 & 0 & 0 & 0 & 0 & 0 & 0 & 0 \\
x_2^2x_3        & 0 & 0 & 0 & 0 & 0 & 0 & 0 & 0 & 0 & 0 & 0 & 0 & 0 & 0 & 0 & 0 & 0 & 0 \\
x_2x_3^2        & 0 & 0 & 0 & 0 & 0 & 0 & 0 & 0 & 0 & 0 & 0 & 0 & 0 & 0 & 0 & 0 & 0 & 0 \\
x_3^3           & 0 & 0 & 0 & 0 & 0 & 0 & 0 & 0 & 0 & 0 & 0 & 0 & 0 & 0 & 0 & 0 & 0 & 0 \\
\end{array}
$}
\]


    The index on the left indicates the degree of the row. 
    
    From this matrix we can easily see the Hilbert function of $R/\operatorname{Ann}(f)$: the space $(0:\m^{3})/(0:\m^2)$ is the space of partials of degree exactly $3$, in this case only the first two columns correspond to partials of degree 3, and since they are linearly independent, $H_A(3)=2$. Now, in order to compute the partials of degree exactly 2, we have to remove the two columns corresponding to degree 3 partials, and we are left with the third, fourth and fifth column, which are linearly independent. Hence, $H_A(2)=3$. Finally, we know that $H_A(1)$ is the number of essential variables of $f$, in this case 3. We reach the same conclusion by eliminating the first 5 columns and looking at the degree 1 partials of the Hankel.  

    For the symmetric decomposition, we take the degree 3 partials and we look at the columns. The two degree 3 partials correspond to columns indexed by degree 1 monomials, so their order is one. Hence, $\Delta_{Q_0}(3)=2$, $\Delta_{Q_1}(3)=\Delta_{Q_2}(3)=0$.  
    The 3 degree 2 partials have order 1,2 and 2. Hence, $\Delta_{Q_0}(2)=2, \; \Delta_{Q_1}(2)=1$. Similarly, $\Delta_{Q_0}(1)=2, \Delta_{Q_1}(1)=1$. We obtain the symmetric decomposition of \Cref{example:table:HFdecomp}.
From this discussion we extract an explicit procedure to compute $H_A$ and its symmetric decomposition.

\begin{table}[h!] 
\centering
\caption{Hilbert function decomposition of $\K[x_1,x_2,x_3]/Ann(f)$} \label{example:table:HFdecomp}
\begin{tabular}{|l|l|l|l|l|l|}
\hline
Degree         & 0 & 1 & 2 & 3 & 4 \\ \hline
$\Delta_{Q_0}$ & 1 & 2 & 2 & 2 & 1 \\ \hline
$\Delta_{Q_1}$ & 0 & 1 & 1 & 0 & 0 \\ \hline
$\Delta_{Q_2}$ & 0 & 0 & 0 & 0 & 0 \\ \hline
Total          & 1 & 3 & 3 & 2 & 1 \\ \hline
\end{tabular}
\end{table}
%One can check that if we compute the symmetric decomposition by hand as in \cite[Example 8]{BJPR} we obtain the same result. 
\end{example}

From this example and the discussion above, we can deduce an algorithm to compute the Hilbert function and its decomposition from any $f\in \K_{dp}[Y_1,\ldots, Y_n]$:



\vspace{0.5cm}
\begin{mdframed}[]
\begin{alg}[Hilbert Function decomposition]\label{alg:SymmetricHF}\end{alg}
%\begin{center}
%\underline{\textbf{Algorithm:} Cactus rank and decomposition}
%\end{center}
\vspace{0.1cm}
\noindent\textbf{Input:} A polynomial $f\in \K_{dp}[Y_1,\ldots, Y_n]$. \\
    \textbf{Output:} The Hilbert function of the associated graded algebra of \\ $\K[x_1,\ldots, x_n]/\operatorname{Ann}(f)$ and its symmetric decomposition.

\begin{enumerate}
    \item Construct the matrix $H_{f}$ indexed by monomials of degree at most $\deg f-1$.
    \item Set $i=\deg f-1$.
    \item\label{alg: item: i} Let $v$ be the vector of columns of $H_F$ with at least one nonzero entry on the rows indexed by monomials of degree $i$.  
    
    \item  For $a$ from $\deg f-1$ to 1 do: 
    \begin{itemize}
        \item Take $w$ a subvector of $v$ of the columns in $H_f$ indexed by monomials of degree $d-i-a$.
        \item Set $\Delta_{Q_a}(i)=\text{rank}(w)$.
        \item Eliminate in $v$ the columns in $w$.
        \item $a\to a-1$.        
        \end{itemize}
        \item Eliminate in $H_f$ the columns in $v$.
        \item $r\to i-1$
        \item Return to step \Cref{alg: item: i}.
        
    \end{enumerate}
\end{mdframed}
\vspace{0.5cm}

\sloppy We can use a reverse engineering strategy to determine a polynomial $f\in \K_{dp}[Y_1,\ldots, Y_n]$ such that $\K[x_1,\ldots, x_n]/\operatorname{Ann}(f)$ has a given Hilbert function. 

\begin{example}
    We can test this with the previous example. We take the Hilbert function decomposition in \Cref{example:table:HFdecomp}, from which we know that the number of variables is 3 and the degree of $f$ is 4. Therefore we consider 
   {\small{ \[
\begin{aligned}
f =\;& a_0 
+ a_1 x_1 + a_2 x_2 + a_3 x_3 \\
&+ a_4 x_1^2 + a_5 x_1 x_2 + a_6 x_1 x_3 + a_7 x_2^2 + a_8 x_2 x_3 + a_9 x_3^2 \\
&+ a_{10} x_1^3 + a_{11} x_1^2 x_2 + a_{12} x_1^2 x_3 + a_{13} x_1 x_2^2 + a_{14} x_1 x_2 x_3 + a_{15} x_1 x_3^2 \\
&+ a_{16} x_2^3 + a_{17} x_2^2 x_3 + a_{18} x_2 x_3^2 + a_{19} x_3^3 \\
&+ a_{20} x_1^4 + a_{21} x_1^3 x_2 + a_{22} x_1^3 x_3 + a_{23} x_1^2 x_2^2 + a_{24} x_1^2 x_2 x_3 + a_{25} x_1^2 x_3^2 \\
&+ a_{26} x_1 x_2^3 + a_{27} x_1 x_2^2 x_3 + a_{28} x_1 x_2 x_3^2 + a_{29} x_1 x_3^3 \\
&+ a_{30} x_2^4 + a_{31} x_2^3 x_3 + a_{32} x_2^2 x_3^2 + a_{33} x_2 x_3^3 + a_{34} x_3^4
\end{aligned}
\]
    }}
 and we build its Hankel operator matrix:
\[\resizebox{\textwidth}{!}{$
\begin{array}{r|cccccccccccccccccccc}
& 1 & x_1 & x_2 & x_3 & x_1^2 & x_1x_2 & x_1x_3 & x_2^2 & x_2x_3 & x_3^2 & x_1^3 & x_1^2x_2 & x_1^2x_3 & x_1x_2^2 & x_1x_2x_3 & x_1x_3^2 & x_2^3 & x_2^2x_3 & x_3^3 \\
\hline 
 & 1 & 2 & 3 & 4 & 5 & 6 & 7 & 8 & 9 & 10 & 11 & 12 & 13 & 14 & 15 & 16 & 17 & 18 & 19 & 20 \\
\hline
1 & a_0   & a_1 & a_2 & a_3 & a_4 & a_5 & a_6 & a_7 & a_8 & a_9 & a_{10} & a_{11} & a_{12} & a_{13} & a_{14} & a_{15} & a_{16} & a_{17} & a_{18} & a_{19} \\
x_1 & a_1 & a_4 & a_5 & a_6 & a_{10} & a_{11} & a_{12} & a_{13} & a_{14} & a_{15} & a_{20} & a_{21} & a_{22} & a_{23} & a_{24} & a_{25} & a_{26} & a_{27} & a_{28} & a_{29} \\
x_2 & a_2 & a_5 & a_7 & a_8 & a_{11} & a_{13} & a_{14} & a_{16} & a_{17} & a_{18} & a_{21} & a_{23} & a_{24} & a_{26} & a_{27} & a_{28} & a_{30} & a_{31} & a_{32} & a_{33} \\
x_3 & a_3 & a_6 & a_8 & a_9 & a_{12} & a_{14} & a_{15} & a_{17} & a_{18} & a_{19} & a_{22} & a_{24} & a_{25} & a_{27} & a_{28} & a_{29} & a_{31} & a_{32} & a_{33} & a_{34} \\
x_1^2 & a_4 & a_{10} & a_{11} & a_{12} & a_{20} & a_{21} & a_{22} & a_{23} & a_{24} & a_{25} & 0 & 0 & 0 & 0 & 0 & 0 & 0 & 0 & 0 & 0 \\
x_1x_2 & a_5 & a_{11} & a_{13} & a_{14} & a_{21} & a_{23} & a_{24} & a_{26} & a_{27} & a_{28} & 0 & 0 & 0 & 0 & 0 & 0 & 0 & 0 & 0 & 0 \\
x_1x_3 & a_6 & a_{12} & a_{14} & a_{15} & a_{22} & a_{24} & a_{25} & a_{27} & a_{28} & a_{29} & 0 & 0 & 0 & 0 & 0 & 0 & 0 & 0 & 0 & 0 \\
x_2^2 & a_7 & a_{13} & a_{16} & a_{17} & a_{23} & a_{26} & a_{27} & a_{30} & a_{31} & a_{32} & 0 & 0 & 0 & 0 & 0 & 0 & 0 & 0 & 0 & 0 \\
x_2x_3 & a_8 & a_{14} & a_{17} & a_{18} & a_{24} & a_{27} & a_{28} & a_{31} & a_{32} & a_{33} & 0 & 0 & 0 & 0 & 0 & 0 & 0 & 0 & 0 & 0 \\
x_3^2 & a_9 & a_{15} & a_{18} & a_{19} & a_{25} & a_{28} & a_{29} & a_{32} & a_{33} & a_{34} & 0 & 0 & 0 & 0 & 0 & 0 & 0 & 0 & 0 & 0 \\
x_1^3 & a_{10} & a_{20} & a_{21} & a_{22} & 0 & 0 & 0 & 0 & 0 & 0 & 0 & 0 & 0 & 0 & 0 & 0 & 0 & 0 & 0 & 0 \\
x_1^2x_2 & a_{11} & a_{21} & a_{23} & a_{24} & 0 & 0 & 0 & 0 & 0 & 0 & 0 & 0 & 0 & 0 & 0 & 0 & 0 & 0 & 0 & 0 \\
x_1^2x_3 & a_{12} & a_{22} & a_{24} & a_{25} & 0 & 0 & 0 & 0 & 0 & 0 & 0 & 0 & 0 & 0 & 0 & 0 & 0 & 0 & 0 & 0 \\
x_1x_2^2 & a_{13} & a_{23} & a_{26} & a_{27} & 0 & 0 & 0 & 0 & 0 & 0 & 0 & 0 & 0 & 0 & 0 & 0 & 0 & 0 & 0 & 0 \\
x_1x_2x_3 & a_{14} & a_{24} & a_{27} & a_{28} & 0 & 0 & 0 & 0 & 0 & 0 & 0 & 0 & 0 & 0 & 0 & 0 & 0 & 0 & 0 & 0 \\
x_1x_3^2 & a_{15} & a_{25} & a_{28} & a_{29} & 0 & 0 & 0 & 0 & 0 & 0 & 0 & 0 & 0 & 0 & 0 & 0 & 0 & 0 & 0 & 0 \\
x_2^3 & a_{16} & a_{26} & a_{30} & a_{31} & 0 & 0 & 0 & 0 & 0 & 0 & 0 & 0 & 0 & 0 & 0 & 0 & 0 & 0 & 0 & 0 \\
x_2^2x_3 & a_{17} & a_{27} & a_{31} & a_{32} & 0 & 0 & 0 & 0 & 0 & 0 & 0 & 0 & 0 & 0 & 0 & 0 & 0 & 0 & 0 & 0 \\
x_2x_3^2 & a_{18} & a_{28} & a_{32} & a_{33} & 0 & 0 & 0 & 0 & 0 & 0 & 0 & 0 & 0 & 0 & 0 & 0 & 0 & 0 & 0 & 0 \\
x_3^3 & a_{19} & a_{29} & a_{33} & a_{34} & 0 & 0 & 0 & 0 & 0 & 0 & 0 & 0 & 0 & 0 & 0 & 0 & 0 & 0 & 0 & 0 \\
\end{array}
$}\]


Imposing that $\Delta_{Q_0}(3)=2$ we have that the matrix 
{\small{
\[
%\resizebox{1.5}{!}{$
\begin{array}{r|cccc}
& 1 & x_1 & x_2 & x_3\\
\hline 
1 & 0   & a_1 & a_2 & a_3\\
x_1 & a_1 & a_4 & a_5 & a_6 \\
x_2 & a_2 & a_5 & a_7 & a_8 \\
x_3 & a_3 & a_6 & a_8 & a_9 \\
x_1^2 & a_4 & a_{10} & a_{11} & a_{12} \\
x_1x_2 & a_5 & a_{11} & a_{13} & a_{14}\\
x_1x_3 & a_6 & a_{12} & a_{14} & a_{15}\\
x_2^2 & a_7 & a_{13} & a_{16} & a_{17} \\
x_2x_3 & a_8 & a_{14} & a_{17} & a_{18}\\
x_3^2 & a_9 & a_{15} & a_{18} & a_{19}\\
x_1^3 & a_{10} & a_{20} & a_{21} & a_{22}\\
x_1^2x_2 & a_{11} & a_{21} & a_{23} & a_{24}\\
x_1^2x_3 & a_{12} & a_{22} & a_{24} & a_{25}\\
x_1x_2^2 & a_{13} & a_{23} & a_{26} & a_{27}\\
x_1x_2x_3 & a_{14} & a_{24} & a_{27} & a_{28}\\
x_1x_3^2 & a_{15} & a_{25} & a_{28} & a_{29} \\
x_2^3 & a_{16} & a_{26} & a_{30} & a_{31} \\
x_2^2x_3 & a_{17} & a_{27} & a_{31} & a_{32}\\
x_2x_3^2 & a_{18} & a_{28} & a_{32} & a_{33} \\
x_3^3 & a_{19} & a_{29} & a_{33} & a_{34} \\
\end{array}\]
}}
has two columns corresponding to degree 3 polynomials that are linearly independent. A possibility is $a_{21}=1$, $a_7=1$, and the rest of the coefficients 0 of the first two columns and the degree 3 part of the last column equal to zero. Note that so far $f$ is determined by coefficients satisfying this choice. We could continue with this process until we recover the coefficients of $f$ (with possibly other solutions having the same Hilbert function).

\end{example}

\begin{comment}
    
\begin{remark}
    If $f$ is a homogeneous polynomial of degree $d$, then if $H_{(\alpha+\beta)}$ is the entry indexed by the row $x^{\alpha}$ and column $x^\beta$, it is clear that $H_{\alpha + \beta}=0$ if $|\alpha+\beta|\neq d$. Moreover, we know from \cite{Iarrobino} that the Hilbert function is symmetric and the symmetric decomposition is just $\Delta_{Q_0}(i)=H_A(i)$, and $\Delta_{Q_a}(i)=0$ for all $a>0$. By the previous arguments we have that $\Delta_{Q_a}(i)$ is the rank of the matrix indexed by monomials 
    %$x^{\alpha}$ with $|\alpha|=d-i $
    of total degree $d-i$ on the columns, and of total degree $i$ on the rows. This matrix is the \emph{catalecticant matrix} of order $i$, $Cat_F(i,d-i)$ (see \cite{IK}). Thus, this method shows that $H_A(i)=\text{rank} Cat_F(i,d-i)$, which is formally proven in \cite{Geramita}.   
\end{remark}

One would like to study deeper this method and explore whether it is possible to obtain a parametrization of local Gorenstein algebras similar to the one for homogeneous polynomials using Catalecticant matrices (see \cite{Geramita}). 

\end{comment}



\bibliographystyle{alphaurl}
\bibliography{bibliography}
\end{document}
